% Le communiqué radio — Allemand 1ère A — Cours signé OnBuch+
% Programme officiel MINESEC. Compiler avec : tectonic ALA17-le-communique-radio.tex
\newcommand{\DOCMATIERE}{Allemand}
\newcommand{\DOCNIVEAU}{1ère A}
\newcommand{\DOCMODULE}{Chapitre 4 : Média et communication}
\newcommand{\DOCLECON}{ALA17}
\newcommand{\DOCTITRE}{Le communiqué radio}
% OnBuch+ — préambule « cours signés » ENRICHI (compilé avec tectonic / XeLaTeX)
% Système visuel « Pop » d'OnBuch+ : fond crème (jamais blanc), bordures encre
% épaisses, ombres « dures » décalées, titres Archivo Black en capitales.
% Version MATHÉMATIQUES : graphes (pgfplots/tikz), boîtes pédagogiques
% (définition, propriété, méthode, attention, le savais-tu, exercice résolu,
% exercice, TP, bilan), page de garde et signature OnBuch+ ancrée en bas.
%
% Macros attendues dans main.tex AVANT \input{preamble} :
%   \DOCMATIERE  \DOCNIVEAU  \DOCMODULE  \DOCLECON (numéro)  \DOCTITRE
\documentclass[11pt]{article}

\usepackage{fontspec}
\usepackage[ngerman,french]{babel} % français = langue principale ; l'allemand via \all{..} / textebox
\usepackage[a4paper,margin=2cm,top=3.4cm,bottom=2.8cm,headheight=1.4cm,footskip=1.3cm]{geometry}
\usepackage{amsmath,amssymb,amsfonts}
\usepackage{mathtools} % \xmapsto, \coloneqq... souvent utilisés par les modèles
\usepackage{cancel} % simplifications barrées (bilans de forces)
% \abs{z} (module, valeur absolue) et \norm{u} (norme), utilisés par les modèles
\providecommand{\abs}{}\let\abs\relax\DeclarePairedDelimiter{\abs}{\lvert}{\rvert}
\providecommand{\norm}{}\let\norm\relax\DeclarePairedDelimiter{\norm}{\lVert}{\rVert}
\usepackage[table]{xcolor} % [table] : fournit aussi \rowcolors (alternance de lignes)
\usepackage{pagecolor}
\usepackage{tikz}
\usetikzlibrary{arrows.meta,positioning,calc,shapes.geometric,shapes.misc,shapes.arrows,decorations.pathmorphing,decorations.pathreplacing,decorations.markings,patterns,shadows,mindmap,trees,fit,backgrounds,matrix,intersections,angles,quotes}
\usepackage{pgfplots}
\usepackage[version=4]{mhchem} % équations nucléaires \ce{^{235}_{92}U}
\usepackage[european,siunitx]{circuitikz} % schémas électriques
\pgfplotsset{compat=1.18}
\usepgfplotslibrary{fillbetween,groupplots} % aires entre courbes (calcul intégral)
\usepackage{enumitem}
\usepackage{fancyhdr}
% [most] charge toutes les bibliothèques tcolorbox (listings, minted, raster,
% documentation...) et gonfle inutilement la mémoire TeX sur un document long
% (~300 boîtes) : on ne charge que ce qui est réellement utilisé.
\usepackage[skins,breakable]{tcolorbox}
\usepackage{graphicx}
\usepackage{colortbl}
\usepackage{booktabs}
\usepackage{tabularx}
\usepackage{diagbox} % case d'en-tête diagonale des tableaux à double entrée
% Colonnes extensibles centrée / gauche / droite (C, L, R), souvent utilisées par les modèles
\newcolumntype{C}{>{\centering\arraybackslash}X}
\newcolumntype{L}{>{\raggedright\arraybackslash}X}
\newcolumntype{R}{>{\raggedleft\arraybackslash}X}
\usepackage{array}
\usepackage{multicol}
\usepackage{multirow}
\usepackage{siunitx}
% parse-numbers=false : accepte aussi \SI{2,5 \times 10^{-3}}{...}
\sisetup{locale=FR,output-decimal-marker={,},group-minimum-digits=4,parse-numbers=false}
\DeclareSIUnit\ohm{\ensuremath{\symup{\Omega}}} % U+2126 absent de la police
\DeclareSIUnit\division{div} % réglages d'oscilloscope (ms/div, V/div)
\DeclareSIUnit\tr{tr} % tours (vitesse de rotation, ex. \SI{1500}{\tr\per\minute})
\DeclareSIUnit\molar{M} % concentration molaire (ex. \SI{3}{\molar}, \SI{1}{\micro\molar})
\DeclareSIUnit\an{an} % année (le modèle confond parfois avec \year, un compteur TeX) — ex. \SI{5}{\kilo\meter\per\an}
\DeclareSIUnit\FCFA{FCFA} % franc CFA (ex. \SI{500}{\FCFA\per\kilo\gram})
\DeclareSIUnit\degreeBrix{\textdegree Brix} % teneur en sucre d'un sirop/jus (réfractométrie)
\DeclareSIUnit\month{mois} % le modèle utilise parfois \month, un compteur TeX, comme unité de durée
\usepackage{qrcode}
% \degree : commande fréquemment utilisée par les modèles pour un angle ou une
% température en dehors de \SI{}{} (non fournie par les paquets chargés).
\providecommand{\degree}{^{\circ}}
% \qed (fin de démonstration), utilisé par les modèles sans amsthm
\providecommand{\qed}{\hfill$\square$}
% \barre : double barre des valeurs interdites dans un tableau de variations (tkz-tab)
\providecommand{\barre}{\ensuremath{\|}}
% environnement proof (amsthm non chargé) utilisé par les modèles
\ifdefined\proof\else\newenvironment{proof}[1][Démonstration]{\par\noindent\textit{#1.}\ }{\qed\par}\fi
\usepackage{lastpage}
\usepackage{hyperref}
\hypersetup{hidelinks}

% ============================================================
% Palette « Pop » OnBuch+ — mêmes hex que la classe P (plus_ui.dart)
% ============================================================
\definecolor{poporange}{HTML}{F59321}
\definecolor{popdark}{HTML}{E07A0C}
\definecolor{popcream}{HTML}{FFF6E8}
\definecolor{popink}{HTML}{1C1714}
\definecolor{poppurple}{HTML}{7A5AE0}
\definecolor{popgreen}{HTML}{2E9E6A}
\definecolor{poppink}{HTML}{E0457B}
\definecolor{popblue}{HTML}{2F7DE1}
\definecolor{popgold}{HTML}{C98A12}
\definecolor{popmuted}{HTML}{8A7F76}
\definecolor{popline}{HTML}{EADFD2}
\definecolor{popgray}{HTML}{8A7F76}
\colorlet{popgrey}{popgray}
% Alias tolérants (le modèle nomme parfois les couleurs différemment)
\colorlet{popwhite}{white}
\colorlet{popblack}{popink}
\colorlet{popred}{poppink}
\colorlet{popyellow}{popgold}
% Teintes claires (fonds de tableaux/encadrés) — le modèle nomme parfois
% les couleurs avec un suffixe L (ex. popblueL) sans les avoir définies.
\colorlet{poporangeL}{poporange!20}
\colorlet{popdarkL}{popdark!20}
\colorlet{poppurpleL}{poppurple!20}
\colorlet{popgreenL}{popgreen!20}
\colorlet{poppinkL}{poppink!20}
\colorlet{popblueL}{popblue!20}
\colorlet{popgoldL}{popgold!20}
\colorlet{popredL}{popred!20}
\colorlet{popgrayL}{popgray!20}
% Teintes claires pour fonds de tableaux / zones
\colorlet{popblueL}{popblue!10!white}
\colorlet{poporangeL}{poporange!14!white}
\colorlet{popgreenL}{popgreen!12!white}
\colorlet{poppurpleL}{poppurple!12!white}
\colorlet{poppinkL}{poppink!12!white}
\colorlet{popgoldL}{popgold!14!white}

\pagecolor{popcream}

% ============================================================
% Typographie
% ============================================================
\defaultfontfeatures{Ligatures=TeX}
\setmainfont{PlusJakartaSans-Medium.ttf}[Path=../../../fonts/,BoldFont=PlusJakartaSans-Bold.ttf]
% Symboles, API et emojis absents de la police principale : repli sur DejaVu Sans (caractères actifs)
\newfontface\symfont{DejaVuSans.ttf}[Path=../../../fonts/]
\makeatletter
\newcommand\symactive[1]{\catcode#1=13 \begingroup\lccode`\~=#1 \lowercase{\endgroup\def~}{{\symfont\char#1}}}
\newcommand\symblank[1]{\catcode#1=13 \begingroup\lccode`\~=#1 \lowercase{\endgroup\def~}{}}
\newcommand\symhyph[1]{\catcode#1=13 \begingroup\lccode`\~=#1 \lowercase{\endgroup\def~}{\mbox{-}}}
\newcount\symc
\newcommand\symrange[2]{\symc=#1 \@whilenum\symc<#2 \do{\expandafter\symactive\expandafter{\the\symc}\advance\symc by 1 }}
\newcommand\symblankrange[2]{\symc=#1 \@whilenum\symc<#2 \do{\expandafter\symblank\expandafter{\the\symc}\advance\symc by 1 }}
\makeatother
\symrange{"0250}{"02FF}\symactive{"2713}\symactive{"2714}\symactive{"2717}\symactive{"2718}\symactive{"2610}\symactive{"2611}\symactive{"2612}
\symhyph{"2011}\symhyph{"2E1A}\symblankrange{"1F300}{"1FAFF}\symblank{"FE0F}
% Maths en sans-serif (Fira Math) avec chiffres et lettres droites de Plus
% Jakarta, pour que les formules se fondent dans le texte.
\usepackage[math-style=ISO,bold-style=ISO]{unicode-math}
\setmathfont{FiraMath-Regular.otf}[Path=../../../fonts/]
\setmathfont{PlusJakartaSans-Medium.ttf}[Path=../../../fonts/,range={up/num,up/{Latin,latin},\mathup}]
\usepackage{icomma} % virgule décimale sans espace en mode math
% Caractères mathématiques Unicode absents des polices de texte (Plus Jakarta,
% Archivo Black) : rendus par leur équivalent mathématique, en texte comme en
% formule — sinon ils s'affichent en carré vide (ex. « Arithmétique dans ℤ »).
\usepackage{newunicodechar}
\newunicodechar{ℝ}{\ensuremath{\mathbb{R}}}
\newunicodechar{ℤ}{\ensuremath{\mathbb{Z}}}
\newunicodechar{ℂ}{\ensuremath{\mathbb{C}}}
\newunicodechar{ℕ}{\ensuremath{\mathbb{N}}}
\newunicodechar{ℚ}{\ensuremath{\mathbb{Q}}}
\newunicodechar{𝔻}{\ensuremath{\mathbb{D}}}
\newunicodechar{α}{\ensuremath{\alpha}}
\newunicodechar{β}{\ensuremath{\beta}}
\newunicodechar{γ}{\ensuremath{\gamma}}
\newunicodechar{δ}{\ensuremath{\delta}}
\newunicodechar{ε}{\ensuremath{\varepsilon}}
\newunicodechar{θ}{\ensuremath{\theta}}
\newunicodechar{λ}{\ensuremath{\lambda}}
\newunicodechar{μ}{\ensuremath{\mu}}
\newunicodechar{σ}{\ensuremath{\sigma}}
\newunicodechar{τ}{\ensuremath{\tau}}
\newunicodechar{φ}{\ensuremath{\varphi}}
\newunicodechar{ψ}{\ensuremath{\psi}}
\newunicodechar{ω}{\ensuremath{\omega}}
\newunicodechar{Σ}{\ensuremath{\Sigma}}
% majuscules produites par \MakeUppercase dans les titres (α-aminés -> Α)
\newunicodechar{Α}{\ensuremath{\alpha}}
\newunicodechar{Β}{\ensuremath{\beta}}
\newunicodechar{Γ}{\ensuremath{\Gamma}}
\newunicodechar{Δ}{\ensuremath{\Delta}}
\newunicodechar{Ε}{\ensuremath{\varepsilon}}
\newunicodechar{Θ}{\ensuremath{\Theta}}
\newunicodechar{Λ}{\ensuremath{\Lambda}}
\newunicodechar{Μ}{\ensuremath{\mu}}
\newunicodechar{Τ}{\ensuremath{\tau}}
\newunicodechar{Φ}{\ensuremath{\Phi}}
\newunicodechar{Ψ}{\ensuremath{\Psi}}
\newunicodechar{Ω}{\ensuremath{\Omega}}
\newunicodechar{↦}{\ensuremath{\mapsto}}
\newunicodechar{∈}{\ensuremath{\in}}
\newunicodechar{∉}{\ensuremath{\notin}}
\newunicodechar{∧}{\ensuremath{\wedge}}
\newunicodechar{⇔}{\ensuremath{\Leftrightarrow}}
\newunicodechar{⟺}{\ensuremath{\Longleftrightarrow}}
\newunicodechar{⇒}{\ensuremath{\Rightarrow}}
\newunicodechar{⟹}{\ensuremath{\Longrightarrow}}
\newunicodechar{∘}{\ensuremath{\circ}}
\newunicodechar{∪}{\ensuremath{\cup}}
\newunicodechar{∩}{\ensuremath{\cap}}
\newunicodechar{≡}{\ensuremath{\equiv}}
\newunicodechar{⊂}{\ensuremath{\subset}}
\newunicodechar{⊥}{\ensuremath{\perp}}
\newunicodechar{∃}{\ensuremath{\exists}}
\newunicodechar{∀}{\ensuremath{\forall}}
\newunicodechar{∛}{\ensuremath{\sqrt[3]{\,}}}
\newunicodechar{∜}{\ensuremath{\sqrt[4]{\,}}}
\newunicodechar{⃗}{\ensuremath{{}^{\to}}}
\newunicodechar{ⁿ}{\ifmmode^{n}\else\textsuperscript{\ensuremath{n}}\fi}
\newunicodechar{ˣ}{\ifmmode^{x}\else\textsuperscript{\ensuremath{x}}\fi}
\newunicodechar{ᵇ}{\ifmmode^{b}\else\textsuperscript{\ensuremath{b}}\fi}
\newunicodechar{ʳ}{\ifmmode^{r}\else\textsuperscript{\ensuremath{r}}\fi}
\newunicodechar{ᵘ}{\ifmmode^{u}\else\textsuperscript{\ensuremath{u}}\fi}
\newunicodechar{ʸ}{\ifmmode^{y}\else\textsuperscript{\ensuremath{y}}\fi}
\newunicodechar{ᵃ}{\ifmmode^{a}\else\textsuperscript{\ensuremath{a}}\fi}
\newunicodechar{ᵉ}{\ifmmode^{e}\else\textsuperscript{\ensuremath{e}}\fi}
\newunicodechar{ᵏ}{\ifmmode^{k}\else\textsuperscript{\ensuremath{k}}\fi}
\newunicodechar{ᵖ}{\ifmmode^{p}\else\textsuperscript{\ensuremath{p}}\fi}
\newunicodechar{ᴺ}{\ifmmode^{N}\else\textsuperscript{\ensuremath{N}}\fi}
\newunicodechar{ᶜ}{\ifmmode^{c}\else\textsuperscript{\ensuremath{c}}\fi}
\newunicodechar{ʰ}{\ifmmode^{h}\else\textsuperscript{\ensuremath{h}}\fi}
\newunicodechar{ᵠ}{\ifmmode^{q}\else\textsuperscript{\ensuremath{q}}\fi}
\newunicodechar{ₙ}{\ifmmode_{n}\else\textsubscript{\ensuremath{n}}\fi}
\newunicodechar{ₐ}{\ifmmode_{a}\else\textsubscript{\ensuremath{a}}\fi}
\newunicodechar{ᵢ}{\ifmmode_{i}\else\textsubscript{\ensuremath{i}}\fi}
\newunicodechar{ᵣ}{\ifmmode_{r}\else\textsubscript{\ensuremath{r}}\fi}
\newunicodechar{ₖ}{\ifmmode_{k}\else\textsubscript{\ensuremath{k}}\fi}
\newunicodechar{ᵦ}{\ifmmode_{\beta}\else\textsubscript{\ensuremath{\beta}}\fi}
\newunicodechar{ₓ}{\ifmmode_{x}\else\textsubscript{\ensuremath{x}}\fi}
\newfontfamily\popdisplay{ArchivoBlack-Regular.ttf}[Path=../../../fonts/]
\newfontfamily\poplogo{SpaceGrotesk-Bold.ttf}[Path=../../../fonts/]
\newfontfamily\popsemi{PlusJakartaSans-SemiBold.ttf}[Path=../../../fonts/]
\newfontfamily\popxbold{PlusJakartaSans-ExtraBold.ttf}[Path=../../../fonts/]
\newfontfamily\popxboldls{PlusJakartaSans-ExtraBold.ttf}[Path=../../../fonts/,LetterSpace=3.0]

\setlist[enumerate]{leftmargin=1.7em,itemsep=5pt,topsep=4pt}
\setlist[itemize]{leftmargin=1.7em,itemsep=4pt,topsep=4pt,label={\color{poporange}\textbullet}}
\setlength{\parindent}{0pt}
\setlength{\parskip}{5pt}
\renewcommand{\arraystretch}{1.25}

% pgfplots : style Pop par défaut
\pgfplotsset{
  every axis/.append style={axis line style={popink,line width=1pt},tick style={popink},
    grid style={popline,line width=0.5pt},label style={font=\small},tick label style={font=\footnotesize}},
  popcurve/.style={poporange,line width=1.8pt,no markers,smooth},
}
% Même style disponible dans un \draw TikZ simple (hors pgfplots)
\tikzset{popcurve/.style={poporange,line width=1.8pt}}
\tikzset{popgrid/.style={popline,very thin}}
\colorlet{popcurve}{poporange} % le nom du style est parfois utilisé comme couleur

% ============================================================
% Pill / squiggle
% ============================================================
\newcommand{\poppill}[3]{%
  \tikz[baseline=-0.55ex]\node[draw=popink,line width=1.2pt,rounded corners=2.6pt,%
    fill=#2,inner xsep=6.5pt,inner ysep=3pt]{%
    \popxboldls\fontsize{7}{7}\selectfont\color{#3}\MakeUppercase{#1}};%
}
\newcommand{\popsquiggle}[1]{%
  \tikz[baseline=-0.4ex]{
    \draw[#1,line width=2.6pt,line cap=round]
      plot [smooth] coordinates {(0,0.09) (0.34,0.01) (0.68,0.21) (1.02,0.01) (1.36,0.10)};
  }%
}
% Mot-clé surligné (marqueur orange sous le texte)
\newcommand{\cle}[1]{{\popxbold\color{popdark}#1}}

% ============================================================
% En-tête / pied
% ============================================================
\pagestyle{fancy}
\fancyhf{}
\renewcommand{\headrulewidth}{0pt}
\renewcommand{\footrulewidth}{0pt}
\lhead{\raisebox{-0.15ex}{\poplogo\fontsize{13}{13}\selectfont\color{popink}OnBuch\color{poporange}+}}
\rhead{\poppill{\DOCMATIERE\ \textbullet\ \DOCNIVEAU\ \textbullet\ Leçon \DOCLECON}{white}{popink}}
\lfoot{\popxbold\fontsize{7.6}{7.6}\selectfont\color{popmuted}\MakeUppercase{Cours signé OnBuch+}\ \textbullet\ {\popsemi\DOCTITRE}}
\rfoot{\tikz[baseline=-0.6ex]\node[rounded corners=8.5pt,fill=popink,minimum height=17pt,minimum width=17pt,inner xsep=5pt,inner sep=0]{\popxbold\fontsize{8}{8}\selectfont\color{popcream}\thepage\,/\,\pageref*{LastPage}};}

% ============================================================
% Titres
% ============================================================
\newcommand{\chaptitre}[2]{%
  \begin{tcolorbox}[enhanced,colback=poporange,colframe=popink,boxrule=2.6pt,arc=11pt,
      drop shadow=popink,left=15pt,right=15pt,top=12pt,bottom=11pt,before skip=3pt,after skip=18pt]
    \poppill{#2}{popink}{popcream}\par\vspace{9pt}
    {\raggedright\hyphenpenalty=10000\exhyphenpenalty=10000\popdisplay\fontsize{22}{24}\selectfont\color{popcream}\MakeUppercase{#1}\par}
  \end{tcolorbox}
}
% Section numérotée : pastille orange avec le numéro + titre Archivo
\newcounter{popsec}
\newcommand{\coursec}[1]{%
  \stepcounter{popsec}\setcounter{popsub}{0}%
  \par\vspace{16pt}\noindent
  \begin{minipage}{\linewidth}
  \tikz[baseline=-0.7ex]\node[draw=popink,line width=1.6pt,fill=poporange,rounded corners=4pt,minimum size=22pt,inner sep=0]{\popdisplay\fontsize{12}{12}\selectfont\color{popcream}\thepopsec};%
  \hspace{8pt}{\popdisplay\fontsize{15}{17}\selectfont\color{popink}\MakeUppercase{#1}}\par
  \vspace{4pt}\noindent{\color{poporange}\rule{36pt}{3.6pt}}
  \end{minipage}\par\nopagebreak\vspace{8pt}
}
\newcounter{popsub}
\newcommand{\courssub}[1]{%
  \stepcounter{popsub}%
  \par\vspace{9pt}\noindent{\popxbold\fontsize{11.5}{13}\selectfont\color{popink}{\color{poporange}\thepopsec.\thepopsub}\hspace{6pt}#1}\par\nopagebreak\vspace{4pt}
}

% ============================================================
% Boîtes pédagogiques — même recette : fond blanc, bordure épaisse colorée,
% ombre dure encre, badge plein. Titre optionnel [#1].
% ============================================================
\tcbset{popbase/.style={breakable,enhanced,colback=white,boxrule=2.3pt,arc=9pt,
  shadow={3pt}{-3pt}{0pt}{popink},left=14pt,right=14pt,top=11pt,bottom=11pt,before skip=11pt,after skip=15pt,
  fontupper=\popsemi\fontsize{10.6}{14.5}\selectfont\color{popink},
  fontlower=\popsemi\fontsize{10.6}{14.5}\selectfont\color{popink}}}
% Coche dessinée (le glyphe ✓ n'existe pas dans les polices du document)
\newcommand{\popcheck}[1]{\tikz[baseline=-0.5ex]\draw[#1,line width=1.6pt,line cap=round,line join=round](-0.13,0.0)--(-0.04,-0.09)--(0.13,0.1);}
\providecommand{\coche}{\popcheck{popgreen}}
\providecommand{\resultat}[1]{\par\smallskip\noindent\fcolorbox{popgreen}{popgreenL}{\parbox{\dimexpr\linewidth-2\fboxsep-2\fboxrule\relax}{\textbf{Résultat.} #1}}\par\smallskip}
\newcommand{\popboxhead}[3]{\poppill{#1}{#2}{white}\ifx\relax#3\relax\else\hspace{6pt}{\popxbold\fontsize{10.6}{12}\selectfont\color{popink}#3}\fi\par\vspace{7pt}}

\newtcolorbox{aretenir}[1][]{popbase,colframe=popgreen,before upper={\popboxhead{À retenir}{popgreen}{#1}}}
% Texte en allemand (document de lecture, dialogue, extrait) : cadre
% d'étude. Un titre optionnel (source, auteur) s'affiche dans l'en-tête.
\newtcolorbox{textebox}[1][]{popbase,colframe=popblue,before upper={\popboxhead{Text}{popblue}{#1}\begin{otherlanguage}{ngerman}},after upper={\end{otherlanguage}}}
% Marques ✓ / ✗ (forme correcte / fautive) souvent utilisées par les modèles
\providecommand{\cross}{\textcolor{poppink}{\ensuremath{\times}}}
\providecommand{\xmark}{\cross}\providecommand{\crossmark}{\cross}\providecommand{\cmark}{\ensuremath{\checkmark}}
% Ligne de réponse / justification à compléter (inventée par les modèles)
\providecommand{\justification}[1]{\par\noindent\textit{Justification~:} \dotfill\par}
% \textipa (package tipa non chargé) : la police ne couvre pas l'API -> simple passage
\providecommand{\textipa}[1]{#1}
% Soulignement (ulem non chargé) : \uline, \ul
\providecommand{\uline}[1]{\underline{#1}}\providecommand{\ul}[1]{\underline{#1}}
% \glossaire{\item ...} inventée par les modèles : liste de glossaire
\providecommand{\glossaire}[1]{\begin{itemize}#1\end{itemize}}
% \alert (beamer) inventée par les modèles : texte mis en évidence
\providecommand{\alert}[1]{\textbf{\textcolor{poppink}{#1}}}
% variantes ASCII d'ordinaux écrites en macro
\providecommand{\iere}{\textsuperscript{re}}\providecommand{\ieme}{\textsuperscript{e}}\providecommand{\ere}{\textsuperscript{re}}\providecommand{\eme}{\textsuperscript{e}}\providecommand{\er}{\textsuperscript{er}}\providecommand{\ie}{\textsuperscript{e}}
% Ordinaux français écrits en macro par les modèles : 1\ière, 3\ième
\newcommand{\ière}{\textsuperscript{re}}\newcommand{\ième}{\textsuperscript{e}}
% \exemple (inventée par les modèles) : étiquette « Exemple : »
\providecommand{\exemple}{\textbf{Exemple~:}\ }
% \square utilisable aussi hors mode math (cases à cocher)
\AtBeginDocument{\let\squareMath\square\renewcommand{\square}{\ensuremath{\squareMath}}}
% Mot ou phrase en allemand à l'intérieur d'un paragraphe français
\newcommand{\all}[1]{\ifmmode\text{\textit{#1}}\else\begingroup\selectlanguage{ngerman}\itshape #1\endgroup\fi}
\let\esp\all % alias toléré
\newtcolorbox{exemplebox}[1][]{popbase,colframe=poporange,before upper={\popboxhead{Exemple}{poporange}{#1}}}
\newtcolorbox{definition}[1][]{popbase,colframe=popblue,before upper={\popboxhead{Définition}{popblue}{#1}}}
\newtcolorbox{propriete}[1][]{popbase,colframe=poppurple,before upper={\popboxhead{Propriété}{poppurple}{#1}}}
\newtcolorbox{methode}[1][]{popbase,colframe=popgold,before upper={\popboxhead{Méthode}{popgold}{#1}}}
\newtcolorbox{attention}[1][]{popbase,colframe=poppink,before upper={\popboxhead{Attention}{poppink}{#1}}}
\newtcolorbox{experience}[1][]{popbase,colframe=popblue,colback=popblueL,before upper={\popboxhead{Expérience}{popblue}{#1}}}
% « Le savais-tu ? » : carte violette pleine (contexte, culture, Cameroun)
\newtcolorbox{savaistu}[1][]{popbase,colback=poppurple,colframe=popink,
  fontupper=\popsemi\fontsize{10.4}{14.2}\selectfont\color{white},
  before upper={\poppill{Le savais-tu\,?}{white}{poppurple}\ifx\relax#1\relax\else\hspace{6pt}{\popxbold\color{white}#1}\fi\par\vspace{7pt}}}
% Exercice résolu : énoncé en haut, \tcblower puis la solution sur fond crème
\newcounter{popexo}\newcounter{popexr}
\newtcolorbox{exoresolu}[1][]{popbase,colframe=poporange,colbacklower=popcream,
  segmentation style={popink,line width=1.2pt,dashed},
  before upper={\stepcounter{popexr}\popboxhead{Exercice résolu \thepopexr}{poporange}{#1}},
  before lower={\poppill{Solution}{popink}{popcream}\par\vspace{6pt}}}
% Exercice (non résolu en place) — la correction est dans la partie Corrigés
\newtcolorbox{exercice}[1][]{popbase,colframe=popink,
  before upper={\stepcounter{popexo}\popboxhead{Exercice \thepopexo}{popink}{#1}}}
% Corrigé
\newtcolorbox{corrige}[1][]{popbase,colframe=popgreen,colback=popgreenL,
  before upper={\popboxhead{Corrigé}{popgreen}{#1}}}
% Objectifs / prérequis / bilan
\newtcolorbox{objectifs}{popbase,colframe=popink,colback=poporangeL,
  before upper={\popboxhead{Objectifs de la leçon}{popink}{}}}
\newtcolorbox{prerequis}{popbase,colframe=popink,colback=popblueL,
  before upper={\popboxhead{Prérequis}{popblue}{}}}
\newtcolorbox{bilan}{popbase,colframe=popink,colback=white,boxrule=2.8pt,
  before upper={\popboxhead{Fiche bilan}{poporange}{L'essentiel à connaître pour l'examen}}}
% Figure encadrée avec légende
\newtcolorbox{popfigure}[1][]{enhanced,breakable=false,colback=white,colframe=popline,boxrule=1.4pt,arc=8pt,
  left=10pt,right=10pt,top=10pt,bottom=8pt,before skip=10pt,after skip=12pt,halign=center,
  lower separated=false,halign lower=center,
  fontlower=\popsemi\fontsize{9}{11.5}\selectfont\color{popmuted},
  #1}
\newcommand{\legende}[1]{\par\vspace{6pt}{\popsemi\fontsize{9}{11.5}\selectfont\color{popmuted}#1}}

% ============================================================
% Signature OnBuch+
% ============================================================
\newcommand{\poplogoinline}[1]{{\poplogo\fontsize{#1}{#1}\selectfont\color{popink}OnBuch\color{poporange}+}}
\newcommand{\signatureonbuch}{%
  \begin{tcolorbox}[enhanced,colback=popink,colframe=popink,boxrule=2.2pt,arc=10pt,
      drop shadow=poporange,left=14pt,right=14pt,top=10pt,bottom=10pt,before skip=0pt,after skip=0pt]
    \tikz[baseline=-0.7ex]\node[circle,fill=popgreen,draw=popcream,line width=1.3pt,minimum size=22pt,inner sep=0]{\popcheck{white}};%
    \hspace{9pt}\begin{minipage}[c]{0.62\linewidth}
      {\popxbold\fontsize{12}{13}\selectfont\color{popcream}Signé\ \textbullet\ {\poplogo OnBuch\color{poporange}+}}\par\vspace{2pt}
      {\raggedright\popsemi\fontsize{8}{10.5}\selectfont\color{popline}\MakeUppercase{Équipe pédagogique OnBuch+}\\\MakeUppercase{Conforme au programme officiel MINESEC}\par}
    \end{minipage}\hfill
    {\poplogo\fontsize{22}{22}\selectfont\color{popcream}OnBuch\color{poporange}+}
  \end{tcolorbox}%
}

% ============================================================
% Page de garde
% #1 titre · #2 matière · #3 niveau · #4 module · #5 n° leçon · #6 durée ·
% #7 description / objectifs (texte ou itemize)
% ============================================================
\newcommand{\pagegarde}[7]{%
  \thispagestyle{empty}
  \newgeometry{a4paper,margin=1.7cm,top=1.6cm,bottom=1.4cm}%
  \begin{tikzpicture}[remember picture,overlay]
    \fill[poporange!18] ([xshift=-3.2cm,yshift=-3.6cm]current page.north east) circle (4.6cm);
    \fill[poppurple!14] ([xshift=2.4cm,yshift=7.6cm]current page.south west) circle (3.8cm);
  \end{tikzpicture}%
  {\poplogo\fontsize{30}{30}\selectfont\color{popink}OnBuch\color{poporange}+}\hfill
  \raisebox{0.6ex}{\poppill{Cours signé \textbullet\ Premium}{popink}{popcream}}\par
  \vspace{1pt}\popsquiggle{poppurple}\par
  \vspace{8pt}
  \begin{tcolorbox}[enhanced,colback=poporange,colframe=popink,boxrule=2.8pt,arc=13pt,
      drop shadow=popink,left=18pt,right=18pt,top=12pt,bottom=13pt,after skip=10pt]
    \poppill{#2 \textbullet\ #3}{popink}{popcream}\hspace{5pt}\poppill{#4}{white}{popink}\par\vspace{8pt}
    {\popdisplay\fontsize{14}{14}\selectfont\color{popink}LEÇON #5}\par\vspace{4pt}
    {\raggedright\hyphenpenalty=10000\exhyphenpenalty=10000\popdisplay\fontsize{25}{27}\selectfont\color{popcream}\MakeUppercase{#1}\par}
    \vspace{8pt}
    {\popxbold\fontsize{9.5}{9.5}\selectfont\color{popink}\MakeUppercase{Durée conseillée : #6}}
  \end{tcolorbox}
  \begin{tcolorbox}[enhanced,colback=white,colframe=popink,boxrule=2.2pt,arc=10pt,
      drop shadow=popink,left=16pt,right=16pt,top=10pt,bottom=10pt,after skip=8pt]
    \poppill{Dans cette leçon}{poppurple}{white}\par\vspace{5pt}
    \popsemi\fontsize{9.3}{12.6}\selectfont\color{popink}#7
  \end{tcolorbox}
  \noindent\begin{minipage}[t]{0.487\linewidth}
    \begin{tcolorbox}[enhanced,colback=popgreen,colframe=popink,boxrule=2.2pt,arc=10pt,
        drop shadow=popink,left=11pt,right=11pt,top=10pt,bottom=10pt,height=3.1cm,valign=center]
      \tcbox[colback=white,colframe=popink,boxrule=1.3pt,arc=5pt,boxsep=0pt,nobeforeafter,
        left=3pt,right=3pt,top=3pt,bottom=3pt,valign=center]{\qrcode[height=1.9cm]{https://play.google.com/store/apps/details?id=com.luvvix.onbuch&hl=fr}}%
      \hspace{8pt}\begin{minipage}[c]{0.5\linewidth}
        {\popdisplay\fontsize{11}{12.5}\selectfont\color{white}\MakeUppercase{Télécharge OnBuch}}\par\vspace{4pt}
        {\raggedright\popsemi\fontsize{8.4}{11}\selectfont\color{white}Gratuit \textbullet\ Google Play\\Tuteur IA, annales, concours, résultats\par}
      \end{minipage}
    \end{tcolorbox}
  \end{minipage}\hfill
  \begin{minipage}[t]{0.487\linewidth}
    \begin{tcolorbox}[enhanced,colback=poppurple,colframe=popink,boxrule=2.2pt,arc=10pt,
        drop shadow=popink,left=11pt,right=11pt,top=10pt,bottom=10pt,height=3.1cm,valign=center]
      \tikz[baseline=-0.7ex]\node[circle,fill=white,draw=popink,line width=1.3pt,minimum size=1.6cm,inner sep=0]{\popdisplay\fontsize{24}{24}\selectfont\color{poppurple}+};%
      \hspace{8pt}\begin{minipage}[c]{0.6\linewidth}
        {\popdisplay\fontsize{11}{12.5}\selectfont\color{white}\MakeUppercase{Passe à OnBuch+}}\par\vspace{4pt}
        {\raggedright\popsemi\fontsize{8.4}{11}\selectfont\color{white}Tous les cours signés, Tuteur IA illimité, fiches PDF — onglet OnBuch+ de l'app\par}
      \end{minipage}
    \end{tcolorbox}
  \end{minipage}
  \vspace{8pt}
  \popcontenu
  \vfill
  \signatureonbuch
  \restoregeometry
  \newpage
}

% Rangée « Ce cours contient » (page de garde)
\newcommand{\poptile}[3]{% #1 couleur #2 gros texte #3 légende
  \begin{tcolorbox}[enhanced,colback=#1,colframe=popink,boxrule=2pt,arc=9pt,shadow={2.5pt}{-2.5pt}{0pt}{popink},
      left=6pt,right=6pt,top=9pt,bottom=9pt,halign=center,valign=center,height=1.75cm,width=\linewidth,nobeforeafter]
    {\popdisplay\fontsize{12.5}{13}\selectfont\color{white}\MakeUppercase{#2}}\par\vspace{3pt}
    {\centering\popsemi\fontsize{7.8}{9.5}\selectfont\color{white}#3\par}
  \end{tcolorbox}}
\newcommand{\popcontenu}{%
  \par\noindent{\popxboldls\fontsize{7.5}{7.5}\selectfont\color{popmuted}CE COURS SIGNÉ CONTIENT}\par\vspace{7pt}
  \noindent\begin{minipage}[t]{0.235\linewidth}\poptile{popblue}{Cours}{complet, illustré, pas à pas}\end{minipage}\hfill
  \begin{minipage}[t]{0.235\linewidth}\poptile{poporange}{Méthodes}{et exercices résolus}\end{minipage}\hfill
  \begin{minipage}[t]{0.235\linewidth}\poptile{poppink}{Exercices}{type examen + corrigés}\end{minipage}\hfill
  \begin{minipage}[t]{0.235\linewidth}\poptile{popgreen}{Fiche bilan}{l'essentiel à retenir}\end{minipage}\par}

% Sommaire
\newtcolorbox{sommaire}{enhanced,colback=white,colframe=popink,boxrule=2.2pt,arc=10pt,
  drop shadow=popink,left=18pt,right=18pt,top=13pt,bottom=13pt,before skip=6pt,after skip=20pt,
  before upper={\poppill{Sommaire}{poppurple}{white}\par\vspace{9pt}\popsemi\fontsize{10.6}{14.5}\selectfont\color{popink}}}

% Signature de fin de cours — poussée tout en bas de la dernière page
\newcommand{\findecours}{%
  \par\vfill\nopagebreak
  \begin{center}\popsquiggle{poporange}\end{center}\vspace{4pt}
  \signatureonbuch
}
\AtBeginDocument{\def\llbracket{\mathopen{[\mkern-3.5mu[}}\def\rrbracket{\mathclose{]\mkern-3.5mu]}}} % intervalles d'entiers (glyphe ⟦ absent de la police)
% Glyphes absents de Fira Math : redéfinis à partir de glyphes disponibles
\AtBeginDocument{%
  \def\setminus{\mathbin{\backslash}}\let\smallsetminus\setminus
  \def\vdots{\mathord{\vbox{\baselineskip3.2pt\lineskiplimit0pt\kern4pt\hbox{.}\hbox{.}\hbox{.}}}}%
  \def\ddots{\mathinner{\mkern1mu\raise6.5pt\hbox{.}\mkern2mu\raise3.6pt\hbox{.}\mkern2mu\raise0.7pt\hbox{.}\mkern1mu}}%
  \def\triangle{\mathord{\scalebox{0.9}{$\symup{\Delta}$}}}%
  \def\nabla{\mathord{\raisebox{\depth}{\scalebox{1}[-1]{$\symup{\Delta}$}}}}%
  \def\checkmark{\popcheck{popdark}}%
  \def\star{\mathbin{\ast}}\def\bigstar{\mathord{\ast}}%
  \def\diamond{\mathbin{\scalebox{0.6}{$\Diamond$}}}%
  \def\bigcup{\mathop{\mathchoice{\raisebox{-0.3em}{\scalebox{1.7}{$\cup$}}}{\scalebox{1.25}{$\cup$}}{\cup}{\cup}}\displaylimits}%
  \def\bigcap{\mathop{\mathchoice{\raisebox{-0.3em}{\scalebox{1.7}{$\cap$}}}{\scalebox{1.25}{$\cap$}}{\cap}{\cap}}\displaylimits}%
  \def\lesseqqgtr{\mathrel{\lessgtr}}\def\bigoplus{\mathop{\oplus}\displaylimits}%
}
\providecommand{\ssi}{si et seulement si} % abréviation fréquente
\AtBeginDocument{\providecommand{\wideparen}{\overparen}} % arc de cercle

%% FIGFIX-BEGIN v5 — correctifs globaux des figures (docs/onbuch-generation/tools/figfix)
\usepackage{adjustbox}\usepackage{etoolbox}\tcbuselibrary{raster}
% toute figure TikZ/pgfplots plus large que la ligne est réduite à la largeur disponible
\BeforeBeginEnvironment{tikzpicture}{\begin{adjustbox}{max width=\linewidth}}
\AfterEndEnvironment{tikzpicture}{\end{adjustbox}}
% légendes pgfplots lisibles par-dessus les courbes
\pgfplotsset{every axis/.append style={legend style={fill=white,fill opacity=0.92,text opacity=1,draw=black!15,inner sep=3pt}}}
% étiquettes d'arêtes (syntaxe edge["..."]) avec fond blanc
\tikzset{every edge quotes/.append style={fill=white,inner sep=1.2pt}}
%% FIGFIX-END
\begin{document}

\pagegarde{Le communiqué radio}{Allemand}{1ère A}{Chapitre 4 : Média et communication}{ALA17}{2 h}{Cette leçon de type « texte » propose l'étude d'un communiqué radio allemand authentique : compréhension globale et détaillée, lexique des médias et de la communication, grammaire (accusatif, datif, prépositions de temps) et production guidée d'un communiqué. Elle prépare l'élève de 1ère A à comprendre et à rédiger des messages d'information courts et précis.
\vspace{4pt}
\begin{multicols}{2}\small
\begin{itemize}
\item À la fin de cette leçon, je sais identifier la structure et la fonction d'un communiqué radio allemand.
\item Je sais comprendre globalement puis en détail un texte d'information radiophonique.
\item Je sais employer correctement le vocabulaire thématique : das Radio, die Mitteilung, die Durchsage, der Sprecher, die Meldung, die Veranstaltung.
\item Je sais former les familles de mots, les articles et les pluriels du lexique des médias.
\item Je sais distinguer et utiliser l'accusatif et le datif après verbes et prépositions.
\item Je sais utiliser les prépositions de temps (am, um, von...bis, seit, vor, nach, in) dans des phrases.
\end{itemize}\end{multicols}}

\begin{sommaire}
\begin{multicols}{2}
\begin{enumerate}
\item Découverte : le texte support
\item Compréhension du texte
\item Lexique : les médias et la radio
\item Grammaire 1 : accusatif et datif
\item Grammaire 2 : les prépositions de temps
\item Méthode du commentaire et production guidée
\item Activité d'intégration
\item Méthodes et exercices résolus
\item Exercices
\item Corrigés des exercices
\item Fiche bilan
\end{enumerate}
\end{multicols}
\end{sommaire}

\begin{objectifs}
\begin{itemize}
\item À la fin de cette leçon, je sais identifier la structure et la fonction d'un communiqué radio allemand.
\item Je sais comprendre globalement puis en détail un texte d'information radiophonique.
\item Je sais employer correctement le vocabulaire thématique : das Radio, die Mitteilung, die Durchsage, der Sprecher, die Meldung, die Veranstaltung.
\item Je sais former les familles de mots, les articles et les pluriels du lexique des médias.
\item Je sais distinguer et utiliser l'accusatif et le datif après verbes et prépositions.
\item Je sais utiliser les prépositions de temps (am, um, von...bis, seit, vor, nach, in) dans des phrases.
\item Je sais répondre par écrit et à l'oral à des questions de compréhension sur un texte.
\item Je sais rédiger un résumé et un court communiqué radio en allemand.
\end{itemize}
\end{objectifs}

\begin{prerequis}
\begin{itemize}
\item Les articles définis et indéfinis au nominatif (der, die, das / ein, eine) vus en Seconde.
\item La conjugaison des verbes réguliers et des verbes forts courants au présent et au perfect.
\item La place du verbe en deuxième position dans la phrase déclarative.
\item Le vocabulaire de base de la vie quotidienne et des horaires (Uhrzeit, Wochentage).
\item Les prépositions simples vues en Seconde (in, an, auf, mit, für).
\end{itemize}
\end{prerequis}

\newpage

\coursec{Découverte : le texte support}

Il est 7 heures à Yaoundé. Dans une maison du quartier Mvog-Ada, la radio annonce le report d'un match du championnat et la fermeture de plusieurs routes après la pluie de la nuit. Au même moment, à Berlin, un passager pressé entend à la gare : \all{Achtung, eine wichtige Durchsage...} Que ce soit au Cameroun ou en Allemagne, le communiqué radio informe, annonce et oriente des milliers d'auditeurs chaque jour. Mais comment ce message est-il construit ? Quels mots et quelles structures le speaker utilise-t-il ? C'est la problématique de notre leçon : comprendre un communiqué radio allemand, en repérer les informations essentielles et apprendre à en produire un nous-mêmes.

\courssub{Qu'est-ce qu'un communiqué radio ?}

\begin{definition}[Le communiqué radio]
Le \cle{communiqué radio} (en allemand \all{die Durchsage} ou \all{die Mitteilung}) est un \textbf{message oral bref et informatif}, lu par un speaker (\all{der Sprecher} / \all{die Sprecherin}) et adressé à un \textbf{large public}. Il annonce un événement (fête, manifestation, match), une perturbation (route fermée, train retardé) ou une information urgente. Il répond toujours à cinq questions : \emph{qui ? quoi ? quand ? où ? pourquoi ?}
\end{definition}

\begin{definition}[Die Durchsage und die Mitteilung]
\all{die Durchsage, -n} : message officiel lu à la radio, à la gare, à l'aéroport ou dans un lieu public (une annonce sonore). \all{die Mitteilung, -en} : communication d'une information, souvent plus générale (une note, un avis). Exemple : \all{Bitte beachten Sie die Durchsage!} « Veuillez écouter l'annonce ! »
\end{definition}

\begin{savaistu}[La radio en Allemagne]
En Allemagne, chaque Land possède ses propres stations de radio publiques : le BR en Bavière, le WDR en Rhénanie-du-Nord-Westphalie, le SWR dans le Sud-Ouest... Elles sont regroupées dans l'ARD, la première chaîne publique. C'est un peu comme la CRTV au Cameroun, qui informe les auditeurs de Yaoundé à Garoua. Les Durchsagen (annonces) y rythment la journée : météo, circulation, événements culturels.
\end{savaistu}

\courssub{Le texte support : Eine wichtige Durchsage}

\begin{methode}[Avant de lire : la stratégie de la première lecture]
\begin{enumerate}
\item Lis le texte \textbf{une première fois en silence}, sans dictionnaire : ton seul but est de saisir le \textbf{thème général} (de quoi parle-t-on ?).
\item Souligne les \textbf{mots transparents} (mots proches du français) : \all{Radio}, \all{Information}, \all{Kultur}, \all{Programm}, \all{Busse}... Ils te donnent déjà la moitié du sens.
\item Entoure les \textbf{chiffres} : dates, heures, durées. Dans un communiqué, ce sont les informations les plus importantes.
\item Formule une \textbf{hypothèse de sens} en une phrase française, puis vérifie-la à la deuxième lecture.
\end{enumerate}
\end{methode}

\textbf{Hypothèses de sens.} Avant même de lire, observons le titre : \all{Eine wichtige Durchsage}. \all{wichtig} ressemble à... rien en français, mais le contexte (une radio) nous aide : un message « important ». Les mots transparents \all{Kulturfest}, \all{informieren}, \all{Frequenz} nous laissent deviner : la radio annonce une fête culturelle et donne des informations pratiques. Vérifions maintenant cette hypothèse en lisant le texte.

\begin{textebox}[Eine wichtige Durchsage — Radio Stadtwelle, 7 Uhr morgens]
Liebe Zuhörer! Hier ist Radio Stadtwelle mit einer wichtigen Durchsage. Am Samstag, dem 14. Oktober, findet in der Innenstadt ein großes Kulturfest statt. Das Fest beginnt um 10 Uhr und dauert bis 18 Uhr. Auf dem Marktplatz spielen drei Musikgruppen, und es gibt Essen aus vielen Ländern. Der Eintritt ist frei!

Wegen des Festes ist die Hauptstraße von 8 Uhr bis 20 Uhr für den Verkehr gesperrt. Die Busse der Linien 4 und 7 fahren über die Ringstraße. Bitte planen Sie mehr Zeit für Ihre Fahrt ein.

Und noch eine Meldung: Am Sonntag findet im Stadtpark ein Fußballspiel für Kinder statt. Anmeldungen nimmt das Jugendamt bis Freitag entgegen.

Wir informieren Sie jede Stunde über die Lage im Verkehr. Bleiben Sie auf unserer Frequenz! Und jetzt kommt die Musik...
\end{textebox}

\textbf{Glossaire en marge} (mots nouveaux à noter dans le cahier de vocabulaire) :

\begin{center}
\begin{tabularx}{\linewidth}{|X|X|}\hline
\rowcolor{popblueL} \textbf{Deutsch} & \textbf{Français} \\ \hline
die Durchsage, -n & l'annonce, le communiqué \\ \hline
der Zuhörer, - / die Zuhörerin, -nen & l'auditeur / l'auditrice \\ \hline
stattfinden (findet statt, fand statt, hat stattgefunden) & avoir lieu \\ \hline
die Innenstadt, -städte & le centre-ville \\ \hline
der Eintritt, -e & l'entrée (le prix d'entrée) \\ \hline
der Verkehr (Sg.) & la circulation \\ \hline
gesperrt (sperren) & fermé(e), bloqué(e) \\ \hline
die Umleitung, -en & la déviation \\ \hline
die Meldung, -en & la nouvelle, l'information \\ \hline
die Frequenz, -en & la fréquence (radio) \\ \hline
\end{tabularx}
\end{center}

\begin{attention}[Das Radio, pas die Radio !]
\all{das Radio} est un nom \textbf{neutre} : on dit \all{das Radio}, jamais \all{die Radio}, même si « la radio » est féminin en français. Pluriel : \all{die Radios}. Autre piège : \all{der Verkehr} signifie « la circulation », pas « le trafic de drogue » ; et attention au verbe à particule \all{stattfinden} : \all{Das Fest findet um 10 Uhr statt.} — la particule \all{statt} va en fin de phrase.
\end{attention}

\courssub{Repérage des éléments d'annonce : les cinq questions}

Un communiqué radio bien construit répond à cinq questions. Cherchons les réponses dans notre texte :

\begin{center}
\begin{tabularx}{\linewidth}{|l|X|X|}\hline
\rowcolor{popblueL} \textbf{Question} & \textbf{Réponse dans le texte} & \textbf{Français} \\ \hline
Wer ? (qui ?) & Radio Stadtwelle, der Sprecher & la radio, le speaker \\ \hline
Was ? (quoi ?) & ein großes Kulturfest & une grande fête culturelle \\ \hline
Wann ? (quand ?) & am Samstag, dem 14. Oktober, 10--18 Uhr & samedi 14 octobre, de 10 h à 18 h \\ \hline
Wo ? (où ?) & in der Innenstadt, auf dem Marktplatz & au centre-ville, sur la place du marché \\ \hline
Warum ? (pourquoi ?) & wegen des Festes: Straße gesperrt & à cause de la fête : rue fermée \\ \hline
\end{tabularx}
\end{center}

\begin{popfigure}
\begin{center}
\begin{tikzpicture}[mindmap, grow cyclic, every node/.style={concept}, concept color=popblue!40, root concept/.append style={concept color=poporange!50, font=\bfseries}, level 1/.append style={level distance=3.2cm, sibling angle=72}, text=popdark]
\node[root concept]{Der Radio\-kommentar}
  child[concept color=popgreen!40]{ node[align=center]{Wer?\\ \footnotesize der Sprecher} }
  child[concept color=popblue!30]{ node[align=center]{Was?\\ \footnotesize das Kulturfest} }
  child[concept color=poppurple!35]{ node[align=center]{Wann?\\ \footnotesize am Samstag} }
  child[concept color=poppink!40]{ node[align=center]{Wo?\\ \footnotesize die Innenstadt} }
  child[concept color=popgold!45]{ node[align=center]{Warum?\\ \footnotesize Information} };
\end{tikzpicture}
\end{center}
\legende{Carte mentale : les cinq questions d'un communiqué radio}
\end{popfigure}

\courssub{De la lecture silencieuse à la lecture à voix haute}

Après la lecture silencieuse, l'enseignant lit le texte \textbf{à voix haute} : c'est le modèle de prononciation. Écoute et observe :

\begin{itemize}
\item \all{Zuhörer} se prononce « tsou-rö-reur » : le \all{Z} allemand sonne « ts ».
\item \all{wichtig} se prononce « vikh-tikh » : le \all{w} allemand sonne comme un « v » français.
\item \all{Frequenz} : « fré-kvents » ; \all{stattfinden} : « chtat-fin-den » (\all{st-} en début de mot se prononce « cht »).
\item Le speaker articule, marque une pause après chaque information et \textbf{met le ton sur les chiffres} : \all{um ZEHN Uhr}, \all{bis ACHTZEHN Uhr}.
\end{itemize}

\begin{popfigure}
\begin{center}
\begin{tikzpicture}[node distance=0.9cm, every node/.style={font=\small, text=popdark}]
\node[draw=popblue, fill=popblueL, rounded corners, minimum width=2.6cm, minimum height=1cm, align=center] (studio) {\textbf{das Studio}\\ Radiostation};
\node[draw=popgreen, fill=popgreenL, rounded corners, minimum width=2.6cm, minimum height=1cm, align=center, right=of studio] (sprecher) {\textbf{der Sprecher}\\ liest den Text};
\node[draw=poporange, fill=poporangeL, rounded corners, minimum width=2.6cm, minimum height=1cm, align=center, right=of sprecher] (mikro) {\textbf{das Mikrofon}\\ sendet die Stimme};
\node[draw=poppurple, fill=poppurpleL, rounded corners, minimum width=2.6cm, minimum height=1cm, align=center, right=of mikro] (zuhoerer) {\textbf{die Zuhörer}\\ hören zu};
\draw[-{Stealth}, very thick, popdark] (studio) -- (sprecher);
\draw[-{Stealth}, very thick, popdark] (sprecher) -- (mikro);
\draw[-{Stealth}, very thick, popdark] (mikro) -- (zuhoerer);
\end{tikzpicture}
\end{center}
\legende{Le parcours du message : du studio aux auditeurs}
\end{popfigure}

\begin{aretenir}[Les mots-clés d'un communiqué]
Un communiqué radio commence souvent par une formule d'appel : \all{Liebe Zuhörer!} « Chers auditeurs ! » ou \all{Achtung, eine Durchsage!} « Attention, une annonce ! ». Il se termine par une formule de clôture : \all{Bleiben Sie auf unserer Frequenz!} « Restez sur notre fréquence ! ». Entre les deux : les faits, les dates, les heures — des phrases courtes et claires.
\end{aretenir}

\begin{exoresolu}[Compréhension globale et détaillée]
\textbf{Énoncé.} Réponds en français, puis complète la phrase allemande.
\begin{enumerate}
\item Quel événement la radio annonce-t-elle, et quand a-t-il lieu ?
\item Quelle conséquence pour les automobilistes ?
\item Complète : \all{Die Hauptstraße ist von 8 Uhr bis 20 Uhr für den Verkehr ...}
\end{enumerate}
\tcblower
\textbf{Solution détaillée.}
\begin{enumerate}
\item La radio annonce une grande fête culturelle (\all{ein großes Kulturfest}) qui a lieu samedi 14 octobre au centre-ville, de 10 h à 18 h. L'entrée est gratuite (\all{der Eintritt ist frei}).
\item Conséquence : la rue principale (\all{die Hauptstraße}) est fermée à la circulation de 8 h à 20 h ; les bus des lignes 4 et 7 passent par la rocade (\all{die Ringstraße}). Il faut prévoir plus de temps.
\item \all{Die Hauptstraße ist von 8 Uhr bis 20 Uhr für den Verkehr gesperrt.} « La rue principale est fermée à la circulation de 8 h à 20 h. » Le participe \all{gesperrt} vient du verbe \all{sperren} (fermer, bloquer).
\end{enumerate}
\end{exoresolu}

\begin{exercice}[À toi de jouer : hypothèses et repérage]
\begin{enumerate}
\item Relis le texte et souligne tous les mots transparents (au moins six). Traduis-les.
\item Relève les cinq informations chiffrées du texte (dates et heures).
\item Imagine : la radio de Douala annonce un match de football au stade de Bépanda. Note en allemand les réponses aux cinq questions (Wer? Was? Wann? Wo? Warum?) en quelques mots.
\end{enumerate}
\end{exercice}

\begin{corrige}[Exercice « À toi de jouer »]
\begin{enumerate}
\item Mots transparents possibles : \all{Radio} (radio), \all{Kulturfest} (fête culturelle), \all{Musikgruppen} (groupes de musique), \all{informieren} (informer), \all{Frequenz} (fréquence), \all{Musik} (musique), \all{Fußballspiel} (match de football), \all{Stadtpark} (parc municipal).
\item Les informations chiffrées : le 14 octobre ; début à 10 h ; fin à 18 h ; rue fermée de 8 h à 20 h ; informations chaque heure (\all{jede Stunde}) ; inscriptions jusqu'à vendredi.
\item Exemple de réponse : \all{Wer? — Radio Douala. Was? — ein Fußballspiel. Wann? — am Sonntag, um 16 Uhr. Wo? — im Stadion in Bépanda. Warum? — das Finale der Schulmeisterschaft.}
\end{enumerate}
\end{corrige}

Nous savons maintenant ce qu'est un communiqué radio et comment en repérer les informations essentielles. Dans la prochaine étape, nous approfondirons la compréhension du texte question par question, avant d'enrichir notre vocabulaire du monde des médias.

\coursec{Compréhension du texte}

Dans la section précédente, nous avons découvert le texte support : un \cle{communiqué radio} (\all{die Durchsage}, \all{die Mitteilung}) annonçant un festival de rue et ses conséquences sur la circulation. Nous l'avons lu et écouté une première fois, sans nous arrêter sur les détails. Il est temps maintenant de passer à la \cle{compréhension} proprement dite : d'abord globale (de quoi parle le texte ? qui parle à qui ?), puis détaillée (dates, horaires, mesures précises). C'est exactement la démarche exigée à l'examen : lire d'abord « en large », puis « en profondeur ».

\courssub{Relecture du texte support}

Avant de répondre aux questions, relisons attentivement le communiqué. Pour faciliter le repérage, les phrases sont numérotées : dans les réponses, on pourra ainsi citer précisément la phrase qui justifie chaque réponse.

\begin{textebox}[Radio Yaoundé — „Die Morgennachrichten“]
(1) Guten Morgen, liebe Hörerinnen und Hörer, hier ist Radio Yaoundé mit einer wichtigen Durchsage. (2) Am Samstag findet in der Hauptstraße unser großes Stadtfest statt. (3) Das Fest beginnt um 10 Uhr und endet um 22 Uhr. (4) Die Hauptstraße ist für den Verkehr von 8 Uhr bis 23 Uhr gesperrt. (5) Die Busse müssen deshalb eine andere Route nehmen: sie fahren über die Avenue Kennedy. (6) Jede Stunde informiert unser Reporter die Hörer über die Situation. (7) Für die Kinder gibt es Spiele und Musik, und am Abend spielt eine Band aus Douala. (8) Besuchen Sie das Fest mit Ihrer Familie! (9) Wir danken den Bürgern für ihr Verständnis und wünschen allen viel Spaß. (10) Und jetzt zurück zum Studio, zu den Nachrichten aus Kamerun und der Welt.
\end{textebox}

\textbf{Glossaire :} \all{die Durchsage, -n} : l'annonce, le communiqué ; \all{das Stadtfest, -e} : la fête de la ville ; \all{stattfinden (findet statt, fand statt, hat stattgefunden)} : avoir lieu ; \all{gesperrt} : fermé(e), bloqué(e) ; \all{der Verkehr} : la circulation ; \all{die Route, -n} : l'itinéraire ; \all{der Hörer, - / die Hörerin, -nen} : l'auditeur / l'auditrice ; \all{das Verständnis} : la compréhension, l'indulgence ; \all{das Studio, -s} : le studio.

\courssub{Compréhension globale}

La compréhension globale répond à trois questions simples : \textbf{de quoi parle le texte ? qui parle ? à qui ?}

\begin{itemize}
\item \textbf{Le sujet} : il s'agit d'un \cle{communiqué radio} qui annonce un festival de rue (\all{das Stadtfest}) à Yaoundé et informe les habitants des perturbations de circulation (rue fermée, bus déviés).
\item \textbf{L'émetteur} : un \all{Sprecher} (speaker) de \all{Radio Yaoundé} s'adresse au micro. On le voit dès la phrase (1) : \all{„hier ist Radio Yaoundé“}.
\item \textbf{Les destinataires} : les auditeurs, c'est-à-dire les habitants de la ville : \all{„liebe Hörerinnen und Hörer“} (phrase 1), puis \all{„die Bürger“} (phrase 9).
\item \textbf{La fonction du texte} : informer (date, horaires, rue fermée), mais aussi inviter (\all{„Besuchen Sie das Fest!“}, phrase 8) et remercier (phrase 9). C'est un texte à la fois \emph{informatif} et \emph{incitatif}, typique du service public radiophonique.
\end{itemize}

\begin{aretenir}[L'essentiel du texte en une phrase]
Radio Yaoundé annonce que le festival de la ville aura lieu samedi dans la rue principale, que celle-ci sera fermée à la circulation et que les bus seront déviés.
\end{aretenir}

\courssub{Compréhension détaillée : les questions}

Passons maintenant aux questions de détail. Avant de répondre, il faut bien distinguer les deux grands types de questions en allemand, car la place du verbe n'est pas la même.

\begin{propriete}[Questions ouvertes et questions fermées]
\begin{itemize}
\item Une \cle{question ouverte} commence par un \textbf{mot interrogatif} (\all{Wann? Wo? Wer? Was? Warum? Wie lange?}), suivi du verbe conjugué en 2\up{e} position : \all{Wann findet das Fest statt?} On ne peut pas répondre par oui ou non.
\item Une \cle{question fermée} (réponse \all{ja} ou \all{nein}) commence directement par le \textbf{verbe conjugué} en 1\up{re} position : \all{Ist die Straße gesperrt?} — \all{Ja, sie ist gesperrt.}
\end{itemize}
Contrairement au français, il n'y a ni « est-ce que », ni inversion compliquée avec tiret : le verbe vient simplement en premier.
\end{propriete}

\begin{exemplebox}[Observer avant de répondre]
\all{Das Fest beginnt um 10 Uhr.} « La fête commence à 10 heures. »\par
\all{Die Hauptstraße ist für den Verkehr gesperrt.} « La rue principale est fermée à la circulation. »\par
\all{Die Busse müssen eine andere Route nehmen.} « Les bus doivent prendre un autre itinéraire. »
\end{exemplebox}

\begin{methode}[Répondre à une question de compréhension en 4 étapes]
\begin{enumerate}
\item \textbf{Repérer le mot-clé} de la question (ex. : \all{Wann} $\rightarrow$ on cherche une date ou une heure ; \all{die Busse} $\rightarrow$ on cherche la phrase qui parle des bus).
\item \textbf{Localiser la phrase du texte} qui contient ce mot-clé ou un synonyme (noter son numéro).
\item \textbf{Extraire l'information} demandée, rien de plus.
\item \textbf{Reformuler} avec ses propres mots, en phrase complète, sans recopier intégralement le texte : on peut changer l'ordre des mots ou utiliser un pronom (\all{er, sie, es}).
\end{enumerate}
\end{methode}

Appliquons cette méthode aux trois questions du programme.

\begin{exoresolu}[Questions de compréhension détaillée]
Réponds par des phrases complètes : a) \all{Wann findet das Fest statt?} b) \all{Wie lange ist die Straße gesperrt?} c) \all{Was müssen die Busse machen?}
\tcblower
a) Mot-clé : \all{wann} $\rightarrow$ phrase (2). \textbf{Réponse :} \all{Das Fest findet am Samstag statt.} (La fête a lieu samedi.)\par
b) Mot-clé : \all{wie lange} (combien de temps) $\rightarrow$ phrase (4) : \all{von 8 Uhr bis 23 Uhr}. \textbf{Réponse :} \all{Die Straße ist von 8 Uhr bis 23 Uhr gesperrt, also fünfzehn Stunden.} (La rue est fermée de 8 h à 23 h, soit quinze heures.)\par
c) Mot-clé : \all{die Busse} $\rightarrow$ phrase (5). \textbf{Réponse :} \all{Sie müssen eine andere Route nehmen und fahren über die Avenue Kennedy.} (Ils doivent prendre un autre itinéraire et passent par l'avenue Kennedy.) Remarque : on a remplacé \all{die Busse} par le pronom \all{sie} pour éviter la répétition — c'est une bonne reformulation.
\end{exoresolu}

\courssub{Vrai ou faux : justifier par une citation}

L'exercice « vrai / faux » (\all{richtig / falsch}) est un classique de l'examen. La règle d'or : une réponse sans justification ne rapporte que la moitié des points. Il faut toujours \textbf{citer la phrase du texte} qui prouve la réponse.

\begin{exoresolu}[Richtig oder falsch ?]
a) \all{Das Fest endet um 20 Uhr.} b) \all{Die Hauptstraße ist für den Verkehr gesperrt.} c) \all{Am Abend spielt eine Band aus Bamenda.}
\tcblower
a) \textbf{Falsch.} Le texte dit : \all{„Das Fest beginnt um 10 Uhr und endet um 22 Uhr.“} (phrase 3). La fête se termine à 22 h, pas à 20 h.\par
b) \textbf{Richtig.} Preuve : \all{„Die Hauptstraße ist für den Verkehr von 8 Uhr bis 23 Uhr gesperrt.“} (phrase 4).\par
c) \textbf{Falsch.} Preuve : \all{„am Abend spielt eine Band aus Douala“} (phrase 7) — le groupe vient de Douala, pas de Bamenda.
\end{exoresolu}

\begin{attention}[Piège du vrai/faux]
Les énoncés « faux » reprennent souvent les mots exacts du texte en changeant \textbf{un seul détail} : une heure, un lieu, un nom de ville. Ne réponds jamais de mémoire : retourne au texte et compare mot à mot. Ici, \all{Douala} et non \all{Bamenda}, \all{22 Uhr} et non \all{20 Uhr}.
\end{attention}

\courssub{Tableau de repérage et marqueurs temporels}

Pour organiser ses réponses, on peut construire un \cle{tableau de repérage} : question, réponse trouvée dans le texte, numéro de la phrase.

\begin{center}
\begin{tabularx}{\linewidth}{|X|X|l|}
\hline
\rowcolor{popblueL} \textbf{Frage} & \textbf{Antwort im Text} & \textbf{Satz} \\
\hline
Wann findet das Fest statt? & am Samstag & (2) \\
\hline
Wann beginnt das Fest? & um 10 Uhr & (3) \\
\hline
Wie lange ist die Straße gesperrt? & von 8 Uhr bis 23 Uhr & (4) \\
\hline
Was müssen die Busse machen? & eine andere Route nehmen & (5) \\
\hline
Wer informiert die Hörer? & der Reporter, jede Stunde & (6) \\
\hline
\end{tabularx}
\end{center}

Ce tableau met en évidence les \cle{marqueurs temporels} du communiqué. Un communiqué radio doit être précis : l'auditeur ne peut pas « relire », donc les informations de temps sont répétées et clairement structurées.

\begin{center}
\begin{tabularx}{\linewidth}{|X|X|X|}
\hline
\rowcolor{popgreenL} \textbf{Marqueur} & \textbf{Fonction} & \textbf{Français} \\
\hline
\all{am Samstag} & jour & le samedi \\
\hline
\all{um 10 Uhr} & heure précise & à 10 heures \\
\hline
\all{von 8 Uhr bis 23 Uhr} & durée (début $\rightarrow$ fin) & de 8 h à 23 h \\
\hline
\all{jede Stunde} & fréquence & toutes les heures \\
\hline
\all{am Abend} & moment de la journée & le soir \\
\hline
\end{tabularx}
\end{center}

\begin{popfigure}
\begin{center}
\begin{tikzpicture}[mindmap, grow cyclic, every node/.style={concept, align=center, font=\small}, root concept/.append style={concept color=popblue, font=\bfseries}, level 1/.append style={level distance=3.4cm, sibling angle=90}]
\node[root concept, align=center]{Zeitmarker\\ im Text}
  child[concept color=poporange]{ node[align=center]{\all{am Samstag}\\ Tag} }
  child[concept color=popgreen]{ node[align=center]{\all{um 10 Uhr}\\ Uhrzeit} }
  child[concept color=poppurple]{ node[align=center]{\all{von ... bis}\\ Dauer} }
  child[concept color=poppink]{ node[align=center]{\all{jede Stunde}\\ Frequenz} };
\end{tikzpicture}
\end{center}
\legende{Carte mentale : les quatre types de marqueurs temporels du communiqué.}
\end{popfigure}

\begin{attention}[Le verbe \all{stattfinden} : particule séparable]
\all{stattfinden} (avoir lieu) est un verbe à \cle{particule séparable}. À la 2\up{e} position, la particule \all{statt} part en \textbf{fin de phrase} : \all{Das Fest findet am Samstag statt.} « La fête a lieu samedi. » Ne dites jamais \all{Das Fest stattfindet am Samstag} : c'est la faute la plus fréquente des francophones, car en français « avoir lieu » ne se sépare pas. Au \all{Perfekt}, la particule se colle au participe : \all{Das Fest hat stattgefunden.} « La fête a eu lieu. »
\end{attention}

\courssub{Reformulation : résumer le communiqué}

Résumer (\all{zusammenfassen}) signifie : garder uniquement les informations essentielles (qui ? quoi ? quand ? où ? quelles conséquences ?) et les exprimer \textbf{avec ses propres mots}, en 2 ou 3 phrases.

\begin{methode}[Résumer un communiqué radio]
\begin{enumerate}
\item Souligner dans le texte les 4 ou 5 informations essentielles (événement, date, lieu, mesures, conséquences).
\item Éliminer les détails secondaires (les jeux pour enfants, les remerciements, la formule de retour au studio).
\item Rédiger 2 ou 3 phrases en changeant la formulation : pronoms, synonymes, ordre des mots.
\end{enumerate}
\end{methode}

\begin{exemplebox}[Un résumé modèle]
\all{Radio Yaoundé informiert die Hörer über das Stadtfest: Es findet am Samstag in der Hauptstraße statt, von 10 Uhr bis 22 Uhr. Die Straße ist gesperrt und die Busse fahren über die Avenue Kennedy.}\par
« Radio Yaoundé informe les auditeurs au sujet de la fête de la ville : elle a lieu samedi dans la rue principale, de 10 h à 22 h. La rue est fermée et les bus passent par l'avenue Kennedy. »
\end{exemplebox}

\begin{exercice}[À toi de jouer]
1) Réponds en phrases complètes : \all{Wer informiert die Hörer jede Stunde? Was gibt es für die Kinder?}\par
2) Richtig oder falsch ? Justifie par une citation : \all{Die Straße ist nur zwei Stunden gesperrt.}\par
3) Résume le communiqué en deux phrases avec tes propres mots.
\end{exercice}

\begin{corrige}[Exercice : compréhension du texte]
1) \all{Der Reporter informiert die Hörer jede Stunde über die Situation. Für die Kinder gibt es Spiele und Musik.}\par
2) \textbf{Falsch.} Preuve : \all{„Die Hauptstraße ist für den Verkehr von 8 Uhr bis 23 Uhr gesperrt.“} — la rue est fermée pendant quinze heures, pas deux.\par
3) Exemple : \all{Am Samstag findet in Yaoundé ein Stadtfest statt. Die Hauptstraße ist gesperrt, und die Busse nehmen eine andere Route.}
\end{corrige}

\courssub{Dégagement du thème : la radio, service public d'information}

Au-delà des détails, ce texte illustre le rôle fondamental de la \cle{radio comme service public} : informer les citoyens rapidement, gratuitement et dans une langue claire. Le communiqué annonce un événement, prévient des perturbations (rue fermée, bus déviés), rassure (\all{„Wir danken den Bürgern für ihr Verständnis“}) et invite à participer. Au Cameroun comme en Allemagne, la radio reste un média très écouté : elle touche les habitants des villes comme des villages, au marché, dans les taxis ou les bendskins.

\begin{savaistu}[La radio en Allemagne et au Cameroun]
En Allemagne, les radios publiques (comme celles de l'ARD) sont financées par une redevance payée par les foyers : elles ont pour mission d'\emph{informer}, d'\emph{éduquer} et de \emph{divertir}, sans dépendre de la publicité. Au Cameroun, la CRTV et de nombreuses radios privées et communautaires jouent un rôle comparable : bulletins d'information, annonces de circulation, avis à la population. Le \all{Sprecher} allemand et l'animateur camerounais utilisent les mêmes formules : saluer les auditeurs, annoncer clairement, remercier, puis passer la parole : \all{„Und jetzt zurück zum Studio!“}
\end{savaistu}

\begin{experience}[Débat guidé]
En groupes de quatre, discutez en allemand avec des phrases simples : \all{Ist das Radio heute noch wichtig?} — Arguments possibles : \all{Das Radio ist schnell. / Das Radio ist gratis. / Viele Leute hören Radio im Bus. / Das Handy ist auch wichtig.} Chaque groupe présente ensuite sa conclusion à la classe en deux phrases.
\end{experience}

\begin{aretenir}[Bilan de la section]
Pour comprendre un communiqué radio : 1) identifier le sujet, l'émetteur et les destinataires ; 2) repérer les mots-clés des questions et localiser les phrases du texte ; 3) justifier les réponses vrai/faux par des citations ; 4) relever les marqueurs temporels (\all{am}, \all{um}, \all{von ... bis}, \all{jede Stunde}) ; 5) résumer avec ses propres mots. Le thème dégagé : la radio est un service public qui informe et guide les citoyens.
\end{aretenir}

\coursec{Lexique : les médias et la radio}

Dans la section précédente, nous avons lu et compris le communiqué radio de Radio Yaoundé. Pour bien parler de ce texte, le commenter et le résumer, il faut maintenant maîtriser son vocabulaire. Cette section organise les mots du thème « médias et radio » en quatre familles logiques : l'appareil, le message, les personnes et l'événement. Pour chaque mot, retiens toujours trois informations : le mot, son article, son pluriel.

\courssub{Observation : un texte pour découvrir le vocabulaire}

Lis d'abord ce court article de presse scolaire. Souligne mentalement tous les mots liés à la radio.

\begin{textebox}[Ein Tag beim Radio — Un jour à la radio]
Es ist sechs Uhr morgens in Yaoundé. In der Redaktion von Radio Mfoundi beginnt der Arbeitstag. Der Journalist Pierre Nkoulou liest zuerst die Nachrichten: eine Meldung über ein Fußballspiel, eine Information über den Verkehr in der Stadt und eine Mitteilung des Bürgermeisters.

Um sieben Uhr macht die Sprecherin eine Durchsage: „Liebe Zuhörer, heute Nachmittag findet auf dem Platz ein großes Fest statt. Die Schule veranstaltet ein Konzert. Die Straße ist ab zwölf Uhr gesperrt.“ Viele Zuhörer hören diese Durchsage im Radio, zu Hause oder im Taxi.

Der Sender sendet auf der Frequenz 98,4. Das Programm des Tages ist interessant: um acht Uhr die Nachrichten, um neun Uhr Musik, um zehn Uhr eine Sendung über die Gesundheit der Kinder. Der Sprecher informiert die Zuhörer über alle Veranstaltungen der Woche: ein Fest in der Kirche, eine Feier im Stadion und ein Konzert im Kulturzentrum.

Am Abend spricht die Sprecherin noch einmal mit den Journalisten der Redaktion. Sie besprechen das Programm von morgen. „Das Radio ist wichtig für die Menschen“, sagt sie. „Wir informieren sie jeden Tag über die Stadt, das Land und die Welt.“
\end{textebox}

\textbf{Glossaire}

\begin{center}
\begin{tabularx}{\linewidth}{|X|X|}
\hline
\rowcolor{popblueL} \textbf{Deutsch} & \textbf{Français} \\
\hline
die Redaktion (-en) & la rédaction \\
\hline
die Nachrichten (Pl.) & les informations (journal) \\
\hline
die Meldung (-en) & la nouvelle brève, le fait divers \\
\hline
die Mitteilung (-en) & la communication officielle, le communiqué \\
\hline
die Durchsage (-n) & l'annonce (micro, haut-parleur) \\
\hline
der Bürgermeister (-) & le maire \\
\hline
stattfinden (findet statt) & avoir lieu \\
\hline
veranstalten & organiser \\
\hline
gesperrt & fermé, bloqué (rue) \\
\hline
die Frequenz (-en) & la fréquence \\
\hline
die Sendung (-en) & l'émission \\
\hline
die Gesundheit & la santé \\
\hline
besprechen & discuter de, mettre au point \\
\hline
\end{tabularx}
\end{center}

\textbf{Questions de compréhension}

\begin{enumerate}
\item \all{Um wie viel Uhr beginnt der Arbeitstag in der Redaktion?} (À quelle heure commence la journée de travail à la rédaction ?)
\item \all{Was macht die Sprecherin um sieben Uhr?} (Que fait la speakerine à sept heures ?)
\item \all{Auf welcher Frequenz sendet der Sender?} (Sur quelle fréquence émet la station ?)
\item \all{Welche Veranstaltungen gibt es in dieser Woche?} (Quelles manifestations ont lieu cette semaine ?)
\end{enumerate}

Réponses attendues : 1. \all{Der Arbeitstag beginnt um sechs Uhr.} 2. \all{Sie macht eine Durchsage.} 3. \all{Der Sender sendet auf der Frequenz 98,4.} 4. \all{Es gibt ein Fest in der Kirche, eine Feier im Stadion und ein Konzert im Kulturzentrum.}

\courssub{Le champ lexical de la radio : l'appareil et l'émission}

\begin{definition}[Vocabulaire : das Radio und das Programm]
\begin{itemize}
\item \all{das Radio, -s} : la radio (l'appareil et le média) — \all{Ich höre Radio.} « J'écoute la radio. »
\item \all{der Sender, -} : la station de radio, l'émetteur — \all{Der Sender sendet jeden Tag.} « La station émet chaque jour. »
\item \all{die Frequenz, -en} : la fréquence — \all{Die Frequenz ist 98,4.} « La fréquence est 98,4. »
\item \all{das Programm, -e} : le programme — \all{Das Programm beginnt um acht Uhr.} « Le programme commence à huit heures. »
\item \all{die Sendung, -en} : l'émission — \all{Die Sendung ist sehr interessant.} « L'émission est très intéressante. »
\end{itemize}
\end{definition}

\begin{attention}[der Sender ou das Radio ?]
\all{das Radio} désigne l'appareil (et le média en général), \all{der Sender} désigne la station qui émet (Radio Mfoundi, CRTV Radio, Deutsche Welle). On dit \all{im Radio hören} « entendre à la radio », mais \all{beim Sender arbeiten} « travailler dans une station ». Ne traduis pas « station » par \all{die Station} : \all{die Station} signifie « la gare / l'étape » (faux ami).
\end{attention}

\courssub{Le champ lexical du message}

\begin{definition}[Vocabulaire : die Botschaft]
\begin{itemize}
\item \all{die Mitteilung, -en} : la communication, le communiqué (officiel) — \all{Der Bürgermeister gibt eine Mitteilung.} « Le maire publie un communiqué. »
\item \all{die Durchsage, -n} : l'annonce (au micro, au haut-parleur) — \all{Die Durchsage kommt um sieben Uhr.} « L'annonce passe à sept heures. »
\item \all{die Meldung, -en} : la nouvelle brève, le fait divers — \all{Eine Meldung über ein Fußballspiel.} « Une brève sur un match de football. »
\item \all{die Nachricht, -en} : la nouvelle, le message ; au pluriel \all{die Nachrichten} : les informations (le journal) — \all{Ich höre die Nachrichten.} « J'écoute les informations. »
\item \all{die Information, -en} : l'information, le renseignement — \all{Die Information ist wichtig.} « L'information est importante. »
\end{itemize}
\end{definition}

\begin{attention}[die Nachricht ≠ die Meldung]
\all{die Nachricht} (singulier) = une nouvelle, un message personnel (\all{Ich bekomme eine Nachricht.} « Je reçois un message. »). Au pluriel, \all{die Nachrichten} = les informations, le journal télévisé ou radiophonique. \all{die Meldung}, elle, est une nouvelle brève, un fait divers annoncé rapidement à la radio ou dans la presse. On ne dit donc pas \all{die Meldungen hören} pour « écouter les informations », mais \all{die Nachrichten hören}.
\end{attention}

\courssub{Le champ lexical des personnes}

\begin{definition}[Vocabulaire : die Personen]
\begin{itemize}
\item \all{der Sprecher, -} : le speaker, l'annonceur — \all{Der Sprecher liest die Nachrichten.} « Le speaker lit les informations. »
\item \all{die Sprecherin, -nen} : la speakerine — \all{Die Sprecherin macht eine Durchsage.} « La speakerine fait une annonce. »
\item \all{der Zuhörer, -} : l'auditeur — \all{Die Zuhörer hören das Programm.} « Les auditeurs écoutent le programme. »
\item \all{der Journalist, -en} : le journaliste (déclinaison faible : \all{den Journalisten, dem Journalisten}) — \all{Der Journalist schreibt einen Text.} « Le journaliste écrit un texte. »
\item \all{die Redaktion, -en} : la rédaction — \all{Die Redaktion arbeitet am Morgen.} « La rédaction travaille le matin. »
\end{itemize}
\end{definition}

\begin{savaistu}[Le féminin en -in]
En allemand, les noms de métiers et de personnes forment presque toujours le féminin avec le suffixe \all{-in} : \all{der Sprecher} → \all{die Sprecherin}, \all{der Journalist} → \all{die Journalistin}, \all{der Lehrer} → \all{die Lehrerin}, \all{der Zuhörer} → \all{die Zuhörerin}. Au pluriel, le féminin prend \all{-nen} : \all{die Sprecherinnen}. Attention à la prononciation : « chpré-ke-rinne ».
\end{savaistu}

\courssub{Le champ lexical des événements et les verbes du thème}

\begin{definition}[Vocabulaire : das Ereignis]
\begin{itemize}
\item \all{die Veranstaltung, -en} : la manifestation, l'événement organisé — \all{Die Veranstaltung beginnt um 15 Uhr.} « La manifestation commence à 15 heures. »
\item \all{das Fest, -e} : la fête, le festival — \all{Das Fest ist auf dem Platz.} « La fête a lieu sur la place. »
\item \all{die Feier, -n} : la cérémonie, la fête (officielle) — \all{Die Feier ist im Stadion.} « La cérémonie a lieu au stade. »
\end{itemize}
Verbes essentiels :
\begin{itemize}
\item \all{stattfinden} (verbe fort à particule séparable : \all{findet statt}, \all{fand statt}, \all{hat stattgefunden}) : avoir lieu — \all{Die Veranstaltung findet morgen statt.} « La manifestation a lieu demain. »
\item \all{ankündigen} (verbe à particule séparable : \all{kündigt an}, \all{hat angekündigt}) : annoncer — \all{Der Sprecher kündigt das Konzert an.} « Le speaker annonce le concert. »
\item \all{informieren} : informer — \all{Das Radio informiert die Zuhörer.} « La radio informe les auditeurs. »
\item \all{sperren} : fermer, bloquer (une rue) — \all{Die Polizei sperrt die Straße.} « La police ferme la rue. »
\end{itemize}
\end{definition}

\begin{exemplebox}[Exemples traduits]
\all{Der Sprecher macht eine Durchsage.} « Le speaker fait une annonce. »

\all{Die Veranstaltung findet morgen statt.} « La manifestation a lieu demain. »

\all{Die Sprecherin kündigt ein Konzert an.} « La speakerine annonce un concert. »

\all{Die Straße ist heute gesperrt.} « La rue est fermée aujourd'hui. »
\end{exemplebox}

\courssub{Carte mentale : quatre familles de mots}

\begin{popfigure}
\begin{tikzpicture}[mindmap, grow cyclic, every node/.style={concept, font=\small},
root concept/.append style={concept color=popblue, font=\small\bfseries},
level 1/.append style={level distance=3.4cm, sibling angle=90},
level 2/.append style={level distance=2.4cm, sibling angle=45, font=\scriptsize}]
\node[root concept]{RADIO UND MEDIEN}
  child[concept color=poporange]{ node {das Gerät}
    child { node {das Radio, -s} }
    child { node {der Sender, -} }
    child { node {die Frequenz, -en} }
  }
  child[concept color=popgreen]{ node {die Botschaft}
    child { node {die Mitteilung} }
    child { node {die Durchsage} }
    child { node {die Meldung} }
    child { node {die Nachricht} }
  }
  child[concept color=poppurple]{ node {die Personen}
    child { node {der Sprecher} }
    child { node {die Sprecherin} }
    child { node {der Zuhörer} }
    child { node {der Journalist} }
  }
  child[concept color=poppink]{ node {das Ereignis}
    child { node {die Veranstaltung} }
    child { node {das Fest, -e} }
    child { node {die Feier, -n} }
  };
\end{tikzpicture}
\legende{Carte mentale du vocabulaire : quatre familles autour du thème « radio et médias ».}
\end{popfigure}

\courssub{Les familles de mots : apprendre vite, retenir longtemps}

\begin{definition}[La famille de mots — die Wortfamilie]
Une \cle{famille de mots} regroupe des mots construits sur la même racine. Exemple : \all{melden} (annoncer) → \all{die Meldung} (la nouvelle brève) → \all{der Melder} (celui qui annonce). Connaître la racine permet de deviner le sens de mots nouveaux.
\end{definition}

Trois familles à connaître pour cette leçon :

\begin{itemize}
\item \all{sprechen} (parler) → \all{der Sprecher} (le speaker), \all{die Sprache} (la langue), \all{die Durchsage} (l'annonce), \all{besprechen} (discuter de).
\item \all{melden} (annoncer, signaler) → \all{die Meldung} (la nouvelle brève).
\item \all{veranstalten} (organiser) → \all{die Veranstaltung} (la manifestation).
\end{itemize}

\begin{popfigure}
\begin{tikzpicture}[
grow=right,
level distance=3.6cm,
level 1/.style={sibling distance=2.6cm},
level 2/.style={sibling distance=1.4cm},
every node/.style={draw=popblue, rounded corners, fill=popblueL, font=\small, align=center},
edge from parent/.style={draw=popdark, thick, -{Stealth}}]
\node[fill=poporangeL, draw=poporange, font=\small\bfseries, align=center]{sprechen\\ (parler)}
  child { node[align=center]{besprechen\\ discuter de} }
  child { node[align=center]{die Durchsage\\ l'annonce} }
  child { node[align=center]{die Sprache\\ la langue} }
  child { node[align=center]{der Sprecher\\ le speaker} };
\end{tikzpicture}
\legende{Arbre de la famille de mots de \all{sprechen} : une racine, quatre dérivés.}
\end{popfigure}

\courssub{Tableau récapitulatif du vocabulaire}

\begin{center}
\begin{tabularx}{\linewidth}{|l|l|X|}
\hline
\rowcolor{popblueL} \textbf{Deutsch (article, pluriel)} & \textbf{Français} & \textbf{Exemple} \\
\hline
das Radio, -s & la radio & \all{Ich höre Radio.} \\
\hline
der Sender, - & la station & \all{Der Sender sendet auf 98,4.} \\
\hline
die Frequenz, -en & la fréquence & \all{Die Frequenz ist gut.} \\
\hline
das Programm, -e & le programme & \all{Das Programm beginnt um 8 Uhr.} \\
\hline
die Sendung, -en & l'émission & \all{Die Sendung ist interessant.} \\
\hline
die Mitteilung, -en & le communiqué & \all{Der Bürgermeister gibt eine Mitteilung.} \\
\hline
die Durchsage, -n & l'annonce & \all{Die Sprecherin macht eine Durchsage.} \\
\hline
die Meldung, -en & la brève & \all{Eine Meldung über den Verkehr.} \\
\hline
die Nachricht, -en & la nouvelle ; les infos (Pl.) & \all{Ich höre die Nachrichten.} \\
\hline
die Information, -en & l'information & \all{Die Information ist wichtig.} \\
\hline
der Sprecher, - / die Sprecherin, -nen & le speaker / la speakerine & \all{Der Sprecher liest die Nachrichten.} \\
\hline
der Zuhörer, - & l'auditeur & \all{Die Zuhörer rufen an.} \\
\hline
der Journalist, -en & le journaliste & \all{Der Journalist schreibt einen Text.} \\
\hline
die Redaktion, -en & la rédaction & \all{Die Redaktion arbeitet am Morgen.} \\
\hline
die Veranstaltung, -en & la manifestation & \all{Die Veranstaltung findet statt.} \\
\hline
das Fest, -e & la fête & \all{Das Fest ist auf dem Platz.} \\
\hline
die Feier, -n & la cérémonie & \all{Die Feier ist im Stadion.} \\
\hline
\end{tabularx}
\end{center}

\courssub{Expressions utiles à réutiliser}

\begin{aretenir}[Trois expressions à savoir par cœur]
\begin{itemize}
\item \all{eine Durchsage machen} : faire une annonce — \all{Der Sprecher macht eine Durchsage.} « Le speaker fait une annonce. »
\item \all{über etwas informieren} : informer sur quelque chose (avec l'accusatif) — \all{Das Radio informiert die Zuhörer über das Fest.} « La radio informe les auditeurs sur la fête. »
\item \all{im Radio hören} : entendre à la radio — \all{Ich höre die Nachrichten im Radio.} « J'écoute les informations à la radio. »
\end{itemize}
\end{aretenir}

\begin{methode}[Comment mémoriser le vocabulaire allemand]
\begin{enumerate}
\item Apprends chaque nom avec son article et son pluriel, jamais le nom seul : \all{das Radio, -s} ; \all{die Meldung, -en} ; \all{der Sprecher, -}.
\item Regroupe les mots par famille logique (appareil, message, personne, événement) : la carte mentale t'aide à réviser.
\item Apprends les mots par familles de mots : si tu connais \all{sprechen}, tu devines \all{der Sprecher} et \all{die Sprache}.
\item Fabrique une phrase avec chaque mot nouveau : un mot utilisé est un mot retenu.
\item Révise à voix haute : la prononciation fixe la mémoire (\all{Sprecher} se prononce « chpré-heur »).
\end{enumerate}
\end{methode}

\begin{exoresolu}[Complète avec le mot qui convient]
Énoncé. Complète les phrases avec : \all{die Nachrichten}, \all{eine Durchsage}, \all{ein Sender}, \all{die Veranstaltung}, \all{der Zuhörer}.

1. \all{Ich höre jeden Morgen \_\_\_\_\_\_\_\_ im Radio.}

2. \all{Die Sprecherin macht \_\_\_\_\_\_\_\_ über das Fest.}

3. \all{Radio Mfoundi ist \_\_\_\_\_\_\_\_ in Yaoundé.}

4. \all{\_\_\_\_\_\_\_\_ findet am Samstag statt.}

5. \all{\_\_\_\_\_\_\_\_ ruft beim Radio an.}
\tcblower
Solution détaillée.

1. \all{Ich höre jeden Morgen \textbf{die Nachrichten} im Radio.} — « les informations » se dit toujours au pluriel : \all{die Nachrichten} (accusatif pluriel, identique au nominatif).

2. \all{Die Sprecherin macht \textbf{eine Durchsage} über das Fest.} — \all{eine Durchsage machen} : « faire une annonce » ; \all{Durchsage} est féminin, l'article indéfini reste \all{eine} à l'accusatif.

3. \all{Radio Mfoundi ist \textbf{ein Sender} in Yaoundé.} — \all{der Sender} est masculin ; après le verbe \all{sein}, le nom reste au nominatif : \all{ein Sender} (la station).

4. \all{\textbf{Die Veranstaltung} findet am Samstag statt.} — \all{stattfinden} se combine avec \all{die Veranstaltung} : « la manifestation a lieu samedi » ; la particule \all{statt} va en fin de phrase.

5. \all{\textbf{Der Zuhörer} ruft beim Radio an.} — celui qui écoute et téléphone à la station, c'est l'auditeur ; \all{anrufen} est aussi un verbe à particule séparable : \all{ruft ... an}.
\end{exoresolu}

\begin{attention}[Erreurs fréquentes des francophones]
\begin{itemize}
\item Oublier la majuscule : tous les noms allemands prennent une majuscule : \all{das Radio}, \all{die Meldung}, jamais \all{das radio}.
\item Confondre les articles : \all{das Radio} (neutre), \all{der Sender} (masculin), \all{die Sendung} (féminin). Un mot français féminin n'est pas forcément féminin en allemand : « la station » = \all{der Sender}.
\item Traduire « avoir lieu » mot à mot : on dit \all{stattfinden}, jamais \all{haben Ort}. \all{Das Fest findet morgen statt.} « La fête a lieu demain. »
\item Oublier la particule séparable : \all{Der Sprecher kündigt das Konzert \textbf{an}.} La particule \all{an-} va à la fin de la phrase.
\end{itemize}
\end{attention}

\begin{exercice}[À toi de jouer]
1. Donne l'article et le pluriel de : \all{Radio}, \all{Meldung}, \all{Sprecher}, \all{Fest}, \all{Frequenz}.

2. Traduis en allemand : a) Le speaker fait une annonce. b) La manifestation a lieu demain. c) J'écoute les informations à la radio. d) La radio informe les auditeurs sur le concert.

3. Trouve le mot de la même famille : \all{veranstalten} → ? ; \all{melden} → ? ; \all{sprechen} → ? (deux réponses).
\end{exercice}

\begin{corrige}[Exercice « À toi de jouer »]
1. \all{das Radio, -s} ; \all{die Meldung, -en} ; \all{der Sprecher, -} ; \all{das Fest, -e} ; \all{die Frequenz, -en}.

2. a) \all{Der Sprecher macht eine Durchsage.} b) \all{Die Veranstaltung findet morgen statt.} c) \all{Ich höre die Nachrichten im Radio.} d) \all{Das Radio informiert die Zuhörer über das Konzert.}

3. \all{veranstalten} → \all{die Veranstaltung} ; \all{melden} → \all{die Meldung} ; \all{sprechen} → \all{der Sprecher}, \all{die Sprache} (aussi \all{die Durchsage}, \all{besprechen}).
\end{corrige}

Tu possèdes maintenant tout le vocabulaire du thème « médias et radio », organisé en familles et mémorisable grâce à la carte mentale. La prochaine section te montrera comment employer ces mots dans des phrases correctes : nous étudierons l'accusatif et le datif, les deux cas qui décident de la forme de l'article après le verbe.

\coursec{Grammaire 1 : L'accusatif et le datif}

Après avoir exploré le vocabulaire spécifique de la radio dans la section précédente (\emph{der Sprecher, die Durchsage, die Meldung}), nous devons maintenant comprendre comment ces mots s'articulent dans la phrase. En allemand, la fonction d'un nom (sujet, objet direct, objet indirect) ne se devine pas à sa place dans la phrase comme en français, mais à sa \textbf{déclinaison} : c'est le système des \cle{cas} (Kasus). Maîtriser l'accusatif et le datif est indispensable pour comprendre un communiqué radio et pour produire vos propres phrases correctes.

\courssub{Le système des cas : rappel du nominatif}

En allemand, tout nom (ou groupe nominal) porte une marque de \cle{cas} qui indique son rôle grammatical. Il y a quatre cas, mais le programme de Première se concentre sur les trois premiers.

\begin{definition}[Les quatre cas allemands]
\begin{itemize}
    \item \textbf{Nominativ (1. Fall) :} Le sujet de la phrase (qui/quoi fait l'action ? \emph{Wer/Was?}).
    \item \textbf{Genitiv (2. Fall) :} Le complément du nom (dont ? \emph{Wessen?}) — approfondi en Terminale.
    \item \textbf{Dativ (3. Fall) :} Le complément d'objet indirect / attributif (à qui ? \emph{Wem?}).
    \item \textbf{Akkusativ (4. Fall) :} Le complément d'objet direct (qui/quoi ? \emph{Wen/Was?}).
\end{itemize}
\end{definition}

\emph{Observation :} Dans le communiqué \all{„Der Sprecher informiert die Zuhörer.“} (Le speaker informe les auditeurs), \all{Der Sprecher} est au nominatif (sujet), \all{die Zuhörer} est à l'accusatif (objet direct). En français, l'ordre Sujet-Verbe-Objet suffit ; en allemand, c'est l'article \all{der} vs \all{den} qui le montre.

\courssub{L'accusatif : le complément d'objet direct (COD)}

L'accusatif marque l'objet \textbf{directement} atteint par l'action du verbe (verbes transitifs directs : \all{sehen, hören, lesen, informieren, erwarten, besuchen, haben}).

\begin{propriete}[L'accusatif : forme des articles définis, indéfinis et possessifs]
\emph{Règle d'or :} \textbf{Seul le masculin singulier change de forme} par rapport au nominatif.
\begin{center}
\begin{tabular}{|l|c|c|c|c|}
\hline
\rowcolor{popblueL}
\textbf{Cas} & \textbf{Masculin} & \textbf{Féminin} & \textbf{Neutre} & \textbf{Pluriel} \\
\hline
Nominativ & der / ein / mein & die / eine / meine & das / ein / mein & die / \textendash{} / meine \\
\hline
Akkusativ & \textbf{den / einen / meinen} & die / eine / meine & das / ein / mein & die / \textendash{} / meine \\
\hline
\end{tabular}
\end{center}
Les possessifs suivent le même schéma que l'article indéfini : \all{mein} $\to$ \all{meinen} (masc.), \all{meine} (fém./pl.), \all{mein} (neutre).
\end{propriete}

\begin{exemplebox}[Accusatif : exemples traduits]
\begin{itemize}
    \item \all{Ich höre den Sprecher.} « J'écoute le speaker. » (masc. : \all{der} $\to$ \all{den})
    \item \all{Wir lesen die Meldung.} « Nous lisons la dépêche. » (fém. : \all{die} inchangé)
    \item \all{Er sieht das Studio.} « Il voit le studio. » (neutre : \all{das} inchangé)
    \item \all{Sie erwarten die Gäste.} « Ils attendent les invités. » (pluriel : \all{die} inchangé)
    \item \all{Ich habe einen Radioapparat.} « J'ai un poste de radio. » (indéfini masc. : \all{ein} $\to$ \all{einen})
    \item \all{Sie begrüßt ihren Onkel.} « Elle salue son oncle. » (possessif masc. : \all{ihr} $\to$ \all{ihren})
\end{itemize}
\end{exemplebox}

\begin{methode}[Test de l'accusatif : \emph{Wen oder Was?}]
Pour vérifier si un complément est à l'accusatif, posez la question \all{Wen?} (qui, pour les personnes) ou \all{Was?} (quoi, pour les choses) au verbe.
\begin{enumerate}
    \item Trouver le verbe conjugué (\all{informiert}).
    \item Demander : \all{Wen informiert der Sprecher?} $\to$ \all{die Zuhörer} (Accusatif).
\end{enumerate}
\end{methode}

\courssub{Le datif : le complément d'objet indirect (COI)}

Le datif marque le \textbf{bénéficiaire} ou le \textbf{destinataire} de l'action (souvent « à qui ? » en français). C'est le cas des verbes « à double objet » (donner qqch \textbf{à qqn}) et de nombreux verbes allemands qui régissent le datif de façon lexicale.

\begin{propriete}[Le datif : forme des articles définis, indéfinis et possessifs]
Au datif, \textbf{tous les genres changent} et le \textbf{pluriel ajoute un \emph{-n} au nom} (s'il ne finit pas déjà par \emph{-n/-s}).
\begin{center}
\begin{tabular}{|l|c|c|c|c|}
\hline
\rowcolor{popgreenL}
\textbf{Cas} & \textbf{Masculin} & \textbf{Féminin} & \textbf{Neutre} & \textbf{Pluriel} \\
\hline
Nominativ & der / ein / mein & die / eine / meine & das / ein / mein & die Kinder / meine Kinder \\
\hline
Dativ & \textbf{dem / einem / meinem} & \textbf{der / einer / meiner} & \textbf{dem / einem / meinem} & \textbf{den Kindern / meinen Kindern} \\
\hline
\end{tabular}
\end{center}
\emph{Attention :} \all{die} (fém. sing. Nom/Akk) devient \all{der} au datif (homographie avec masc. nominatif). Le pluriel datif est toujours \all{den} + \emph{-n} au nom (\all{den Zuhörern, den Sendungen, meinen Freunden}).
\end{propriete}

\begin{exemplebox}[Datif : exemples traduits]
\begin{itemize}
    \item \all{Ich helfe dem Sprecher.} « J'aide le speaker. » (masc. : \all{der} $\to$ \all{dem})
    \item \all{Wir danken der Moderatorin.} « Nous remercions l'animatrice. » (fém. : \all{die} $\to$ \all{der})
    \item \all{Das gefällt dem Kind.} « Cela plaît à l'enfant. » (neutre : \all{das} $\to$ \all{dem})
    \item \all{Wir antworten den Zuhörern.} « Nous répondons aux auditeurs. » (pluriel : \all{die Zuhörer} $\to$ \all{den Zuhörern})
    \item \all{Ich gebe einem Freund ein Radio.} « Je donne une radio à un ami. » (indéfini masc. : \all{ein} $\to$ \all{einem})
    \item \all{Er schreibt seiner Mutter.} « Il écrit à sa mère. » (possessif fém. : \all{seine} $\to$ \all{seiner})
\end{itemize}
\end{exemplebox}

\begin{methode}[Test du datif : \emph{Wem?}]
Posez la question \all{Wem?} (à qui ?) au verbe.
\begin{enumerate}
    \item Verbe : \all{helfen}.
    \item Question : \all{Wem helfe ich?} $\to$ \all{dem Sprecher} (Datif).
\end{enumerate}
\end{methode}

\courssub{Verbes à construction datif (Verben mit Dativ)}

C'est un piège majeur pour les francophones : certains verbes sont transitifs \textbf{directs} en français (COD) mais transitifs \textbf{indirects} en allemand (Datif).

\begin{aretenir}[Verbes fréquents régissant le datif]
\begin{center}
\begin{tabularx}{\linewidth}{|X|X|X|}
\hline
\rowcolor{poppurpleL}
\textbf{Allemand} & \textbf{Français (souvent COD)} & \textbf{Exemple (Datif)} \\
\hline
\all{helfen} & aider & \all{Ich helfe \textbf{dir}.} (Je t'aide.) \\
\all{danken} & remercier & \all{Wir danken \textbf{Ihnen}.} (Nous vous remercions.) \\
\all{antworten} & répondre & \all{Er antwortet \textbf{mir}.} (Il me répond.) \\
\all{gefallen} & plaire & \all{Das gefällt \textbf{mir}.} (Cela me plaît.) \\
\all{folgen} & suivre (obéir) & \all{Wir folgen \textbf{ihm}.} (Nous le suivons / obéissons.) \\
\all{glauben} & croire (qqn) & \all{Ich glaube \textbf{dir}.} (Je te crois.) \\
\all{gratulieren} & féliciter & \all{Ich gratuliere \textbf{dir}.} (Je te félicite.) \\
\all{gehören} & appartenir (à qqn) & \all{Das Radio gehört \textbf{mir}.} (La radio m'appartient.) \\
\hline
\end{tabularx}
\end{center}
\end{aretenir}

\begin{attention}[Piège n°1 : \emph{helfen} vs \emph{aider}]
En français : \emph{J'aide \textbf{mon frère}} (COD). En allemand : \all{Ich helfe \textbf{meinem Bruder}} (Datif : \all{mein\textbf{em}}). \textbf{Ne dites jamais} \all{*ich helfe meinen Bruder} (accusatif incorrect).
\end{attention}

\begin{savaistu}[Le cas de \emph{informieren} et \emph{fragen}]
\all{informieren} est un verbe transitif direct (accusatif) pour la personne : \all{Ich informiere \textbf{den Zuhörer}}. Mais pour le \textbf{sujet} de l'information, on utilise une préposition : \all{informieren \textbf{über} + Akkusativ} (\all{Ich informiere ihn \textbf{über das Fest}}). Double accusatif ! \\
De même, \all{fragen} (demander) prend l'accusatif pour la personne (\all{Ich frage \textbf{ihn}}) et une préposition pour le sujet : \all{fragen \textbf{nach} + Dativ} (\all{Ich frage ihn \textbf{nach dem Weg}}).
\end{savaistu}

\courssub{Prépositions à cas fixe (Präpositionen mit festem Kasus)}

Certaines prépositions imposent \textbf{toujours} le même cas, quel que soit le contexte (mouvement ou position). Il faut les apprendre par cœur.

\begin{propriete}[Prépositions + Datif \emph{vs} + Accusatif]
\begin{center}
\begin{tabularx}{\linewidth}{|X|X|X|}
\hline
\rowcolor{poporangeL}
\textbf{+ DATIF (9 prépositions)} & \textbf{+ ACCUSATIF (5 prépositions)} & \textbf{Mnémotechnique} \\
\hline
\all{mit} (avec) & \all{f\"ur} (pour) & \emph{« Mit dem Bus f\"ur den Vater.»} \\
\all{von} (de, depuis) & \all{durch} (par, à travers) & \\
\all{zu} (chez, vers) & \all{gegen} (contre, vers) & \\
\all{nach} (après, vers [pays/villes sans art.]) & \all{ohne} (sans) & \\
\all{bei} (chez, près de, au travail) & \all{um} (autour de, pour [heure]) & \\
\all{seit} (depuis - temps) & & \\
\all{aus} (hors de, en provenance de) & & \\
\all{gegen\"uber} (en face de) & & \\
\all{außer} (sauf) & & \\
\hline
\end{tabularx}
\end{center}
\end{propriete}

\begin{exemplebox}[Prépositions en contexte radio]
\begin{itemize}
    \item \all{Wir hören \textbf{mit} dem Radio.} (Datif : \all{dem}) — « Nous écoutons avec la radio. »
    \item \all{Die Sendung ist \textbf{f\"ur} \textbf{die} Jugend.} (Accusatif : \all{die}) — « L'émission est pour la jeunesse. »
    \item \all{Der Reporter kommt \textbf{aus} dem Studio.} (Datif : \all{dem}) — « Le reporter vient du studio. »
    \item \all{Er geht \textbf{ohne} \textbf{den} Kopfh\"orer.} (Accusatif : \all{den}) — « Il part sans le casque. »
    \item \all{Seit \textbf{dem} Morgen sendet der Sender.} (Datif : \all{dem}) — « Depuis ce matin, la station émet. »
    \item \all{Wir fahren \textbf{zu} \textbf{dem} (=\textbf{zum}) Sender.} (Datif : \all{dem}) — « Nous allons à la station. »
\end{itemize}
\end{exemplebox}

\courssub{Prépositions à deux cas : \emph{in, an, auf} (Wechselpräpositionen)}

Ces prépositions indiquent un \textbf{lieu}. Le cas dépend de la question qu'on leur pose : \textbf{position statique (Wo?) $\to$ Datif} / \textbf{mouvement vers un but (Wohin?) $\to$ Accusatif}.

\begin{propriete}[Règle \emph{Wo?} (Datif) / \emph{Wohin?} (Accusatif)]
\begin{itemize}
    \item \textbf{Wo? (Où ? Position) :} \all{in} + \textbf{Datif}, \all{an} + \textbf{Datif}, \all{auf} + \textbf{Datif}.
    \item \textbf{Wohin? (Où vers ? Mouvement) :} \all{in} + \textbf{Accusatif}, \all{an} + \textbf{Accusatif}, \all{auf} + \textbf{Accusatif}.
\end{itemize}
\end{propriete}

\begin{exemplebox}[Alternance \emph{in / an / auf}]
\begin{tabular}{lll}
\hline
\rowcolor{poppinkL}
\textbf{Allemand} & \textbf{Cas} & \textbf{Français} \\
\hline
\all{Wir sind \textbf{in der} Stadt.} & Dativ (Wo?) & Nous sommes \textbf{dans} la ville. \\
\all{Wir fahren \textbf{in die} Stadt.} & Akkusativ (Wohin?) & Nous allons \textbf{dans} la ville. \\
\all{Das Bild h\"angt \textbf{an der} Wand.} & Dativ (Wo?) & Le tableau pend \textbf{au} mur. \\
\all{Er h\"angt das Bild \textbf{an die} Wand.} & Akkusativ (Wohin?) & Il accroche le tableau \textbf{au} mur. \\
\all{Das Buch liegt \textbf{auf dem} Tisch.} & Dativ (Wo?) & Le livre est \textbf{sur} la table. \\
\all{Er legt das Buch \textbf{auf den} Tisch.} & Akkusativ (Wohin?) & Il pose le livre \textbf{sur} la table. \\
\hline
\end{tabular}
\end{exemplebox}

\begin{popfigure}
\begin{tikzpicture}[
    node distance=1.2cm and 1.5cm,
    box/.style={draw=popdark, thick, rounded corners, fill=popcream, text width=3.2cm, align=center, font=\small},
    arrow/.style={-Stealth, thick, popblue},
    label/.style={font=\tiny, text=popmuted, above, sloped}
]
% Nodes
\node[box, align=center] (nom){\textbf{Nominativ (Sujet)}\\\all{Der Sprecher}};
\node[box, right=of nom, align=center] (akk){\textbf{Akkusativ (COD)}\\\all{die Zuhörer}};
\node[box, below=of akk, align=center] (dat){\textbf{Dativ (COI / Prépos.)}\\\all{dem Fest}};

% Arrows
\draw[arrow] (nom) -- node[label] {informiert} (akk);
\draw[arrow] (akk) -- node[label, right] {\all{über} + Akk.} (dat);

% Preposition info
\node[draw=poporange, thick, rounded corners, fill=poporange!10, font=\small, align=center, below=0.5cm of dat, text width=5cm] (prep) {\all{informieren jdn (Akk) }\\\all{über etw (Akk)}};
\end{tikzpicture}
\legende{Schéma de structure : \all{Der Sprecher (Nom) informiert die Zuhörer (Akk) über das Fest (Akk nach Präp.)}.}
\end{popfigure}

\courssub{Les prépositions de temps (Temporale Präpositionen)}

\emph{Point de programme : Grammaire 2.} Pour situer un événement dans le temps (horaire d'une émission, date d'un reportage, durée), l'allemand utilise des prépositions qui régissent le datif ou l'accusatif.

\begin{propriete}[Prépositions de temps : cas et emploi]
\begin{center}
\begin{tabularx}{\linewidth}{|X|X|X|X|}
\hline
\rowcolor{popblueL}
\textbf{Préposition} & \textbf{Cas} & \textbf{Emploi} & \textbf{Exemple} \\
\hline
\all{an} & Dativ & Jour de la semaine, partie de la journée & \all{am Montag, am Morgen} \\
\all{in} & Dativ & Mois, saison, année, moment de la journée (sauf nuit) & \all{im Januar, im Sommer, im Jahr 2025, nachmittags} \\
\all{um} & Dativ & Heure précise & \all{um 14.30 Uhr} \\
\all{seit} / \all{seither} & Dativ & Depuis (point de départ dans le passé, action qui dure) & \all{seit Montag, seit zwei Stunden} \\
\all{ab} & Dativ & À partir de (futur ou présent) & \all{ab nächstem Monat} \\
\all{von ... bis} & Dativ & De ... à ... (durée bornée) & \all{von 10 bis 12 Uhr} \\
\all{vor} & Dativ & Il y a (durée écoulée, action terminée) & \all{vor einer Woche} \\
\all{f\"ur} & Akkusativ & Pour (durée future) & \all{f\"ur zwei Tage} \\
\all{innerhalb} / \all{außerhalb} & Genitiv (formel) / Dativ (parlé) & Dans un délai de / hors délai & \all{innerhalb einer Stunde} \\
\hline
\end{tabularx}
\end{center}
\end{propriete}

\begin{exemplebox}[Prépositions de temps dans un communiqué]
\begin{itemize}
    \item \all{Die Sendung beginnt \textbf{um 20.00 Uhr}.} (Heure exacte)
    \item \all{Wir senden \textbf{am Montag} \textbf{im Studio}.} (Jour + Lieu)
    \item \all{Der Reporter arbeitet \textbf{seit 2010} hier.} (Depuis 2010, encore maintenant)
    \item \all{Das Interview dauert \textbf{f\"ur zehn Minuten}.} (Durée prévue : pour 10 minutes)
    \item \all{Die Meldung kam \textbf{vor einer Stunde}.} (Il y a une heure)
    \item \all{Wir sind \textbf{von 8 bis 16 Uhr} im Dienst.} (De 8h à 16h)
\end{itemize}
\end{exemplebox}

\begin{attention}[Contractions obligatoires (Artikelverschmelzung)]
À l'oral comme à l'écrit courant, les prépositions \textbf{contractent} l'article défini suivant (Datif ou Accusatif) :
\begin{center}
\begin{tabular}{|l|l|l|}
\hline
\rowcolor{popgoldL}
\textbf{Préposition + Article} & \textbf{Forme contractée} & \textbf{Exemple} \\
\hline
\all{an + dem} & \all{am} & \all{am Montag, am Morgen} \\
\all{in + dem} & \all{im} & \all{im Januar, im Studio} \\
\all{in + das} & \all{ins} & \all{ins Studio (Wohin?)} \\
\all{zu + dem} & \all{zum} & \all{zum Sender} \\
\all{zu + der} & \all{zur} & \all{zur Sendung} \\
\all{von + dem} & \all{vom} & \all{vom 1. bis zum 3. Mai} \\
\all{bei + dem} & \all{beim} & \all{beim Warten} \\
\hline
\end{tabular}
\end{center}
\textbf{Jamais} \all{*an dem Montag} (sauf insistance démonstrative), toujours \all{am Montag}.
\end{attention}

\courssub{Récapitulatif visuel : Tableau des articles (Nominativ, Akkusativ, Dativ)}

Le tableau suivant synthétise les formes de l'article défini et indéfini pour les trois cas vus jusqu'ici. C'est votre « pense-bête » pour les exercices.

\begin{center}
\begin{tabular}{|l|c|c|c|c|}
\hline
\rowcolor{popblueL}
\textbf{Cas / Genre} & \textbf{Maskulin} & \textbf{Feminin} & \textbf{Neutrum} & \textbf{Plural} \\
\hline
\rowcolor{popblueL}
\textbf{Nominativ (def.)} & der & die & das & die \\
\hline
\textbf{Akkusativ (def.)} & \textbf{den} & die & das & die \\
\hline
\rowcolor{popblueL}
\textbf{Dativ (def.)} & \textbf{dem} & \textbf{der} & \textbf{dem} & \textbf{den} + \emph{-n} \\
\hline
\hline
\rowcolor{popgreenL}
\textbf{Nominativ (indef.)} & ein & eine & ein & keine \\
\hline
\textbf{Akkusativ (indef.)} & \textbf{einen} & eine & ein & keine \\
\hline
\rowcolor{popgreenL}
\textbf{Dativ (indef.)} & \textbf{einem} & \textbf{einer} & \textbf{einem} & \textbf{keinen} + \emph{-n} \\
\hline
\end{tabular}
\end{center}
\emph{En gras : les formes qui changent par rapport au nominatif. Notez l'identité Feminin Nominativ/Akkusativ (\all{die}/\all{eine}) et Neutre Nominativ/Akkusativ (\all{das}/\all{ein}).}

\courssub{Exercice résolu : Identifier les cas dans un communiqué}

\begin{exoresolu}[Analyse de cas]
\textbf{Énoncé :} Lisez les phrases extraites d'un communiqué radio. Pour chaque groupe nominal en \textbf{gras}, indiquez le \textbf{cas} (Nominativ, Akkusativ, Dativ) et justifiez par la question (\emph{Wer? Wen? Wem?} ou préposition).

\begin{enumerate}
    \item \all{\textbf{Der Moderator} begrüßt \textbf{die Zuhörer}.}
    \item \all{Er informiert \textbf{die Zuhörer} \textbf{über das Wetter}.}
    \item \all{\textbf{Dem Reporter} gefällt \textbf{die Sendung}.}
    \item \all{Wir danken \textbf{dem Team} \textbf{für die Arbeit}.}
    \item \all{Die Kinder sitzen \textbf{im} (=\textbf{in dem}) \textbf{Studio}.}
    \item \all{Der Techniker geht \textbf{ins} (=\textbf{in das}) \textbf{Studio}.}
    \item \all{Die Sendung beginnt \textbf{um 20.00 Uhr}.}
    \item \all{Wir hören \textbf{seit einer Stunde} zu.}
\end{enumerate}

\tcblower
\textbf{Solution détaillée :}

\begin{enumerate}
    \item \all{\textbf{Der Moderator}} : \textbf{Nominativ} (Sujet : \emph{Wer begrüßt?}). \\
          \all{\textbf{die Zuhörer}} : \textbf{Akkusativ} (COD de \emph{begrüßen} : \emph{Wen begrüßt er?}).
    \item \all{\textbf{die Zuhörer}} : \textbf{Akkusativ} (COD de \emph{informieren} : \emph{Wen informiert er?}). \\
          \all{\textbf{über das Wetter}} : \textbf{Akkusativ} (Préposition \all{über} impose l'accusatif). \emph{Note : \all{das} ne change pas au neutre.}
    \item \all{\textbf{Dem Reporter}} : \textbf{Dativ} (Objet logique de \emph{gefallen} : \emph{Wem gefällt?}). \\
          \all{\textbf{die Sendung}} : \textbf{Nominativ} (Sujet grammatical : \emph{Was gefällt?}).
    \item \all{\textbf{dem Team}} : \textbf{Dativ} (Verbe \emph{danken} + Dativ : \emph{Wem danken wir?}). \\
          \all{\textbf{für die Arbeit}} : \textbf{Akkusativ} (Préposition \all{für} + Akk).
    \item \all{\textbf{im Studio}} = \all{\textbf{in dem Studio}} : \textbf{Dativ} (Question \emph{Wo?} position statique : \emph{Wo sitzen sie?}). \emph{Article contracté \all{in dem} $\to$ \all{im}.}
    \item \all{\textbf{ins Studio}} = \all{\textbf{in das Studio}} : \textbf{Akkusativ} (Question \emph{Wohin?} mouvement : \emph{Wohin geht er?}). \emph{Article contracté \all{in das} $\to$ \all{ins}.}
    \item \all{\textbf{um 20.00 Uhr}} : \textbf{Dativ} (Préposition \all{um} + heure précise $\to$ Dativ).
    \item \all{\textbf{seit einer Stunde}} : \textbf{Dativ} (Préposition \all{seit} + Dativ).
\end{enumerate}
\end{exoresolu}

\courssub{Exercice d'application : Compléter les articles}

\begin{exercice}[Mettez l'article correct (défini ou indéfini) au cas convenable]
Complétez les phrases du communiqué suivant avec la forme correcte de l'article entre parenthèses. Attention aux contractions !

\begin{enumerate}
    \item \all{\_\_\_\_\_\_\_\_ (der) Sprecher liest \_\_\_\_\_\_\_\_ (die) Meldung.}
    \item \all{Wir hören \_\_\_\_\_\_\_\_ (die) Nachricht \_\_\_\_\_\_\_\_ (mit) \_\_\_\_\_\_\_\_ (der) Aufmerksamkeit.}
    \item \all{\_\_\_\_\_\_\_\_ (dem) Reporter gefällt \_\_\_\_\_\_\_\_ (das) Thema nicht.}
    \item \all{Er antwortet \_\_\_\_\_\_\_\_ (der) Kollegin \_\_\_\_\_\_\_\_ (auf) \_\_\_\_\_\_\_\_ (die) Frage.}
    \item \all{Die Zuhörer warten \_\_\_\_\_\_\_\_ (auf) \_\_\_\_\_\_\_\_ (das) Ergebnis.}
    \item \all{Ein Kind sitzt \_\_\_\_\_\_\_\_ (in) \_\_\_\_\_\_\_\_ (das) Auto \_\_\_\_\_\_\_\_ (und) hört \_\_\_\_\_\_\_\_ (das) Radio.}
    \item \all{Die Mutter steigt \_\_\_\_\_\_\_\_ (in) \_\_\_\_\_\_\_\_ (das) Auto \_\_\_\_\_\_\_\_ (ein).}
    \item \all{Der Sender sendet \_\_\_\_\_\_\_\_ (seit) \_\_\_\_\_\_\_\_ (der) Morgen.}
    \item \all{Das Konzert ist \_\_\_\_\_\_\_\_ (am = an dem) \_\_\_\_\_\_\_\_ (der) Samstag.}
    \item \all{Wir treffen uns \_\_\_\_\_\_\_\_ (um) \_\_\_\_\_\_\_\_ (die) Mittagszeit.}
\end{enumerate}
\end{exercice}

\begin{corrige}[Exercice d'application]
\begin{enumerate}
    \item \all{\textbf{Der} Sprecher liest \textbf{die} Meldung.} (Nom. masc / Akk. fém)
    \item \all{Wir hören \textbf{die} Nachricht \textbf{mit} \textbf{der} Aufmerksamkeit.} (\all{hören} + Akk \rightarrow \all{die} Nachricht ; \all{mit} + Dat \rightarrow \all{der} Aufmerksamkeit).
    \item \all{\textbf{Dem} Reporter gefällt \textbf{das} Thema nicht.} (Dativ pour \all{gefallen} / Nom. neutre).
    \item \all{Er antwortet \textbf{der} Kollegin \textbf{auf} \textbf{die} Frage.} (Dativ \all{antworten} / Akk après \all{auf} $\to$ Wohin? mouvement vers la question).
    \item \all{Die Zuhörer warten \textbf{auf} \textbf{das} Ergebnis.} (\all{warten auf} + Akk).
    \item \all{Ein Kind sitzt \textbf{im} (\textbf{in dem}) Auto und hört \textbf{das} Radio.} (\all{Wo?} $\to$ Dativ \all{in dem} $\to$ \all{im} ; \all{hören} + Akk neutre \all{das}).
    \item \all{Die Mutter steigt \textbf{ins} (\textbf{in das}) Auto \textbf{ein}.} (\all{Wohin?} $\to$ Akk \all{in das} $\to$ \all{ins} ; verbe à particule \all{einsteigen}).
    \item \all{Der Sender sendet \textbf{seit} \textbf{dem} Morgen.} (\all{seit} + Dativ).
    \item \all{Das Konzert ist \textbf{am} \textbf{der} Samstag.} (\all{an} + Dativ jour $\to$ \all{am} ; fém. Dativ \all{der}).
    \item \all{Wir treffen uns \textbf{um} \textbf{die} Mittagszeit.} (\all{um} + Dativ heure $\to$ \all{der} Mittagszeit ; fém. Dativ \all{der}). \emph{Note : \all{Mittagszeit} est féminin.}
\end{enumerate}
\end{corrige}

\begin{attention}[Erreurs fréquentes à éviter (Check-list pour le bac)]
\begin{itemize}
    \item \textbf{Oublier le \emph{-n} du pluriel datif :} \all{mit den Kind\textbf{ern}} (pas \all{*mit den Kinder}).
    \item \textbf{Confondre \emph{die} (fém. Nom/Akk) et \emph{der} (fém. Dat / masc. Nom) :} \all{Ich sehe \textbf{die} Frau} (Akk) vs \all{Ich helfe \textbf{der} Frau} (Dat).
    \item \textbf{Traiter \emph{helfen, danken, gefallen, folgen, gratulieren} comme des verbes à COD :} \all{*Ich helfe ihn} $\to$ \textbf{Faux}. Correct : \all{Ich helfe \textbf{ihm}}.
    \item \textbf{Ne pas contracter les prépositions + article :} \all{in dem} $\to$ \all{im}, \all{in das} $\to$ \all{ins}, \all{zu dem} $\to$ \all{zum}, \all{an dem} $\to$ \all{am}, \all{bei dem} $\to$ \all{beim}. C'est la norme à l'oral et à l'écrit courant.
    \item \textbf{Prépositions de temps :} \all{am} (tag), \all{im} (Monat/Jahr), \all{um} (Uhrzeit) $\to$ \textbf{Toujours Dativ}. \all{für} (Dauer) $\to$ \textbf{Akkusativ}.
\end{itemize}
\end{attention}

\coursec{Grammaire 2 : les prépositions de temps}

\courssub{Transition depuis l'accusatif et le datif}

Dans la section précédente, nous avons vu comment identifier l'\textbf{accusatif} (complément d'objet direct, après certaines prépositions, mesure de temps \emph{sans} préposition) et le \textbf{datif} (complément d'objet indirect, après certaines prépositions, position/lieu \emph{statique}). Les prépositions de temps que nous allons étudier maintenant gouvernent presque toutes le \textbf{datif} (sauf \all{um} + accusatif et \all{bis} + accusatif). Rappelons les contractions fréquentes : \all{am = an + dem}, \all{im = in + dem}, \all{zum = zu + dem}, \all{zur = zu + der}. Maîtriser le datif (articles, pronoms, adjectifs) est donc la clé pour employer correctement ces prépositions.

\courssub{Observation à partir d'un communiqué radio authentique}

Lisons d'abord un court extrait de journal parlé (\all{Hörfunknachrichten}) typique d'une station régionale allemande. Repérez tous les compléments de temps et leurs prépositions.

\begin{textebox}[Extrait de journal parlé — SWR Aktuell, 14. Oktober]
\all{Guten Morgen, hier ist SWR Aktuell. Am 14. Oktober, um 06:30 Uhr, beginnen wir mit den wichtigsten Meldungen. Seit Montag gilt auf der A61 zwischen Koblenz und Ludwigshafen eine neue Geschwindigkeitsbegrenzung. Die Baustelle besteht voraussichtlich bis Freitag. Vor einer Woche gab es dort einen schweren Unfall. Nach der Verkehrsministerkonferenz im Juli wurde die Maßnahme beschlossen. In den nächsten Tagen erwarten wir starken Verkehr. Bleiben Sie dran!}
\end{textebox}

\begin{definition}[Glossaire de l'extrait]
\all{die Meldung, -en} — la nouvelle, le bulletin d'information ; \all{gelten} (gilt, galt, hat gegolten) — s'appliquer, être en vigueur ; \all{die Geschwindigkeitsbegrenzung, -en} — la limitation de vitesse ; \all{die Baustelle, -n} — le chantier ; \all{bestehen} (besteht, bestand, hat bestanden) — exister, durer ; \all{die Verkehrsministerkonferenz, -en} — la conférence des ministres des transports ; \all{die Maßnahme, -n} — la mesure ; \all{beschließen} (beschließt, beschloss, hat beschlossen) — décider ; \all{stark} — fort, important ; \all{der Verkehr} — la circulation ; \all{dranbleiben} — rester à l'écoute.
\end{definition}

\courssub{Les prépositions de temps : forme, cas, emploi}

\begin{propriete}[Règle générale — Cas gouvernés]
Toutes les prépositions de temps de cette leçon (sauf \all{um} et \all{bis}) sont suivies du \textbf{DATIF}. 
\begin{itemize}
\item \all{am} = \all{an} + \all{dem} (datif) → jour, date, partie de la journée.
\item \all{im} = \all{in} + \all{dem} (datif) → mois, année, saison, période future.
\item \all{von ... bis}, \all{bis}, \all{seit}, \all{vor}, \all{nach}, \all{in} (sans article) → datif.
\item \all{um} + heure précise → \textbf{ACCUSATIF}.
\item \all{bis} + jour/date (sans article) → \textbf{ACCUSATIF} (ellipse de \all{bis zum / bis zur}).
\end{itemize}
\end{propriete}

\begin{center}
\begin{tabularx}{\linewidth}{|X|X|X|X|}
\hline
\rowcolor{popblueL}
\textbf{Préposition} & \textbf{Cas} & \textbf{Emploi principal} & \textbf{Exemple + traduction} \\
\hline
\all{am} & Datif & Jour de la semaine, date calendaire, partie de la journée & \all{am Montag} (lundi), \all{am 14. Oktober} (le 14 octobre), \all{am Morgen} (le matin) \\
\hline
\all{um} & Accusatif & Heure précise (horloge) & \all{um 10 Uhr} (à 10 h), \all{um Mitternacht} (à minuit), \all{um 14.30 Uhr} (à 14 h 30) \\
\hline
\all{von ... bis} & Datif / Acc. & Durée délimitée : début — fin & \all{von 8 Uhr bis 20 Uhr} (de 8 h à 20 h), \all{von Montag bis Freitag} (du lundi au vendredi) \\
\hline
\all{bis} & Accusatif & Terminus (jusqu'à) & \all{bis Freitag} (jusqu'à vendredi), \all{bis 20 Uhr} (jusqu'à 20 h) \\
\hline
\all{seit} & Datif & Point de départ dans le passé qui continue jusqu'au présent & \all{seit Montag} (depuis lundi), \all{seit zwei Wochen} (depuis deux semaines) \\
\hline
\all{vor} & Datif & Antériorité : « il y a … » (durée révolue) & \all{vor einer Stunde} (il y a une heure), \all{vor zwei Jahren} (il y a deux ans) \\
\hline
\all{nach} & Datif & Postériorité : « après … » (événement ou moment) & \all{nach dem Essen} (après le repas), \all{nach der Sendung} (après l'émission) \\
\hline
\all{in} & Datif & Futur (dans …) ou période (mois, année, saison) & \all{in einer Stunde} (dans une heure), \all{im Juli} (en juillet), \all{im Jahr 2024} (en 2024) \\
\hline
\end{tabularx}
\end{center}

\courssub{Étude détaillée de chaque préposition}

\textbf{1. \all{am} — Jour, date, partie de la journée (Datif)}\par
\all{am} est la contraction de \all{an dem}. Elle s'emploie devant :
\begin{itemize}
\item les jours de la semaine : \all{am Montag}, \all{am Dienstag}…
\item les dates (chiffre + point + mois) : \all{am 14. Oktober}, \all{am 1. Januar}
\item les parties de la journée (sauf \all{die Nacht} → \all{in der Nacht}) : \all{am Morgen}, \all{am Vormittag}, \all{am Mittag}, \all{am Nachmittag}, \all{am Abend}
\end{itemize}

\begin{exemplebox}[Exemples \all{am}]
\all{Am Samstag fahre ich nach Yaoundé.} — Samedi, je vais à Yaoundé. \\
\all{Die Prüfung ist am 20. Mai.} — L'examen est le 20 mai. \\
\all{Am Morgen höre ich Radio.} — Le matin, j'écoute la radio.
\end{exemplebox}

\textbf{2. \all{um} — Heure précise (Accusatif)}\par
\all{um} n'accepte que l'heure exacte (horloge). L'heure suit la règle de l'accusatif (pas d'article, donc pas de changement visible, mais cas accusatif sous-jacent).

\begin{exemplebox}[Exemples \all{um}]
\all{Der Zug fährt um 10:15 Uhr ab.} — Le train part à 10 h 15. \\
\all{Wir treffen uns um Mitternacht.} — Nous nous retrouvons à minuit. \\
\all{Um wie viel Uhr beginnt die Sendung?} — À quelle heure commence l'émission ?
\end{exemplebox}

\textbf{3. \all{von ... bis} / \all{bis} — Durée et terminus}\par
\begin{itemize}
\item \all{von ... bis} délimite un intervalle : début (\all{von} + datif) — fin (\all{bis} + accusatif/datif selon présence d'article). Avec article : \all{bis zum} (masc./neutre), \all{bis zur} (féminin).
\item \all{bis} seul (sans article) = « jusqu'à » + jour/date/heure : \all{bis Freitag}, \all{bis 20 Uhr}, \all{bis zum 15. Mai}.
\end{itemize}

\begin{exemplebox}[Exemples \all{von ... bis} / \all{bis}]
\all{Das Büro ist von Montag bis Freitag geöffnet.} — Le bureau est ouvert du lundi au vendredi. \\
\all{Die Straße ist von 8 Uhr bis 20 Uhr gesperrt.} — La rue est fermée de 8 h à 20 h. \\
\all{Ich bleibe bis Sonntag.} — Je reste jusqu'à dimanche.
\end{exemplebox}

\textbf{4. \all{seit} — « Depuis » (point de départ continu) — Datif}\par
\cle{Point crucial} : En allemand, \all{seit} + datif s'emploie avec le \textbf{PRÉSENT} (ou le \all{Perfekt} pour action achevée mais effet présent) là où le français utilise « depuis » + présent ou passé composé.
\begin{itemize}
\item \all{seit Montag} (depuis lundi — jour)
\item \all{seit zwei Wochen} (depuis deux semaines — durée)
\item \all{seit dem 1. Januar} (depuis le 1er janvier — date)
\end{itemize}

\begin{exemplebox}[Exemples \all{seit}]
\all{Seit Montag höre ich diesen Sender.} — Depuis lundi, j'écoute cette station. \\
\all{Er wohnt seit zwei Jahren in Berlin.} — Il habite à Berlin depuis deux ans. \\
\all{Seit wann lernen Sie Deutsch?} — Depuis quand apprenez-vous l'allemand ?
\end{exemplebox}

\begin{attention}[Piège n°1 : \all{seit} ≠ \all{für}]
Ne traduisez jamais « depuis » (durée non révolue) par \all{für}. 
\begin{itemize}
\item \all{für zwei Wochen} = « pour deux semaines » (durée prévue, future ou révolue) : \all{Ich fahre für zwei Wochen nach Hamburg.} (Je pars pour deux semaines à Hambourg.)
\item \all{seit zwei Wochen} = « depuis deux semaines » (début dans le passé, continue) : \all{Ich bin seit zwei Wochen in Hamburg.} (Je suis à Hambourg depuis deux semaines.)
\end{itemize}
\end{attention}

\textbf{5. \all{vor} — « Il y a … » (antériorité révolue) — Datif}\par
\all{vor} + datif exprime une durée écoulée depuis un événement passé. Le verbe est au \textbf{PRÄTERITUM} ou \textbf{PERFEKT} (passé).

\begin{exemplebox}[Exemples \all{vor}]
\all{Vor einer Stunde gab es eine Durchsage.} — Il y a une heure, il y a eu une annonce. \\
\all{Vor zwei Jahren habe ich Deutsch angefangen.} — Il y a deux ans, j'ai commencé l'allemand. \\
\all{Wir haben uns vor langer Zeit kennengelernt.} — Nous nous sommes connus il y a longtemps.
\end{exemplebox}

\textbf{6. \all{nach} — « Après … » (postériorité) — Datif}\par
\all{nach} + datif suit un nom (événement, moment, repas) ou un pronom. Ne pas confondre avec l'adverbe \all{danach} (ensuite) ni la conjonction \all{nachdem} (après que + subordonné).

\begin{exemplebox}[Exemples \all{nach}]
\all{Nach der Sendung rufe ich an.} — Après l'émission, j'appelle. \\
\all{Nach dem Essen gehen wir spazieren.} — Après le repas, nous faisons une promenade. \\
\all{Nach Berlin fahre ich erst nächste Woche.} — À Berlin, j'irai seulement la semaine prochaine. (Ici \all{nach} = direction, pas temps — attention au contexte !)
\end{exemplebox}

\textbf{7. \all{in} — « Dans … » (futur) ou « En … » (mois, année, saison) — Datif}\par
Deux emplois distincts :
\begin{itemize}
\item \textbf{Futur / durée avant événement} : \all{in} + durée → \textbf{Datif} : \all{in einer Stunde}, \all{in zwei Tagen}, \all{in einem Jahr}. (Le verbe est au présent ou futur I.)
\item \textbf{Période calendaire} (mois, année, saison) : \all{im} (= \all{in dem}) + datif : \all{im Januar}, \all{im Sommer}, \all{im Jahr 2024}.
\end{itemize}

\begin{exemplebox}[Exemples \all{in}]
\all{Der Bus kommt in zehn Minuten.} — Le bus arrive dans dix minutes. \\
\all{In zwei Wochen fahren wir nach Douala.} — Dans deux semaines, nous allons à Douala. \\
\all{Im Juli ist es in Deutschland oft heiß.} — En juillet, il fait souvent chaud en Allemagne. \\
\all{Im Jahr 2025 feiern wir Jubiläum.} — En 2025, nous fêtons l'anniversaire.
\end{exemplebox}

\courssub{Frise chronologique : \all{vor} — \all{seit} — \all{jetzt} — \all{nach} / \all{in}}

\begin{popfigure}
\begin{tikzpicture}[
  timeline/.style={very thick, popblue},
  tick/.style={thick, popdark},
  labelnode/.style={align=center, font=\small, popdark},
  examplenode/.style={align=center, font=\footnotesize, popmuted},
  arrow/.style={-Stealth, poporange, thick}
]
% Ligne de temps
\draw[timeline] (-6,0) -- (6,0);
% Repères
\foreach \x/\zeit in {-5/vor, -2.5/seit, 0/jetzt, 2.5/nach, 5/in} {
  \draw[tick] (\x,0.2) -- (\x,-0.2);
  \node[labelnode] at (\x,-0.7) {\textbf{\zeit}};
}
% Flèches directionnelles
\draw[arrow] (-5,0.5) -- (-2.5,0.5);
\draw[arrow] (-2.5,0.5) -- (0,0.5);
\draw[arrow] (0,0.5) -- (2.5,0.5);
\draw[arrow] (2.5,0.5) -- (5,0.5);
% Exemples
\node[examplenode, align=center] at (-5,1.2){\all{vor einer Woche}\\(il y a une semaine)};
\node[examplenode, align=center] at (-2.5,1.2){\all{seit Montag}\\(depuis lundi)};
\node[examplenode, popgreen, align=center] at (0,1.2){\textbf{JETZT}\\(maintenant)};
\node[examplenode, align=center] at (2.5,1.2){\all{nach der Sendung}\\(après l'émission)};
\node[examplenode, align=center] at (5,1.2){\all{in einer Stunde}\\(dans une heure)};
% Marqueurs passé/présent/futur
\node[labelnode, poppurple] at (-5,-1.5) {PASSÉ RÉVOLU};
\node[labelnode, popgreen] at (-2.5,-1.5) {PASSÉ → PRÉSENT};
\node[labelnode, popblue] at (0,-1.5) {PRÉSENT};
\node[labelnode, poporange] at (2.5,-1.5) {APRÈS ÉVÉNEMENT};
\node[labelnode, poppink] at (5,-1.5) {FUTUR PROCHE};
\end{tikzpicture}
\legende{Repères temporels : \all{vor} (passé révolu), \all{seit} (début passé continu), \all{nach} (après événement), \all{in} (futur). Le point zéro = \all{jetzt}.}
\end{popfigure}

\courssub{L'ordre des compléments de temps : règle TE-KA-MO-LO}

En allemand, l'ordre canonique des compléments dans la phrase (position centrale, après le verbe conjugué en position 2) suit le principe \textbf{TE-KA-MO-LO} :

\begin{center}
\begin{tabular}{|l|l|l|}
\hline
\textbf{Lettre} & \textbf{Allemand} & \textbf{Français} \\
\hline
T & \textbf{T}emporal (Quand ?) & Temps \\
K & \textbf{K}ausal (Warum ?) & Cause \\
M & \textbf{M}odal (Wie ?) & Manière \\
L & \textbf{L}okal (Wo ? / Wohin ?) & Lieu \\
\hline
\end{tabular}
\end{center}

\emph{Autrement dit : le complément de temps vient \textbf{avant} le complément de cause, qui vient avant le complément de manière, qui vient avant le complément de lieu.}

\begin{exemplebox}[Application TE-KA-MO-LO]
\all{Ich fahre \textbf{morgen} (T) \textbf{wegen des Wetters} (K) \textbf{mit dem Bus} (M) \textbf{nach Yaoundé} (L).} \\
→ Je vais \textbf{demain} (T) \textbf{à cause de la météo} (K) \textbf{en bus} (M) \textbf{à Yaoundé} (L).

\all{Der Sprecher hat \textbf{um 10 Uhr} (T) \textbf{aus aktuellem Anlass} (K) \textbf{sehr deutlich} (M) \textbf{im Studio} (L) \textbf{eine Durchsage gemacht}.} \\
→ Le speaker a \textbf{à 10 h} (T) \textbf{pour une raison d'actualité} (K) \textbf{très clairement} (M) \textbf{au studio} (L) \textbf{fait une annonce}.
\end{exemplebox}

\begin{attention}[Piège n°2 : \all{um} vs \all{am} — Ne pas confondre jour et heure]
\begin{itemize}
\item \all{am Montag} = \textbf{lundi} (jour) — \all{am} + datif
\item \all{um 8 Uhr} = \textbf{à 8 heures} (heure précise) — \all{um} + accusatif
\item \textbf{Combinés} : \all{am Montag um 8 Uhr} = \textbf{lundi à 8 heures} (ordre : jour AVANT heure).
\end{itemize}
\emph{Erreur fréquente :} \all{*um Montag} (faux), \all{*am 8 Uhr} (faux).
\end{attention}

\courssub{Méthode : le trio des dates — \all{am} / \all{um} / \all{in}}

\begin{methode}[Retenir les trois prépositions de base pour situer un événement dans le temps]
\begin{enumerate}
\item \textbf{Jour / Date précise} → \all{am} + datif \\ \all{am Montag}, \all{am 14. Oktober}, \all{am Feiertag}
\item \textbf{Heure précise (horloge)} → \all{um} + accusatif \\ \all{um 10 Uhr}, \all{um 14.30 Uhr}, \all{um Mitternacht}
\item \textbf{Mois / Année / Saison / Période future} → \all{in} + datif (souvent contracté \all{im}) \\ \all{im Juli}, \all{im Jahr 2024}, \all{im Sommer}, \all{in einer Stunde}, \all{in zwei Tagen}
\end{enumerate}
\emph{Astuce mnémotechnique :} \textbf{A–U–I} (comme le début de \all{Am}, \all{Um}, \all{In}) — Jour, Heure, Mois/Année/Futur.
\end{methode}

\courssub{Exercice résolu : compléter un communiqué avec les bonnes prépositions}

\begin{exoresolu}[Compléter le communiqué radio]
\textbf{Énoncé :} Complétez le texte suivant avec les prépositions de temps adaptées (\all{am}, \all{um}, \all{von}, \all{bis}, \all{seit}, \all{vor}, \all{nach}, \all{in} / \all{im}). Attention aux contractions (\all{am}, \all{im}, \all{zum}, \all{zur}) et aux cas !

\all{„Guten Abend, hier ist das Abendprogramm. \_\_\_\_\_\_\_\_\_\_\_\_ 15. Oktober, \_\_\_\_\_\_\_\_\_\_\_\_ 20:00 Uhr, beginnt unsere Sendung „Kulturwelt“. Die Sendung dauert \_\_\_\_\_\_\_\_\_\_\_\_ 20:00 Uhr \_\_\_\_\_\_\_\_\_\_\_\_ 21:30 Uhr. \_\_\_\_\_\_\_\_\_\_\_\_ drei Wochen ist dieses Thema \_\_\_\_\_\_\_\_\_\_\_\_ den Nachrichten. \_\_\_\_\_\_\_\_\_\_\_\_ einem Monat gab es \_\_\_\_\_\_\_\_\_\_\_\_ eine spezielle Ausgabe. \_\_\_\_\_\_\_\_\_\_\_\_ der Sendung können Sie uns anrufen. \_\_\_\_\_\_\_\_\_\_\_\_ November planen wir eine neue Reihe.“}

\tcblower
\textbf{Solution détaillée :}

\begin{enumerate}
\item \all{Am} 15. Oktober — Date → \all{am} + datif (\all{dem} 15. Oktober).
\item \all{um} 20:00 Uhr — Heure précise → \all{um} + accusatif.
\item \all{von} 20:00 Uhr — Début de durée → \all{von} + datif.
\item \all{bis} 21:30 Uhr — Fin de durée (sans article) → \all{bis} + accusatif.
\item \all{Seit} drei Wochen — Point de départ continu (présent : \all{ist}) → \all{seit} + datif.
\item \all{in} den Nachrichten — Période/Lieu « dans les infos » → \all{in} + datif pluriel (\all{den} Nachrichten).
\item \all{Vor} einem Monat — « Il y a un mois » (passé révolu : \all{gab}) → \all{vor} + datif.
\item \all{Nach} der Sendung — « Après l'émission » → \all{nach} + datif (\all{der} Sendung).
\item \all{Im} November — Mois → \all{in} + datif → contraction \all{im} (\all{in dem}).
\end{enumerate}

\textbf{Texte complet corrigé :}
\all{„Guten Abend, hier ist das Abendprogramm. \textbf{Am} 15. Oktober, \textbf{um} 20:00 Uhr, beginnt unsere Sendung „Kulturwelt“. Die Sendung dauert \textbf{von} 20:00 Uhr \textbf{bis} 21:30 Uhr. \textbf{Seit} drei Wochen ist dieses Thema \textbf{in} den Nachrichten. \textbf{Vor} einem Monat gab es \textbf{eine} spezielle Ausgabe. \textbf{Nach} der Sendung können Sie uns anrufen. \textbf{Im} November planen wir eine neue Reihe.“}
\end{exoresolu}

\courssub{Exercices d'entraînement}

\begin{exercice}[Entraînement 1 — Choisir la bonne préposition]
Complétez par \all{am}, \all{um}, \all{von}, \all{bis}, \all{seit}, \all{vor}, \all{nach}, \all{in} / \all{im} (avec contraction si nécessaire).
\begin{enumerate}
\item \_\_\_\_\_\_\_\_\_\_\_\_ Montag fahre ich nach Douala.
\item Der Zug kommt \_\_\_\_\_\_\_\_\_\_\_\_ 14:30 Uhr an.
\item Ich lerne Deutsch \_\_\_\_\_\_\_\_\_\_\_\_ zwei Jahren.
\item \_\_\_\_\_\_\_\_\_\_\_\_ einer Stunde beginnt der Film.
\item Die Ausstellung ist \_\_\_\_\_\_\_\_\_\_\_\_ 10 Uhr \_\_\_\_\_\_\_\_\_\_\_\_ 18 Uhr geöffnet.
\item \_\_\_\_\_\_\_\_\_\_\_\_ dem Essen gehen wir ins Kino.
\item \_\_\_\_\_\_\_\_\_\_\_\_ Juli ist es heiß.
\item Wir warten \_\_\_\_\_\_\_\_\_\_\_\_ Freitag auf die Antwort.
\end{enumerate}
\end{exercice}

\begin{exercice}[Entraînement 2 — Traduction thème (français → allemand)]
Traduisez en allemand en respectant l'ordre TE-KA-MO-LO.
\begin{enumerate}
\item Depuis lundi, j'écoute la radio le matin.
\item Il y a une heure, il y a eu une annonce importante.
\item Lundi à 10 heures, le speaker fait une annonce au studio.
\item Dans deux semaines, nous allons au Cameroun.
\item Après l'émission, je rentre à la maison en taxi.
\end{enumerate}
\end{exercice}

\begin{corrige}[Exercice 1]
1. \textbf{Am} Montag — 2. \textbf{um} 14:30 Uhr — 3. \textbf{seit} zwei Jahren — 4. \textbf{In} einer Stunde — 5. \textbf{von} 10 Uhr \textbf{bis} 18 Uhr — 6. \textbf{Nach} dem Essen — 7. \textbf{Im} Juli — 8. \textbf{bis} Freitag
\end{corrige}

\begin{corrige}[Exercice 2]
1. \textbf{Seit Montag höre ich am Morgen Radio.} (T: \all{seit Montag} / \all{am Morgen} — pas de K, M, L)
2. \textbf{Vor einer Stunde gab es eine wichtige Durchsage.} (T: \all{vor einer Stunde})
3. \textbf{Am Montag um 10 Uhr macht der Sprecher im Studio eine Durchsage.} (T: \all{am Montag um 10 Uhr} — L: \all{im Studio})
4. \textbf{In zwei Wochen fahren wir nach Kamerun.} (T: \all{in zwei Wochen} — L: \all{nach Kamerun})
5. \textbf{Nach der Sendung fahre ich mit dem Taxi nach Hause.} (T: \all{nach der Sendung} — M: \all{mit dem Taxi} — L: \all{nach Hause})
\end{corrige}

\courssub{Synthèse visuelle : tableau récapitulatif complet}

\begin{center}
\begin{tabularx}{\linewidth}{|X|X|X|X|}
\hline
\rowcolor{popblueL}
\textbf{Préposition} & \textbf{Cas} & \textbf{Question} & \textbf{Exemples clés} \\
\hline
\all{am} (= \all{an dem}) & Datif & \all{Wann?} (Jour/Date) & \all{am Montag}, \all{am 14.10.}, \all{am Abend} \\
\hline
\all{um} & Accusatif & \all{Wann?} (Heure exacte) & \all{um 10 Uhr}, \all{um Mitternacht} \\
\hline
\all{von ... bis} & Datif / Acc. & \all{Wie lange?} (Durée) & \all{von 8 bis 20 Uhr}, \all{von Mo. bis Fr.} \\
\hline
\all{bis} & Accusatif & \all{Bis wann?} (Terminus) & \all{bis Freitag}, \all{bis 20 Uhr}, \all{bis zum 15.5.} \\
\hline
\all{seit} & Datif & \all{Seit wann?} (Depuis → présent) & \all{seit Montag}, \all{seit 2 Wochen}, \all{seit dem 1.1.} \\
\hline
\all{vor} & Datif & \all{Wann?} (Il y a … → passé) & \all{vor 1 Stunde}, \all{vor 2 Jahren}, \all{vor langer Zeit} \\
\hline
\all{nach} & Datif & \all{Nach wem/was?} (Après) & \all{nach der Sendung}, \all{nach dem Essen} \\
\hline
\all{in} / \all{im} & Datif & \all{Wann?} (Dans / En mois/année) & \all{in 1 Stunde}, \all{im Juli}, \all{im Jahr 2024} \\
\hline
\end{tabularx}
\end{center}

\begin{savaistu}[Le saviez-vous ? — L'heure officielle en Allemagne]
En Allemagne (comme en France), on utilise l'horaire de 24 h à l'écrit et dans les contextes formels (horaires de train, communiqués radio, administration) : \all{14:30 Uhr}, \all{20:00 Uhr}. À l'oral familier, on utilise souvent le système de 12 h avec \all{vormittags} / \all{nachmittags} / \all{abends} / \all{nachts} : \all{„Es ist halb drei“} (13:30 ou 0:30 selon le contexte), \all{„viertel nach zehn“} (10:15). À la radio, l'heure est toujours annoncée en 24 h : \all{„Es ist 14 Uhr 30“}.
\end{savaistu}

\begin{aretenir}[À retenir pour l'évaluation]
\begin{itemize}
\item \textbf{7 prépositions de temps} : \all{am}, \all{um}, \all{von...bis}, \all{bis}, \all{seit}, \all{vor}, \all{nach}, \all{in/im}.
\item \textbf{Cas} : presque toutes \textbf{DATIF} (sauf \all{um} + ACC, \all{bis} + ACC).
\item \textbf{Contractions} : \all{am} = \all{an dem}, \all{im} = \all{in dem}.
\item \textbf{Seit + présent} : action commencée dans le passé, continue maintenant.
\item \textbf{Vor + passé} : action révolue (« il y a … »).
\item \textbf{Ordre TE-KA-MO-LO} : Temps avant Cause avant Manière avant Lieu.
\item \textbf{Trio date} : \all{am} (jour) — \all{um} (heure) — \all{in/im} (mois/année/futur).
\end{itemize}
\end{aretenir}

\coursec{Méthode du commentaire et production guidée}

Dans la section précédente, nous avons étudié les prépositions de temps (\all{am}, \all{um}, \all{von ... bis}, \all{seit}, \all{ab}, \all{bis}) qui permettent de situer un événement dans le temps. Ces prépositions sont précisément l'outil indispensable du communiqué radio : une annonce sans date ni heure ne sert à rien ! Nous allons maintenant mettre en pratique tout ce que nous avons appris dans cette leçon — le vocabulaire des médias, l'accusatif et le datif, les prépositions de temps — pour atteindre deux objectifs : \cle{commenter} un texte de manière méthodique et \cle{produire} nous-mêmes un communiqué radio, à l'écrit puis à l'oral.

\courssub{La méthode du commentaire de texte en quatre étapes}

Au lycée, on vous demandera souvent de « commenter » un texte allemand. Beaucoup d'élèves improvisent et perdent des points. Or le commentaire est un exercice très codifié : il suffit de suivre quatre étapes, toujours dans le même ordre.

\begin{methode}[Commenter un texte en 4 étapes]
\begin{enumerate}
\item \textbf{Présenter} le texte : de quel \cle{type} de texte s'agit-il (article, lettre, communiqué radio, dialogue) ? Quelle est sa \cle{source} (un journal, une radio, une école) ? Quel est son \cle{thème} général ? Une ou deux phrases suffisent.
\item \textbf{Résumer} le contenu avec vos propres mots : qui fait quoi, quand, où, et avec quelles conséquences ? Le résumé ne doit pas dépasser un tiers du texte original.
\item \textbf{Analyser} la langue : relevez le \cle{lexique} du thème (ici : \all{die Durchsage}, \all{der Sprecher}, \all{die Meldung}...), les \cle{structures} grammaticales remarquables (cas, prépositions, place du verbe) et le \cle{ton} du texte (informatif, urgent, festif, officiel).
\item \textbf{Juger} le texte : quel est son intérêt ? Est-il d'actualité ? Faites une comparaison avec la situation au Cameroun (par exemple : les annonces à la radio CRTV, les communiqués lus au porte-voix dans les quartiers de Yaoundé ou de Douala).
\end{enumerate}
\end{methode}

\begin{popfigure}
\begin{center}
\begin{tikzpicture}[
  node distance=0.9cm,
  etape/.style={rectangle, rounded corners=4pt, draw=popblue, fill=popblueL, thick, text width=3.4cm, align=center, minimum height=1.1cm, font=\small},
  fleche/.style={-{Stealth[length=3mm]}, thick, poporange}
]
\node[etape, align=center] (p){\textbf{1. Présentation}\\ type, source, thème};
\node[etape, right=of p, align=center] (r){\textbf{2. Résumé}\\ qui ? quoi ? quand ? où ?};
\node[etape, right=of r, align=center] (a){\textbf{3. Analyse}\\ lexique, structures, ton};
\node[etape, right=of a, align=center] (j){\textbf{4. Jugement}\\ intérêt, actualité, Cameroun};
\draw[fleche] (p) -- (r);
\draw[fleche] (r) -- (a);
\draw[fleche] (a) -- (j);
\end{tikzpicture}
\end{center}
\legende{Organigramme du commentaire de texte : quatre étapes dans un ordre fixe.}
\end{popfigure}

\begin{attention}[L'erreur classique au commentaire]
Ne commencez jamais votre commentaire par l'analyse grammaticale ! Le correcteur attend d'abord la preuve que vous avez \emph{compris} le texte (étapes 1 et 2). L'analyse de la langue ne vient qu'ensuite. Autre piège : recopier des phrases entières du texte au lieu de résumer avec vos propres mots — cela montre que vous n'avez pas compris.
\end{attention}

\courssub{Commentaire modèle rédigé}

Observons d'abord un commentaire complet, rédigé en allemand simple (niveau A2), sur le texte « Eine wichtige Durchsage » étudié dans cette leçon. Remarquez comment chaque étape est visible.

\begin{exemplebox}[Commentaire modèle : „Eine wichtige Durchsage“]
\all{Der Text ist ein Radiokommuniqué. Ein Sprecher informiert die Zuhörer über eine wichtige Durchsage für die Schule. Das Thema ist eine Veranstaltung: ein Sportfest.}

\all{Der Sprecher sagt, wann und wo das Fest stattfindet: am Freitag um 9 Uhr im Schulhof. Er sagt auch, dass es am Nachmittag keinen Unterricht gibt.}

\all{Der Text benutzt das Vokabular der Medien: „die Durchsage“, „die Zuhörer“, „informieren“. Man findet Präpositionen der Zeit: „am Freitag“, „um 9 Uhr“. Der Ton ist freundlich und klar.}

\all{Das Kommuniqué ist nützlich und aktuell. In Kamerun hören viele Schüler auch Radio, zum Beispiel CRTV. Solche Durchsagen sind überall wichtig.}

\emph{Traduction :} Le texte est un communiqué radio. Un speaker informe les auditeurs d'une annonce importante pour l'école. Le thème est une manifestation : une fête sportive. / Le speaker dit quand et où la fête a lieu : vendredi à 9 heures dans la cour de l'école. Il dit aussi qu'il n'y a pas cours l'après-midi. / Le texte utilise le vocabulaire des médias : « l'annonce », « les auditeurs », « informer ». On trouve des prépositions de temps : « vendredi », « à 9 heures ». Le ton est amical et clair. / Ce communiqué est utile et d'actualité. Au Cameroun, beaucoup d'élèves écoutent aussi la radio, par exemple la CRTV. De telles annonces sont importantes partout.
\end{exemplebox}

\begin{aretenir}[Les phrases toutes faites du commentaire]
Mémorisez ces amorces, réutilisables pour n'importe quel texte :
\begin{itemize}
\item \all{Der Text ist ein Kommuniqué / ein Artikel / ein Dialog.} « Le texte est un communiqué / un article / un dialogue. »
\item \all{Das Thema ist ...} « Le thème est ... »
\item \all{Der Sprecher informiert die Zuhörer über ...} (+ accusatif) « Le speaker informe les auditeurs de ... »
\item \all{Der Ton ist informativ / freundlich / dringend.} « Le ton est informatif / amical / urgent. »
\item \all{In Kamerun ist die Situation ähnlich / anders.} « Au Cameroun, la situation est semblable / différente. »
\end{itemize}
\end{aretenir}

\courssub{Le plan type du résumé}

Le résumé est l'étape 2 du commentaire, mais c'est aussi un exercice autonome fréquent à l'examen. Sa structure est très simple :

\begin{propriete}[Structure d'un résumé efficace]
Un bon résumé = \textbf{une phrase de thème} + \textbf{les informations essentielles} :
\begin{enumerate}
\item \textbf{Phrase de thème} : \all{In diesem Kommuniqué geht es um ein Sportfest an der Schule.} « Dans ce communiqué, il s'agit d'une fête sportive à l'école. »
\item \textbf{Quoi ?} (\all{was}) : l'événement annoncé.
\item \textbf{Quand ?} (\all{wann}) : la date et l'heure — ici, vos prépositions de temps !
\item \textbf{Où ?} (\all{wo}) : le lieu.
\item \textbf{Conséquences} : ce que les auditeurs doivent faire ou savoir (\all{es gibt keinen Unterricht}).
\end{enumerate}
Règle d'or : pas de détails secondaires, pas de citations, pas d'opinion personnelle dans le résumé (le jugement vient à l'étape 4).
\end{propriete}

\courssub{Rappel : les prépositions de temps principales}

Avant de passer à la production, rappelons les formes les plus fréquentes pour situer l'événement dans le communiqué.

\begin{aretenir}[Prépositions de temps — aperçu]
\begin{center}
\begin{tabularx}{\linewidth}{|X|X|X|}
\hline
\rowcolor{popblueL} \textbf{Préposition} & \textbf{Cas} & \textbf{Exemple (allemand / français)} \\
\hline
\all{am} (= \all{an dem}) & Datif & \all{am Freitag} / \all{am 20. März} — « le vendredi / le 20 mars » \\
\hline
\all{um} & Accusatif & \all{um 9 Uhr} / \all{um 14.30 Uhr} — « à 9 heures / à 14 h 30 » \\
\hline
\all{von ... bis} & Datif & \all{von 8 bis 12 Uhr} — « de 8 à 12 heures » \\
\hline
\all{seit} & Datif & \all{seit Montag} / \all{seit 2020} — « depuis lundi / depuis 2020 » \\
\hline
\all{ab} & Datif & \all{ab nächster Woche} — « à partir de la semaine prochaine » \\
\hline
\all{bis} & Accusatif & \all{bis Freitag} / \all{bis 18 Uhr} — « jusqu'à vendredi / jusqu'à 18 h » \\
\hline
\all{in} (durée future) & Datif & \all{in einer Woche} — « dans une semaine » \\
\hline
\end{tabularx}
\end{center}
\emph{Attention :} \all{wegen} (à cause de) n'est \textbf{pas} une préposition de temps, mais une préposition causale qui régit le \textbf{génitif} (\all{wegen des Festes}).
\end{aretenir}

\courssub{Production guidée : rédiger un communiqué radio}

Passons maintenant à l'expression écrite. Votre tâche type : rédiger un communiqué radio de 5 à 6 phrases annonçant un événement scolaire (fête de l'école, match de football, examen reporté).

\begin{methode}[Un bon communiqué radio en 30 secondes]
Un communiqué radio professionnel répond en 30 secondes à cinq questions :
\begin{enumerate}
\item \all{Wer?} Qui parle, qui organise ? (\all{die Schulleitung}, \all{der Sportverein})
\item \all{Was?} Quel est l'événement ? (\all{ein Sportfest}, \all{ein Fußballspiel})
\item \all{Wann?} Quand ? — préposition de temps : \all{am Freitag}, \all{um 9 Uhr}
\item \all{Wo?} Où ? (\all{im Schulhof}, \all{in der Aula})
\item \all{Was müssen die Zuhörer tun?} Que doivent faire les auditeurs ? (\all{kommt zahlreich!}, \all{bringt eure Hefte mit!})
\end{enumerate}
Si votre texte répond aux cinq questions, il est complet.
\end{methode}

\begin{popfigure}
\begin{center}
\begin{tikzpicture}[
  node distance=0.55cm,
  bloc/.style={rectangle, rounded corners=4pt, draw=popgreen, fill=popgreenL, thick, text width=7.6cm, align=center, minimum height=0.85cm, font=\small},
  fleche/.style={-{Stealth[length=3mm]}, thick, popdark}
]
\node[bloc] (s) {\textbf{Salutation} : „Liebe Zuhörer!“};
\node[bloc, below=of s] (i) {\textbf{Information} : l'événement (wer ? was ?)};
\node[bloc, below=of i] (d) {\textbf{Détails pratiques} : date, heure, lieu (wann ? wo ?)};
\node[bloc, below=of d] (c) {\textbf{Consignes} : ce que les auditeurs doivent faire};
\node[bloc, below=of c] (f) {\textbf{Clôture} : „Bleiben Sie dran!“};
\draw[fleche] (s) -- (i);
\draw[fleche] (i) -- (d);
\draw[fleche] (d) -- (c);
\draw[fleche] (c) -- (f);
\end{tikzpicture}
\end{center}
\legende{La structure en cinq blocs d'un communiqué radio scolaire.}
\end{popfigure}

\begin{exemplebox}[Les formules types du speaker]
\begin{center}
\begin{tabularx}{\linewidth}{|X|X|}
\hline
\rowcolor{popblueL} \textbf{Deutsch} & \textbf{Français} \\
\hline
\all{Liebe Zuhörer!} & Chers auditeurs ! \\
\hline
\all{Wir informieren Sie über ...} (+ Akk.) & Nous vous informons sur ... \\
\hline
\all{Am Freitag findet ein Fest statt.} & Vendredi, une fête aura lieu. \\
\hline
\all{Es beginnt um 9 Uhr im Schulhof.} & Elle commence à 9 heures dans la cour. \\
\hline
\all{Wegen des Festes gibt es keinen Unterricht.} & À cause de la fête, il n'y a pas cours. \\
\hline
\all{Kommt zahlreich!} & Venez nombreux ! \\
\hline
\all{Bleiben Sie auf unserer Frequenz!} & Restez à l'écoute ! / Ne changez pas de station ! \\
\hline
\all{Bleiben Sie dran!} & Restez avec nous ! \\
\hline
\end{tabularx}
\end{center}
\end{exemplebox}

\begin{textebox}[Modèle de production : un communiqué radio scolaire]
\all{Liebe Zuhörer! Am Freitag, dem 20. März, findet an unserer Schule ein Sportfest statt. Es beginnt um 9 Uhr im Schulhof. Alle Schüler und ihre Eltern sind herzlich eingeladen. Wegen des Festes gibt es am Nachmittag keinen Unterricht. Wir informieren Sie morgen über die Ergebnisse. Bleiben Sie dran!}

\emph{Traduction :} Chers auditeurs ! Vendredi 20 mars, une fête sportive aura lieu dans notre école. Elle commence à 9 heures dans la cour de l'école. Tous les élèves et leurs parents sont chaleureusement invités. À cause de la fête, il n'y aura pas cours l'après-midi. Nous vous informerons demain des résultats. Restez avec nous !

\emph{Glossaire :} \all{das Sportfest, -e} = la fête sportive ; \all{der Schulhof, -höfe} = la cour de l'école ; \all{eingeladen} = invité ; \all{das Ergebnis, -se} = le résultat ; \all{wegen} (+ génitif) = à cause de.
\end{textebox}

Analysons ce modèle : il suit exactement les cinq blocs du schéma. Il contient deux prépositions de temps (\all{am Freitag}, \all{um 9 Uhr}), un accusatif correct (\all{über die Ergebnisse}), un datif correct (\all{an unserer Schule} — datif car il s'agit d'un lieu, question \all{wo ?}), et le verbe est toujours en deuxième position, même après \all{am Freitag} qui ouvre la phrase : \all{Am Freitag \textbf{findet} ... statt.}

\begin{exoresolu}[Rédiger un communiqué : examen reporté]
Énoncé : Vous êtes le speaker de la radio de votre lycée à Yaoundé. Rédigez un communiqué de 5-6 phrases pour annoncer que l'examen d'allemand est reporté de lundi à jeudi, à 8 heures, dans la salle 12. Consignes : saluez, informez, donnez les détails, dites aux élèves quoi faire, concluez.
\tcblower
Solution détaillée :

\all{Liebe Zuhörer! Wir informieren Sie über eine wichtige Meldung: Die Deutschprüfung am Montag findet nicht statt. Sie findet jetzt am Donnerstag um 8 Uhr im Saal 12 statt. Alle Schüler müssen ihre Hefte und Stifte mitbringen. Bitte kommt um 7.45 Uhr zur Schule! Bleiben Sie auf unserer Frequenz!}

« Chers auditeurs ! Nous vous informons d'une nouvelle importante : l'examen d'allemand de lundi n'aura pas lieu. Il aura lieu jeudi à 8 heures en salle 12. Tous les élèves doivent apporter leurs cahiers et leurs stylos. Venez à l'école à 7 h 45, s'il vous plaît ! Restez à l'écoute ! »

Vérification : salutation (phrase 1), information (phrase 2), détails avec deux prépositions de temps \all{am Donnerstag} et \all{um 8 Uhr} (phrase 3), consignes (phrases 4-5), clôture (phrase 6). Cas corrects : \all{über eine wichtige Meldung} (accusatif après \all{über}), \all{im Saal 12} (= \all{in dem Saal 12}, datif lieu), \all{zur Schule} (= \all{zu der}, datif). Verbe en 2e position partout.
\end{exoresolu}

\begin{attention}[Avant de rendre votre copie : la double vérification]
Deux contrôles obligatoires dans toute production écrite :
\begin{enumerate}
\item \textbf{La place du verbe} : en allemand, le verbe conjugué est toujours en \cle{deuxième position}. Si vous commencez par une date, le verbe suit immédiatement : \all{Am Freitag \textbf{findet} ein Fest statt.} — et jamais \emph{Am Freitag ein Fest findet statt} (calque du français « Vendredi, une fête aura lieu »).
\item \textbf{Le cas après chaque préposition de temps} : \all{am} = \all{an} + \all{dem} (datif), \all{um} + accusatif, \all{seit} + datif, \all{von ... bis} + datif, \all{ab} + datif, \all{bis} + accusatif. Relisez chaque préposition et demandez-vous : « quel cas ai-je mis, et pourquoi ? »
\end{enumerate}
Ces deux erreurs sont les plus fréquentes chez les francophones et les plus sanctionnées à l'examen.
\end{attention}

\courssub{Grille d'auto-évaluation et entraînement oral}

Avant de rendre votre communiqué, évaluez-vous vous-même avec cette grille :

\begin{center}
\begin{tabularx}{\linewidth}{|X|l|}
\hline
\rowcolor{popblueL} \textbf{Critère} & \textbf{Points} \\
\hline
Structure d'annonce complète (salutation, information, détails, consignes, clôture) & / 5 \\
\hline
Vocabulaire du thème (\all{die Durchsage}, \all{die Zuhörer}, \all{die Meldung}...) & / 3 \\
\hline
Au moins 2 prépositions de temps correctes (\all{am}, \all{um}, \all{von ... bis}) & / 4 \\
\hline
Au moins 1 accusatif correct (après \all{über}, \all{für}...) & / 2 \\
\hline
Au moins 1 datif correct (après \all{an}, \all{in}, \all{zu}...) & / 2 \\
\hline
Place du verbe en 2e position dans chaque phrase & / 4 \\
\hline
\textbf{Total} & \textbf{/ 20} \\
\hline
\end{tabularx}
\end{center}

\begin{experience}[Lire son communiqué « comme un speaker »]
Par groupes de quatre, chaque élève lit son communiqué à voix haute, comme un vrai speaker de radio. Les autres ferment les yeux et vérifient avec la grille d'auto-évaluation. Conseils de speaker :
\begin{enumerate}
\item \textbf{L'intonation} : la voix monte légèrement sur l'information importante, puis redescend en fin de phrase.
\item \textbf{Les pauses} : marquez un temps d'arrêt après la salutation \all{Liebe Zuhörer!} et avant la consigne finale.
\item \textbf{La clarté} : articulez chaque mot ; prononcez le \all{w} allemand comme un « v » français (\all{Wir} = « vir ») et le \all{v} comme un « f » (\all{Veranstaltung} = « fér-an-chtal-oung »).
\end{enumerate}
Ensuite, le groupe élit le meilleur speaker, qui lit son texte devant la classe entière.
\end{experience}

\begin{exercice}[À vous de jouer]
Rédigez un communiqué radio de 5 à 6 phrases pour annoncer la fête de fin d'année de votre lycée : samedi 15 juin, 10 heures, dans la cour, avec de la musique et un match de football ; les parents sont invités. Utilisez la grille d'auto-évaluation avant de rendre votre copie, puis lisez votre texte à voix haute chez vous en vous enregistrant.
\end{exercice}


\begin{corrige}[Éléments de correction]
Un communiqué possible : \all{Liebe Zuhörer! Am Samstag, dem 15. Juni, findet an unserem Gymnasium die große Schulfeier statt. Sie beginnt um 10 Uhr im Schulhof. Es gibt Musik und ein Fußballspiel. Alle Eltern sind herzlich eingeladen. Kommt zahlreich und bleibt dran!}

Points à vérifier : verbe en 2e position (\all{findet}, \all{beginnt}, \all{gibt}, \all{sind}) ; prépositions de temps \all{am Samstag} et \all{um 10 Uhr} ; datif \all{an unserem Gymnasium}, \all{im Schulhof} ; accusatif éventuel \all{über die Feier} si la formule \all{wir informieren Sie über ...} est utilisée ; salutation et clôture présentes.
\end{corrige}

\begin{savaistu}[Le communiqué au baccalauréat]
Au baccalauréat camerounais, l'expression écrite d'allemand est souvent notée selon trois critères : le \cle{respect de la consigne} (avez-vous fait ce qui était demandé ?), la \cle{correction grammaticale} (cas, place du verbe, conjugaison) et la \cle{richesse du vocabulaire}. Un communiqué radio bien structuré — salutation, information, détails, consignes, clôture — coche automatiquement les trois cases : il respecte la consigne, il impose des phrases simples donc correctes, et il utilise naturellement tout le vocabulaire des médias appris dans cette leçon. C'est l'un des sujets les plus « rentables » de l'examen !
\end{savaistu}

\begin{aretenir}[L'essentiel de la section]
\begin{itemize}
\item Le commentaire de texte suit quatre étapes fixes : présentation, résumé, analyse, jugement.
\item Le résumé = une phrase de thème + quoi, quand, où, conséquences.
\item Le communiqué radio suit cinq blocs : salutation, information, détails pratiques, consignes, clôture.
\item Il répond en 30 secondes à cinq questions : \all{wer? was? wann? wo? was müssen die Zuhörer tun?}
\item Avant de rendre la copie : vérifier la place du verbe (2e position) et le cas après chaque préposition.
\end{itemize}
\end{aretenir}

\newpage

\coursec{Activité d'intégration}

\courssub{Objectifs de l'activité}
\begin{itemize}
    \item Produire un communiqué radio authentique de 6 à 8 phrases en mobilisant le lexique des médias (\cle{die Durchsage}, \cle{der Sprecher}, \cle{die Meldung}, \cle{die Veranstaltung}, \cle{die Mitteilung}).
    \item Appliquer correctement les cas \textbf{accusatif} (objet direct, prépositions \emph{durch, für, gegen, ohne, um}) et \textbf{datif} (objet indirect, verbes \emph{helfen, danken, gefallen}, prépositions \emph{aus, bei, mit, nach, seit, von, zu}).
    \item Employer au moins trois prépositions de temps différentes (\cle{am}, \cle{um}, \cle{von … bis}, \cle{seit}, \cle{bis}, \cle{ab}) pour situer l'événement.
    \item Respecter la structure type d'un bulletin radio : formule d'ouverture (\all{Liebe Zuhörer!}), corps du message, formule de clôture (\all{Wir wünschen Ihnen einen schönen Tag!}).
    \item Développer des compétences orales : lecture expressive (intonation de speaker), écoute active (fiche \emph{Wer? Was? Wann? Wo?}), interaction (question de compréhension).
\end{itemize}

\courssub{Situation}
\begin{textebox}[Contexte : Le club radio du Lycée Bilingue de Yaoundé]
Le \cle{Radioclub} du Lycée Bilingue de Yaoundé prépare une émission spéciale \all{„Schulnachrichten“} (Nouvelles de l'école) qui sera diffusée vendredi matin sur la sono de l'établissement. Vous êtes membres de l'équipe de rédaction. Votre mission : rédiger et présenter un communiqué officiel (\cle{die Durchsage}) annonçant un événement important de la vie scolaire. L'événement peut être réel (tournoi de football inter-classes, journée culturelle \emph{Kulturtag}, visite d'une délégation du Goethe-Institut Yaoundé, fermeture de la route devant le lycée pour travaux) ou imaginaire mais plausible (inauguration d'un nouveau laboratoire de langues, concert de l'orchestre scolaire). Le communiqué sera lu en direct par un \cle{Sprecher} (speaker) désigné par le groupe. Les autres élèves jouent le rôle d'auditeurs attentifs : ils remplissent une fiche d'écoute structurée et posent ensuite une question pour vérifier leur compréhension.
\end{textebox}

\courssub{Consignes}
\begin{enumerate}
    \item \textbf{Formation des groupes (3–4 élèves).} Chaque groupe choisit un thème parmi la liste ci-dessous ou propose son propre événement (validation par l'enseignant) :
    \begin{itemize}
        \item Tournoi de football inter-classes (finale le mercredi 15 mai à 14 h sur le terrain synthétique).
        \item Journée culturelle \emph{Kulturtag} (vendredi 24 mai, 9 h–16 h, salle polyvalente : chants, danses, théâtre, exposition d'art).
        \item Visite d'une délégation du Goethe-Institut Yaoundé (mardi 11 juin, 10 h, amphithéâtre : présentation des programmes d'échange et bourses PASCH).
        \item Fermeture de la route principale devant le lycée (du lundi 3 juin au vendredi 7 juin, déviation par l'avenue Kennedy).
        \item Inauguration du nouveau laboratoire de langues (jeudi 20 juin, 11 h, bâtiment C, en présence du Proviseur et de l'Ambassadeur d'Allemagne).
    \end{itemize}
    \item \textbf{Rédaction collaborative (20 min).} Rédigez le communiqué en allemand (6–8 phrases). Contraintes obligatoires :
    \begin{itemize}
        \item Formule d'ouverture : \all{Liebe Zuhörer!} ou \all{Liebe Mitschülerinnen und Mitschüler!}
        \item Au moins \textbf{4 mots} du lexique de la leçon (ex. \cle{die Durchsage}, \cle{der Sprecher}, \cle{die Meldung}, \cle{die Veranstaltung}, \cle{die Mitteilung}, \cle{der Hörer}, \cle{der Termin}, \cle{der Ort}).
        \item Au moins \textbf{3 prépositions de temps différentes} (ex. \all{am Mittwoch}, \all{um 14 Uhr}, \all{von Montag bis Freitag}, \all{seit 9 Uhr}, \all{bis 16 Uhr}, \all{ab 10 Uhr}).
        \item Au moins \textbf{un verbe régissant le datif} (\cle{helfen}, \cle{danken}, \cle{gefallen}, \cle{gehören}, \cle{folgen}, \cle{glauben}) utilisé correctement avec son objet au datif.
        \item Formule de clôture type : \all{Wir danken für Ihre Aufmerksamkeit.} / \all{Wir wünschen Ihnen einen schönen Tag.} / \all{Bleiben Sie dran!}
        \item Respect de l'ordre des mots (verbe en position 2, compléments : temps – manière – lieu) et des majuscules aux noms.
    \end{itemize}
    \item \textbf{Répétition orale (5 min).} Le \cle{Sprecher} désigné s'entraîne à lire le texte avec une diction claire, un débit modéré, une intonation montante à l'ouverture et descendante à la clôture, en marquant les pauses (virgules, points). Les autres membres du groupe l'aident à corriger la prononciation des mots difficiles (ex. \all{Veranstaltung} [fɛɐ̯ˈʔanʃtaltʊŋ], \all{Durchsage} [ˈdʊʁçˌzaːɡə], \all{Mittwoch} [ˈmɪtvox]).
    \item \textbf{Présentation devant la classe (3 min par groupe).} Le \cle{Sprecher} lit le communiqué debout, face à la classe, sans lire mot à mot (regard vers l'auditoire). Les autres groupes remplissent la fiche d'écoute :
    \begin{center}
    \begin{tabular}{|l|l|}
    \hline
    \textbf{Wer?} (Qui parle / Quelle institution ?) & \ldots \\
    \hline
    \textbf{Was?} (Quel événement ?) & \ldots \\
    \hline
    \textbf{Wann?} (Date, heure, durée) & \ldots \\
    \hline
    \textbf{Wo?} (Lieu précis) & \ldots \\
    \hline
    \end{tabular}
    \end{center}
    \item \textbf{Questions de compréhension (2 min par groupe).} Chaque groupe auditeur pose \textbf{une question} en allemand au groupe présentateur (ex. \all{Wann findet das Turnier statt?} – \all{Wo genau ist die Veranstaltung?} – \all{Wem danken die Organisatoren?}). Le groupe présentateur répond brièvement en allemand.
    \item \textbf{Vote et bilan (5 min).} La classe vote à main levée pour le meilleur bulletin selon trois critères : \emph{contenu complet}, \emph{correction grammaticale (cas, prépositions)}, \emph{qualité de l'oral (intonation, fluidité)}. L'enseignant commente les points forts et les erreurs récurrentes.
\end{enumerate}

\courssub{Ressources à mobiliser}
\begin{propriete}[Rappel : Accusatif et Datif – formes et emplois principaux]
\begin{center}
\begin{tabularx}{\linewidth}{|X|X|X|}
\hline
\rowcolor{popblueL} \textbf{Cas} & \textbf{Article défini (m/f/n/pl)} & \textbf{Emplois clés (niveau A1+/A2)} \\
\hline
\textbf{Nominatif} & der / die / das / die & Sujet : \all{Der Sprecher liest die Meldung.} \\
\hline
\textbf{Accusatif} & \textbf{den} / die / das / die & Objet direct : \all{Ich höre die Durchsage.} \\
 & & Prépositions \emph{durch, für, gegen, ohne, um} : \all{für die Veranstaltung} \\
\hline
\textbf{Datif} & \textbf{dem} / \textbf{der} / \textbf{dem} / \textbf{den (+ n)} & Objet indirect : \all{Ich helfe dem Sprecher.} \\
 & & Verbes \emph{helfen, danken, gefallen, folgen, glauben, gehören} : \all{Wir danken den Organisatoren.} \\
 & & Prépositions \emph{aus, bei, mit, nach, seit, von, zu} : \all{am (an dem) Mittwoch}, \all{von 10 bis 12 Uhr} \\
\hline
\end{tabularx}
\end{center}
\textbf{Attention :} Au pluriel datif, les noms prennent un \emph{-n} s'ils ne l'ont pas déjà (\all{den Hörern}, \all{den Veranstaltungen}, mais \all{den Meldungen}).
\end{propriete}

\begin{propriete}[Rappel : Prépositions de temps – sens et cas]
\begin{center}
\begin{tabularx}{\linewidth}{|l|l|l|X|}
\hline
\rowcolor{popblueL} \textbf{Préposition} & \textbf{Cas} & \textbf{Sens} & \textbf{Exemples} \\
\hline
\cle{am} (= an dem) & Datif & Jour calendaire, date & \all{am Montag}, \all{am 15. Mai} \\
\cle{um} & Datif & Heure précise & \all{um 14 Uhr}, \all{um 09.30 Uhr} \\
\cle{von … bis} & Datif & Début – fin (étendue) & \all{von Montag bis Freitag}, \all{von 10 bis 12 Uhr} \\
\cle{seit} & Datif & Depuis (début passé, action encore en cours) & \all{seit 9 Uhr}, \all{seit Montag} \\
\cle{bis} & Accusatif (sans article) / Datif & Jusqu'à (échéance) & \all{bis Freitag}, \all{bis 16 Uhr} \\
\cle{ab} & Datif & À partir de (début futur) & \all{ab 10 Uhr}, \all{ab Montag} \\
\cle{in} + Datif & Datif & Dans (durée future) & \all{in einer Stunde}, \all{in zwei Tagen} \\
\hline
\end{tabularx}
\end{center}
\end{propriete}

\begin{definition}[Lexique utile – Mots de la leçon (article, pluriel, traduction)]
\begin{itemize}
    \item \cle{die Durchsage, die Durchsagen} – l'annonce, le communiqué (radio)
    \item \cle{der Sprecher, die Sprecher} – le speaker, l'animateur (radio)
    \item \cle{die Meldung, die Meldungen} – la nouvelle, le bulletin d'information
    \item \cle{die Veranstaltung, die Veranstaltungen} – l'événement, la manifestation
    \item \cle{die Mitteilung, die Mitteilungen} – la communication, l'avis
    \item \cle{der Hörer, die Hörer} – l'auditeur
    \item \cle{der Termin, die Termine} – la date, le rendez-vous
    \item \cle{der Ort, die Orte} – le lieu
    \item \cle{der Veranstalter, die Veranstalter} – l'organisateur
    \item \cle{die Absage, die Absagen} – l'annulation
    \item \cle{verschieben} (verbe fort, \emph{a} → \emph{ie}) – reporter, déplacer (date) : \all{Das Turnier wird verschoben.}
\end{itemize}
\end{definition}

\begin{methode}[Structure type d'un communiqué radio (Modell)]
\begin{enumerate}
    \item \textbf{Ouverture} : \all{Liebe Zuhörer!} / \all{Liebe Mitschülerinnen und Mitschüler!}
    \item \textbf{Identification} : \all{Hier spricht der Radioclub des Lycée Bilingue de Yaoundé.} / \all{Das ist eine Durchsage der Schulleitung.}
    \item \textbf{Annonce de l'événement (Was?)} : \all{Wir möchten Sie über ein wichtiges Ereignis informieren.} + nom au nominatif/accusatif.
    \item \textbf{Détails temporels (Wann?)} : \all{Die Veranstaltung findet am Mittwoch, dem 15. Mai, um 14 Uhr statt.} (prépositions de temps + datif).
    \item \textbf{Détails spatiaux (Wo?)} : \all{Der Ort ist der Sportplatz hinter dem Hauptgebäude.} (préposition \emph{in} + datif : \all{in der Aula}).
    \item \textbf{Informations complémentaires / Appel à participation} : \all{Alle Klassen sind herzlich eingeladen.} / \all{Wir brauchen noch Helfer.}
    \item \textbf{Remerciements (verbe + datif)} : \all{Wir danken den Organisatoren für die Vorbereitung.}
    \item \textbf{Clôture} : \all{Wir wünschen Ihnen einen schönen Tag und gute Unterhaltung!} / \all{Bleiben Sie dran!}
\end{enumerate}
\end{methode}

\courssub{Proposition de production et corrigé commenté}
\begin{exoresolu}[Communiqué modèle – Tournoi de football inter-classes]
\textbf{Production modèle (groupe 3 – Sprecher : Jean) :}
\all{
Liebe Zuhörer! \\
Hier spricht der Radioclub des Lycée Bilingue de Yaoundé. \\
Wir haben eine wichtige Meldung für Sie. \\
Am Mittwoch, dem 15. Mai, findet das große Fußballturnier der Klassen statt. \\
Das Spiel beginnt um 14 Uhr auf dem Kunstrasenplatz hinter dem Hauptgebäude. \\
Von 13 Uhr bis 16 Uhr ist der Platz für den Verkehr gesperrt. \\
Wir danken den Organisatoren und den Schiedsrichtern für ihre Hilfe. \\
Wir wünschen allen Teams viel Erfolg und den Zuschauern spannende Spiele! \\
Bleiben Sie dran!
}

\tcblower

\textbf{Commentaire des choix de langue :}
\begin{enumerate}
    \item \textbf{Ouverture standard} : \all{Liebe Zuhörer!} (vocatif, nominatif pluriel) – formule figée, polie, inclusive.
    \item \textbf{Identification de la source} : \all{Hier spricht der Radioclub …} – structure \emph{hier + verbe + sujet (nominatif)} typique des annonces radio. \all{des Lycée Bilingue de Yaoundé} : génitif apposé (niveau A2, accepté comme bloc figé).
    \item \textbf{Annonce} : \all{Wir haben eine wichtige Meldung für Sie.} – \cle{eine Meldung} (accusatif féminin, objet direct de \emph{haben}), \cle{für Sie} (préposition \emph{für} + accusatif). \cle{wichtig} (adjectif fort, accusatif féminin : \emph{-e}).
    \item \textbf{Temps – Jour} : \all{Am Mittwoch, dem 15. Mai} – \cle{am} = \emph{an dem} (préposition \emph{an} + datif masculin/neutre pour les jours). \cle{dem 15. Mai} : apposition de date au datif (article \emph{dem} + nombre ordinal \emph{15.} + mois). \textbf{Piège :} pas \emph{an den 15. Mai} (accusatif) car pas de mouvement vers le jour.
    \item \textbf{Temps – Heure} : \all{um 14 Uhr} – \cle{um} + datif pour l'heure précise.
    \item \textbf{Lieu} : \all{auf dem Kunstrasenplatz hinter dem Hauptgebäude} – \emph{auf} + datif (position statique : \emph{wo?}). \cle{dem Kunstrasenplatz} (datif masculin), \cle{dem Hauptgebäude} (datif neutre, préposition \emph{hinter} + datif pour position).
    \item \textbf{Durée / Étendue} : \all{Von 13 Uhr bis 16 Uhr} – \cle{von … bis} + datif (heures). \all{ist der Platz für den Verkehr gesperrt} – passif statal (\emph{sein} + participe passé) : \cle{der Platz} (sujet nominatif), \cle{für den Verkehr} (préposition \emph{für} + accusatif masculin).
    \item \textbf{Remerciements – Verbe au datif} : \all{Wir danken den Organisatoren und den Schiedsrichtern für ihre Hilfe.} – \cle{danken} + \textbf{datif} (\cle{den Organisatoren}, \cle{den Schiedsrichtern} – pluriel datif en \emph{-n}). \cle{für ihre Hilfe} : \emph{für} + accusatif (\cle{die Hilfe}).
    \item \textbf{Souhaits – Accusatif} : \all{Wir wünschen allen Teams viel Erfolg und den Zuschauern spannende Spiele!} – \cle{wünschen} + double objet : \cle{allen Teams} (datif pluriel, article \emph{allen}), \cle{viel Erfolg} (accusatif singulier, pas d'article), \cle{den Zuschauern} (datif pluriel), \cle{spannende Spiele} (accusatif pluriel, adjectif fort \emph{-e}).
    \item \textbf{Clôture} : \all{Bleiben Sie dran!} – impératif formel (\emph{Sie}-Form) + particule \emph{dran} (figé pour „rester à l'écoute“).
\end{enumerate}

\textbf{Variantes acceptées (selon l'événement choisi) :}
\begin{itemize}
    \item \textit{Journée culturelle} : \all{Am Freitag, dem 24. Mai, findet unser Kulturtag statt. Von 9 bis 16 Uhr in der Aula. Wir helfen den Lehrern bei der Vorbereitung.}
    \item \textit{Visite Goethe-Institut} : \all{Am Dienstag, dem 11. Juni, um 10 Uhr besucht eine Delegation des Goethe-Instituts unsere Schule. Im Amphitheater informieren sie über Stipendien. Wir danken dem Direktor für die Einladung.}
    \item \textit{Fermeture route} : \all{Von Montag, dem 3. Juni, bis Freitag, dem 7. Juni, ist die Hauptstraße vor dem Lycée gesperrt. Eine Umleitung führt über die Avenue Kennedy. Wir bitten die Autofahrer um Verständnis.}
\end{itemize}
\end{exoresolu}

\courssub{Bilan de l'activité}
\begin{aretenir}[Ce que vous devez retenir de cette activité d'intégration]
\begin{itemize}
    \item Un communiqué radio (\cle{die Durchsage}) suit une \textbf{structure rituelle} : ouverture formelle → identification de l'émetteur → annonce claire (quoi ?) → détails temporels (quand ?) et spatiaux (où ?) → informations pratiques / appel à action → remerciements (souvent avec verbe au \textbf{datif} : \cle{danken}, \cle{helfen}) → formule de clôture.
    \item Les \textbf{prépositions de temps} structurent l'information chronologique : \cle{am} + jour (datif), \cle{um} + heure (datif), \cle{von … bis} + étendue (datif), \cle{seit} + début passé (datif), \cle{ab} + début futur (datif), \cle{bis} + échéance (accusatif/datif). \textbf{Erreur fréquente :} confondre \emph{am} (jour) et \emph{um} (heure) ; oublier le datif après \emph{von, bis, seit, ab}.
    \item L'\textbf{accusatif} marque l'objet direct (\cle{die Meldung}, \cle{die Durchsage}) et suit les prépositions \emph{durch, für, gegen, ohne, um}. Le \textbf{datif} marque l'objet indirect / bénéficiaire (\cle{dem Sprecher}, \cle{den Organisatoren}) et suit les verbes \emph{helfen, danken, gefallen} ainsi que les prépositions \emph{aus, bei, mit, nach, seit, von, zu, an (lieu/temps), in (lieu), auf (lieu), hinter, vor, unter, über, zwischen} (position statique \emph{wo?}).
    \item À l'oral, le \cle{Sprecher} adopte un \textbf{rythme métronomique}, articule les noms composés (\all{Ver-an-stal-tung}, \all{Kun-stra-sen-platz}), marque les frontières syntaxiques par des pauses (virgules = courte pause, point = pause longue + intonation descendante).
    \item L'\textbf{écoute active} (fiche \emph{Wer? Was? Wann? Wo?}) entraîne la compréhension sélective : repérer les mots-clés (noms propres, chiffres, prépositions de temps) sans tout comprendre.
\end{itemize}
\end{aretenir}

\begin{savaistu}[La radio au Cameroun et dans l'espace germanophone]
Au Cameroun, la radio reste le média le plus accessible : la \emph{CRTV Radio} (publique) et de nombreuses radios privées (\emph{Radio Balafon, Radio Environnement, Magic FM}) émettent en français, anglais et langues nationales. Les \emph{radios communautaires} (ex. \emph{Radio Medumba} à Bafang) jouent un rôle clé en zone rurale. En Allemagne, le système \emph{dual} combine radios publiques (\emph{ARD} avec \emph{Deutschlandfunk}, \emph{WDR}, \emph{NDR}…) financées par la redevance (\cle{der Rundfunkbeitrag}) et radios privées financées par la publicité. L'\emph{Österreichischer Rundfunk (ORF)} en Autriche et la \emph{Schweizer Radio und Fernsehen (SRF)} en Suisse fonctionnent sur un modèle similaire. Le terme \cle{der Hörer} (l'auditeur) est au cœur de la conception \emph{öffentlich-rechtlich} : informer, éduquer, divertir (\emph{informieren, bilden, unterhalten}). Au lycée, le \emph{Radioclub} reproduit ce service public à l'échelle de l'établissement.
\end{savaistu}

\begin{attention}[Erreurs fréquentes observées – Pièges à éviter]
\begin{itemize}
    \item \textbf{Confusion accusatif / datif après prépositions de lieu :} \all{in der Aula} (datif, position \emph{wo?}) mais \all{in die Aula} (accusatif, direction \emph{wohin?}). Ici, toujours \textbf{datif} (lieu de l'événement).
    \item \textbf{Oubli du \emph{-n} au pluriel datif :} \all{den Organisator} $\times$ → \all{den Organisatoren} \checkmark ; \all{den Zuschauer} $\times$ → \all{den Zuschauern} \checkmark.
    \item \textbf{Mauvaise préposition de temps :} \all{an 14 Uhr} $\times$ → \all{um 14 Uhr} \checkmark ; \all{in Montag} $\times$ → \all{am Montag} \checkmark.
    \item \textbf{Ordre des mots (TeKaMoLo) :} \all{Wir danken für die Hilfe den Organisatoren} $\times$ → \all{Wir danken den Organisatoren für die Hilfe} \checkmark (datif avant prépositionnel).
    \item \textbf{Majuscules manquantes :} \all{die durchsage} $\times$ → \all{die Durchsage} \checkmark ; \all{am mittwoch} $\times$ → \all{am Mittwoch} \checkmark.
    \item \textbf{Faux ami :} \all{aktuell} ne signifie pas „actuel“ mais „d'actualité, récent“ ; \all{eventuell} = „éventuellement“, pas „événement“ (\cle{die Veranstaltung}).
    \item \textbf{Prononciation :} \all{ß} = [s] fort (pas [b]) ; \all{eu} = [ɔʏ] (\all{Veranstaltung}) ; \all{sp/st} en début de syllabe = [ʃp/ʃt] (\all{Sprecher} [ˈʃpʁɛçɐ], \all{Stunde} [ˈʃtʊndə]).
\end{itemize}
\end{attention}

\begin{experience}[Prolongement possible – Projet « Radio Schulhof »]
Sur plusieurs séances, les groupes affinent leur communiqué, l'enregistrent sur smartphone (application \emph{Voice Recorder} ou \emph{Audacity}), ajoutent un indicatif musical court (sonnerie d'école, extrait de \emph{„Ein bißchen Frieden“} de Nicole, hymne non officiel des élèves germanistes), et publient les fichiers audio sur l'espace numérique de l'établissement (Padlet, Google Drive, WhatsApp groupe classe). Une séance d'écoute comparée permet de travailler l'auto-correction phonétique et la métacognition sur la production orale.
\end{experience}

\newpage

\coursec{Méthodes et exercices résolus}

\courssub{Méthodes et exercices résolus}

\begin{methode}[Comprendre un communiqué radio (écoute/lecture globale et détaillée)]
\textbf{Objectif :} Extraire l'information essentielle (qui, quoi, quand, où, action attendue) d'un texte court, oral ou écrit, typique des annonces radiophoniques.
\begin{enumerate}
    \item \textbf{Identifier le type et la source :} Repérer l'indicatif de la station (\all{Hier ist Radio Campus Yaoundé.}) et le mot-clé \cle{die Durchsage} / \cle{die Meldung}.
    \item \textbf{Écouter/Lire pour le sens global (gist) :} Quel est le \textbf{sujet principal} ? (Invitation, avertissement, information pratique ?). Ne pas bloquer sur les mots inconnus.
    \item \textbf{Repérer les \cle{W-Fragen} (questions clés) :}
    \begin{itemize}
        \item \all{Wer?} (Qui parle ? Qui est concerné ?) $\rightarrow$ \all{der Sprecher}, \all{die Zuhörer}, \all{die Schüler}.
        \item \all{Was?} (Quoi ? Quel événement ?) $\rightarrow$ \all{die Veranstaltung}, \all{der Informationstag}.
        \item \all{Wann?} (Quand ?) $\rightarrow$ Prépositions de temps (\all{am}, \all{um}, \all{von ... bis}).
        \item \all{Wo?} (Où ?) $\rightarrow$ \all{im Goethe-Institut}, \all{in der Aula}.
        \item \all{Was soll man tun?} (Action attendue) $\rightarrow$ Verbes modaux ou infinitifs (\all{sich anmelden}, \all{kommen}, \all{zuhören}).
    \end{itemize}
    \item \textbf{Vérifier les chiffres et noms propres :} Dates, heures, numéros de téléphone, adresses e-mail (souvent dictés lentement).
    \item \textbf{Reformuler en français :} Résumer en 2-3 phrases pour valider la compréhension.
\end{enumerate}
\textbf{Réflexe Bac :} Les questions de compréhension portent souvent sur l'\textbf{intention communicative} (informer, inviter, alerter) et les \textbf{détails chiffrés} (heure, date, lieu). Souligne-les dans le texte.
\end{methode}

\begin{exoresolu}[Compréhension détaillée : « Informationstag am Goethe-Institut »]
\textbf{Texte support (lu deux fois à l'oral / lu en classe) :}
\begin{textebox}[Communiqué radio — Radio Campus Yaoundé]
„Hier ist Radio Campus Yaoundé. Wir bringen eine wichtige Durchsage für alle Schüler der Klassen Première und Terminale. Das Goethe-Institut Kamerun in Yaoundé veranstaltet am Samstag, dem 15. Juni, einen Informationstag zum Thema ‚Studieren in Deutschland‘. Die Veranstaltung beginnt um 10.00 Uhr und endet gegen 16.00 Uhr. Der Referent Herr Dr. Müller hält um 10.30 Uhr einen Vortrag über die Zulassungsvoraussetzungen. Anschließend beantwortet er Ihre Fragen. Interessierte Schülerinnen und Schüler melden sich bitte bis zum 10. Juni per E-Mail an: info@goethe.de. Der Eintritt ist frei. Wir wünschen allen Teilnehmern viel Erfolg und interessante Einblicke.“
\end{textebox}

\textbf{Traduction du texte :} « Ici Radio Campus Yaoundé. Nous vous apportons une annonce importante pour tous les élèves des classes de Première et de Terminale. Le Goethe-Institut Cameroun à Yaoundé organise le samedi 15 juin une journée d'information sur le thème "Étudier en Allemagne". La manifestation commence à 10 h et se termine vers 16 h. L'intervenant, M. Dr Müller, donne à 10 h 30 une conférence sur les conditions d'admission. Ensuite, il répond à vos questions. Les élèves intéressés sont priés de s'inscrire avant le 10 juin par e-mail : info@goethe.de. L'entrée est gratuite. Nous souhaitons à tous les participants beaucoup de succès et des découvertes intéressantes. »

\textbf{Exercice :} Répondez en allemand (phrases complètes) aux questions suivantes :
\begin{enumerate}
    \item Wer sendet diese Durchsage?
    \item Was für eine Veranstaltung findet statt?
    \item Wann und wo findet der Informationstag statt?
    \item Was macht Herr Dr. Müller um 10.30 Uhr?
    \item Wie und bis wann muss man sich anmelden?
    \item Wie viel kostet der Eintritt?
\end{enumerate}

\tcblower
\textbf{Solution détaillée et commentée :}

\textbf{Analyse grammaticale et lexicale préalable :}
\begin{itemize}
    \item \all{die Durchsage} (f., Pl. \all{die Durchsagen}) = l'annonce, le communiqué (radio/gare).
    \item \all{veranstalten} (faible, transitif + Akkusativ) = organiser. \all{die Veranstaltung} = la manifestation, l'événement.
    \item \all{am Samstag, dem 15. Juni} : \all{an + dem = am} (datif de temps, jour calendaire) ; \all{dem 15. Juni} = apposition au datif.
    \item \all{um 10.00 Uhr} : \all{um} + heure précise (accusatif).
    \item \all{sich anmelden} (verbe réfléchi, séparable) : \all{melden ... an} ; ici avec \all{per} + moyen (\all{per E-Mail}).
    \item \all{bis zum 10. Juni} : \all{bis} seul régit l'accusatif (\all{bis Freitag}), mais devant un article on dit \all{bis zu} + \textbf{datif} $\rightarrow$ \all{bis zum} (= \all{bis zu dem}).
    \item \all{Der Eintritt ist frei} = L'entrée est gratuite (faux ami : \all{frei} signifie ici « gratuit », et non « libre »).
\end{itemize}

\textbf{Réponses rédigées (niveau A2) :}
\begin{enumerate}
    \item \all{Radio Campus Yaoundé sendet diese Durchsage.} « Radio Campus Yaoundé diffuse cette annonce. » (Sujet en position 1, verbe conjugué en position 2, accusatif \all{diese Durchsage}.)
    \item \all{Ein Informationstag zum Thema ‚Studieren in Deutschland‘ findet statt.} « Une journée d'information sur le thème "Étudier en Allemagne" a lieu. » (Verbe séparable \all{stattfinden} : particule \all{statt} en fin de phrase.)
    \item \all{Der Informationstag findet am Samstag, dem 15. Juni, um 10.00 Uhr im Goethe-Institut in Yaoundé statt.} « La journée d'information a lieu le samedi 15 juin à 10 h au Goethe-Institut à Yaoundé. » (Ordre : jour — date — heure — lieu — particule \all{statt}.)
    \item \all{Herr Dr. Müller hält um 10.30 Uhr einen Vortrag über die Zulassungsvoraussetzungen.} « M. Dr Müller donne à 10 h 30 une conférence sur les conditions d'admission. » (Accusatif : \all{einen Vortrag} ; \all{über} + accusatif : \all{die Zulassungsvoraussetzungen}.)
    \item \all{Man muss sich bis zum 10. Juni per E-Mail anmelden.} « Il faut s'inscrire avant le 10 juin par e-mail. » (Impersonnel \all{man} + modal \all{müssen} + infinitif réfléchi \all{sich anmelden} en fin de phrase.)
    \item \all{Der Eintritt ist frei.} « L'entrée est gratuite. » (Adjectif prédicatif \all{frei}, sans majuscule.)
\end{enumerate}

\textbf{Conclusion :} Ce communiqué suit la structure canonique : Indicatif $\rightarrow$ Cible (\all{für alle Schüler...}) $\rightarrow$ Information noyau (Quoi, Quand, Où) $\rightarrow$ Détails (Programme) $\rightarrow$ Modalités d'action (\all{Anmeldung}) $\rightarrow$ Coût $\rightarrow$ Formule de politesse finale.
\end{exoresolu}

\begin{methode}[Maîtriser le vocabulaire thématique : familles de mots, articles, pluriels]
\textbf{Objectif :} Construire un lexique actif et orthographié correctement (majuscules, Umlaute, pluriels) autour des médias et de la radio.
\begin{enumerate}
    \item \textbf{Noter le nom avec son article et son pluriel :} \all{das Radio, die Radios} ; \all{die Sendung, die Sendungen} ; \all{der Sprecher, die Sprecher} ; \all{die Meldung, die Meldungen}.
    \item \textbf{Travailler par familles morphologiques (Wortfamilien) :}
    \begin{itemize}
        \item Verbe de base $\rightarrow$ noms d'action, d'agent, d'objet.
        \item \all{senden} $\rightarrow$ \all{die Sendung} (émission), \all{der Sender} (émetteur, station), \all{der Absender} (expéditeur).
        \item \all{melden} $\rightarrow$ \all{die Meldung} (communication, bulletin), \all{die Anmeldung} (inscription), \all{sich anmelden} (s'inscrire), \all{sich abmelden} (se désinscrire).
        \item \all{sprechen} $\rightarrow$ \all{der Sprecher} (locuteur), \all{die Sprecherin} (locutrice), \all{die Sprache} (langue), \all{der Vortrag} (exposé, conférence), \all{besprechen} (discuter).
    \end{itemize}
    \item \textbf{Associer les verbes et leurs régimes (Kasusregierung) :}
    \begin{itemize}
        \item \all{zuhören} + \textbf{datif} : \all{Ich höre dem Sprecher zu.} « J'écoute attentivement le speaker. »
        \item \all{sich anmelden} + \textbf{bei} + datif (personne/institution) / \textbf{für} + accusatif (événement) / \textbf{per} + accusatif (moyen).
        \item \all{teilnehmen} + \textbf{an} + datif : \all{an der Veranstaltung teilnehmen} « participer à la manifestation ».
    \end{itemize}
    \item \textbf{Entraînement :} Transformer \all{der Sprecher} $\rightarrow$ \all{die Sprecherin} (féminin) ; \all{die Meldung} $\rightarrow$ \all{melden} (verbe) ; \all{veranstalten} $\rightarrow$ \all{der Veranstalter} (organisateur).
\end{enumerate}
\textbf{Astuce :} Crée une fiche « Radio \& Médias » avec 3 colonnes : \textbf{Nom (article + pluriel)} | \textbf{Verbe associé + cas} | \textbf{Exemple en contexte}.
\end{methode}

\begin{exoresolu}[Lexique : familles de mots, articles, cas régis]
\textbf{Exercice :} Complétez le tableau ci-dessous. Donnez le pluriel, le verbe de la famille et la préposition/cas régi par le verbe principal s'il y a lieu.

\begin{center}
\begin{tabularx}{\linewidth}{|X|X|X|X|}
\hline
\rowcolor{popblueL} \textbf{Nom (singulier + pluriel)} & \textbf{Verbe de la famille} & \textbf{Préposition + cas} & \textbf{Phrase exemple (allemand)} \\
\hline
die Durchsage, \dots & \dots & \dots & \all{Wir hören die Durchsage.} \\
\hline
\dots , die Sprecher & sprechen & \dots + Dativ & \all{Ich höre \textbf{dem Sprecher} zu.} \\
\hline
die Anmeldung, \dots & sich anmelden & bei + \dots ; für + \dots & \all{Ich melde mich \textbf{für den Kurs} an.} \\
\hline
\dots , die Veranstaltungen & veranstalten & \dots + Akkusativ & \all{Der Verein veranstaltet \textbf{ein Fest}.} \\
\hline
die Meldung, \dots & melden & \dots + Akkusativ & \all{Er meldet \textbf{den Schaden}.} \\
\hline
das Radio, \dots & senden & \dots + Akkusativ & \all{Der Sender sendet \textbf{Musik}.} \\
\hline
\end{tabularx}
\end{center}

\tcblower
\textbf{Solution détaillée :}

\begin{center}
\begin{tabularx}{\linewidth}{|X|X|X|X|}
\hline
\rowcolor{popblueL} \textbf{Nom (singulier + pluriel)} & \textbf{Verbe de la famille} & \textbf{Préposition + cas} & \textbf{Phrase exemple} \\
\hline
die Durchsage, \textbf{die Durchsagen} & \textbf{ansagen / durchsagen} & --- (transitif direct : Akkusativ) & \all{Wir hören die Durchsage.} \\
\hline
\textbf{der Sprecher}, die Sprecher & sprechen & \textbf{zuhören} + \textbf{Dativ} & \all{Ich höre dem Sprecher zu.} \\
\hline
die Anmeldung, \textbf{die Anmeldungen} & \textbf{sich anmelden} (réfl. séparable) & \textbf{bei} + \textbf{Dativ} (institution) ; \textbf{für} + \textbf{Akkusativ} (événement) & \all{Ich melde mich für den Kurs an.} \\
\hline
\textbf{die Veranstaltung}, die Veranstaltungen & veranstalten & --- + \textbf{Akkusativ} (transitif direct) & \all{Der Verein veranstaltet ein Fest.} \\
\hline
die Meldung, \textbf{die Meldungen} & melden & --- + \textbf{Akkusativ} (transitif direct) & \all{Er meldet den Schaden.} \\
\hline
das Radio, \textbf{die Radios} & senden & --- + \textbf{Akkusativ} (transitif direct) & \all{Der Sender sendet Musik.} \\
\hline
\end{tabularx}
\end{center}

\textbf{Justifications grammaticales :}
\begin{enumerate}
    \item \textbf{Pluriels :} les noms féminins en \all{-e}, \all{-ung}, \all{-tion} prennent \all{-n/-en} au pluriel : \all{die Durchsagen, die Anmeldungen, die Veranstaltungen, die Meldungen}. \all{der Sprecher} $\rightarrow$ \all{die Sprecher} (pluriel sans marque ni Umlaut, comme \all{der Lehrer} $\rightarrow$ \all{die Lehrer}). \textbf{Attention :} les mots étrangers en \all{-o} prennent \all{-s} : \all{das Radio} $\rightarrow$ \all{die Radios}.
    \item \textbf{Cas régis :}
    \begin{itemize}
        \item \all{zuhören} (écouter attentivement) est intransitif et régit le \textbf{datif} : \all{Ich höre dem Sprecher zu.} \emph{Ne pas confondre avec \all{hören} (entendre) + accusatif : \all{Ich höre den Sprecher.}}
        \item \all{sich anmelden für} + accusatif (l'événement visé) ; \all{sich anmelden bei} + datif (l'organisme).
        \item \all{veranstalten, melden, senden} sont transitifs directs $\rightarrow$ \textbf{accusatif}.
    \end{itemize}
\end{enumerate}
\end{exoresolu}

\begin{methode}[Appliquer l'accusatif et le datif (cas, articles, pronoms, prépositions)]
\textbf{Objectif :} Choisir le bon cas (Akkusativ ou Dativ) selon la fonction syntaxique (objet direct / objet indirect / préposition) et décliner correctement articles, adjectifs et pronoms.
\begin{enumerate}
    \item \textbf{Identifier le verbe et son régime (Kasusregierung) :}
    \begin{itemize}
        \item \textbf{Verbes transitifs directs (accusatif) :} \all{haben, lesen, hören, sehen, kaufen, besuchen, veranstalten, melden, senden, verstehen}. Question : \all{Wen oder was?} (Qui ? Quoi ?)
        \item \textbf{Verbes ditransitifs (datif + accusatif) :} \all{geben, schenken, zeigen, erklären, empfehlen, schicken, schreiben, sagen, erzählen}. Ordre standard avec deux noms : \textbf{datif avant accusatif}. \all{Ich gebe \textbf{dem Mann} (Dat.) \textbf{das Buch} (Akk.).} « Je donne le livre à l'homme. »
        \item \textbf{Verbes « datif seul » (pièges fréquents) :} \all{helfen, danken, folgen, glauben, gratulieren, antworten, gefallen, gehören, passen, schmecken, fehlen, begegnen}. Question : \all{Wem?} (À qui ?)
    \end{itemize}
    \item \textbf{Appliquer les prépositions :}
    \begin{itemize}
        \item \textbf{Prépositions \cle{accusatives} (toujours Akk.) :} \all{durch, für, gegen, ohne, um, bis, entlang}. Moyen mnémotechnique : \textbf{FUDGOB} (für, um, durch, gegen, ohne, bis).
        \item \textbf{Prépositions \cle{datives} (toujours Dat.) :} \all{aus, bei, mit, nach, seit, von, zu, gegenüber}. Moyen mnémotechnique : \textbf{bei mit nach seit von zu aus} (à réciter comme une comptine).
        \item \textbf{Prépositions \cle{à double cas} (Wechselpräpositionen) :} \all{an, auf, hinter, in, neben, über, unter, vor, zwischen}.
        \textbf{Règle :} mouvement, direction $\rightarrow$ \all{Wohin?} = \textbf{accusatif}. Position, état $\rightarrow$ \all{Wo?} = \textbf{datif}.
    \end{itemize}
    \item \textbf{Décliner l'article défini :}
    \begin{center}
    \begin{tabular}{|c|c|c|c|c|}\hline
    \rowcolor{popblueL} & \textbf{Masculin} & \textbf{Féminin} & \textbf{Neutre} & \textbf{Pluriel} \\ \hline
    \textbf{Nominativ} & der & die & das & die \\ \hline
    \textbf{Akkusativ} & \textbf{den} & die & das & die \\ \hline
    \textbf{Dativ} & \textbf{dem} & \textbf{der} & \textbf{dem} & \textbf{den} (+ n au nom) \\ \hline
    \end{tabular}
    \end{center}
    \item \textbf{Vérifier les pronoms personnels :} \all{ich} $\rightarrow$ \all{mich} (Akk.) / \all{mir} (Dat.) ; \all{du} $\rightarrow$ \all{dich} / \all{dir} ; \all{er} $\rightarrow$ \all{ihn} / \all{ihm} ; \all{sie} $\rightarrow$ \all{sie} / \all{ihr} ; \all{wir} $\rightarrow$ \all{uns} / \all{uns}.
\end{enumerate}
\textbf{Réflexe Bac :} Si le verbe est \all{helfen}, \textbf{jamais} d'accusatif : \all{*ich helfe ihn} $\rightarrow$ \all{ich helfe \textbf{ihm}}. Si \all{in} + mouvement $\rightarrow$ \all{in das} = \all{ins} (Akk.) ; si position $\rightarrow$ \all{in dem} = \all{im} (Dat.).
\end{methode}

\begin{exoresolu}[Cas : accusatif ou datif dans le contexte radio]
\textbf{Exercice :} Mettez les mots entre parenthèses au bon cas (accusatif ou datif). Attention aux articles, pronoms et adjectifs.

\begin{enumerate}
    \item Der Sprecher dankt \_\_\_\_\_\_\_\_\_\_\_\_ (der Zuhörer, masc., sing.) für die Aufmerksamkeit.
    \item Wir hören \_\_\_\_\_\_\_\_\_\_\_\_ (die Meldung, fém., sing.) im Radio.
    \item Ich helfe \_\_\_\_\_\_\_\_\_\_\_\_ (mein Freund, masc., sing.) bei der Anmeldung.
    \item Die Durchsage gilt \_\_\_\_\_\_\_\_\_\_\_\_ (alle Schüler, pluriel).
    \item Wir fahren \_\_\_\_\_\_\_\_\_\_\_\_ (in die Schule, fém., sing.) — mouvement.
    \item Wir sind \_\_\_\_\_\_\_\_\_\_\_\_ (in der Schule, fém., sing.) — position.
    \item Herr Müller schickt \_\_\_\_\_\_\_\_\_\_\_\_ (ich, pronom) \_\_\_\_\_\_\_\_\_\_\_\_ (die Information, fém., sing.).
    \item Ohne \_\_\_\_\_\_\_\_\_\_\_\_ (dein Radio, neutre, sing.) hören wir nichts.
    \item Seit \_\_\_\_\_\_\_\_\_\_\_\_ (lange Zeit, fém., sing.) gibt es Radio.
    \item Der Veranstalter gibt \_\_\_\_\_\_\_\_\_\_\_\_ (die Teilnehmer, pluriel) \_\_\_\_\_\_\_\_\_\_\_\_ (ein Programmheft, neutre, sing.).
\end{enumerate}

\tcblower
\textbf{Solution détaillée avec justifications :}

\begin{enumerate}
    \item \all{dem Zuhörer} — \textbf{datif}. Le verbe \all{danken} (remercier) régit le datif (\all{Wem danke ich?}). Article défini masc. dat. = \all{dem}. « Le speaker remercie l'auditeur de son attention. »
    \item \all{die Meldung} — \textbf{accusatif}. \all{hören} (entendre) est transitif direct (\all{Wen oder was?}). Article défini fém. acc. = \all{die} (identique au nominatif). « Nous écoutons la communication à la radio. »
    \item \all{meinem Freund} — \textbf{datif}. \all{helfen} régit le datif. Possessif \all{mein} + masc. dat. $\rightarrow$ \all{meinem}. \all{Freund} est un nom fort : pas de \all{-n} au datif singulier. « J'aide mon ami pour l'inscription. »
    \item \all{allen Schülern} — \textbf{datif}. \all{gelten} (concerner, s'appliquer à) + datif. Pluriel datif : \all{alle} $\rightarrow$ \all{allen}, et le nom prend \textbf{-n} : \all{Schülern}. « L'annonce concerne tous les élèves. »
    \item \all{in die Schule} — \textbf{accusatif}. \all{in} + \textbf{mouvement} (\all{Wohin?}). Article défini fém. acc. = \all{die}. « Nous allons à l'école (en voiture). »
    \item \all{in der Schule} — \textbf{datif}. \all{in} + \textbf{position} (\all{Wo?}). Article défini fém. dat. = \all{der}. « Nous sommes à l'école. »
    \item \all{mir} — \textbf{datif} (bénéficiaire) ; \all{die Information} — \textbf{accusatif} (objet direct). \all{schicken} est ditransitif : \all{schicken [Dat. : à qui ?] [Akk. : quoi ?]}. Pronom \all{ich} $\rightarrow$ \all{mir}. « M. Müller m'envoie l'information. »
    \item \all{dein Radio} — \textbf{accusatif}. \all{ohne} $\rightarrow$ toujours accusatif. Neutre acc. : possessif \all{dein} sans terminaison. « Sans ta radio, nous n'entendons rien. »
    \item \all{langer Zeit} — \textbf{datif}. \all{seit} $\rightarrow$ toujours datif. Sans article, l'adjectif prend la déclinaison forte : fém. dat. \all{-er} $\rightarrow$ \all{langer Zeit}. « La radio existe depuis longtemps. »
    \item \all{den Teilnehmern} (dat. pluriel) — \all{ein Programmheft} (acc. singulier). Verbe \all{geben} : datif avant accusatif. Pluriel dat. : \all{die} $\rightarrow$ \all{den} + \textbf{-n} au nom (\all{Teilnehmern}). Neutre acc. : \all{ein} sans terminaison. « L'organisateur donne un livret-programme aux participants. »
\end{enumerate}

\textbf{Points de vigilance :}
\begin{itemize}
    \item \all{Freund} n'est pas un nom faible (pas de \all{-n} au datif singulier), contrairement à \all{der Student, der Mensch, der Junge, der Kunde} qui prennent \all{-n/-en} à tous les cas sauf le nominatif singulier.
    \item \all{seit} + datif (depuis, action qui continue) vs \all{vor} + datif (il y a..., événement ponctuel passé).
    \item Contractions : \all{in das = ins}, \all{in dem = im}, \all{zu dem = zum}, \all{an dem = am}, \all{von dem = vom}, \all{bei dem = beim}. Obligatoires à l'oral, fortement recommandées à l'écrit.
\end{itemize}
\end{exoresolu}

\begin{methode}[Utiliser correctement les prépositions de temps (Zeitpräpositionen)]
\textbf{Objectif :} Situer un événement dans le temps (point, durée, fréquence) avec la préposition et le cas corrects.
\begin{enumerate}
    \item \textbf{Point dans le temps (Wann genau ?) :}
    \begin{itemize}
        \item \textbf{\all{am} (= an dem)} + \textbf{datif} : jours de la semaine, dates calendaires, parties de la journée (\all{am Montag, am 15. Juni, am Morgen}). \emph{Exception :} \all{in der Nacht} (la nuit).
        \item \textbf{\all{um}} + \textbf{accusatif} : heure précise (\all{um 10.00 Uhr, um Mitternacht}).
        \item \textbf{\all{im} (= in dem)} + \textbf{datif} : mois, saisons, années (\all{im Juni, im Sommer, im Jahr 2024}).
    \end{itemize}
    \item \textbf{Durée / intervalle :}
    \begin{itemize}
        \item \textbf{\all{von ... bis (zu)}} : début et fin (\all{von 10 bis 16 Uhr, vom 1. bis zum 15. Juni}). \all{von} + datif ; \all{bis} seul + accusatif, \all{bis zu} + datif.
        \item \textbf{\all{für}} + \textbf{accusatif} : durée prévue (\all{für zwei Wochen}).
        \item \textbf{\all{seit}} / \textbf{\all{ab}} + \textbf{datif} : depuis / à partir de (\all{seit Montag, ab Juni}).
        \item \textbf{\all{bis}} + \textbf{accusatif} (sans article) : jusqu'à (\all{bis Freitag}) ; devant un article : \all{bis zu} + datif (\all{bis zum 10. Juni}).
    \end{itemize}
    \item \textbf{Fréquence / répétition :} adverbes en \all{-s} (minuscule) : \all{montags, morgens, abends, werktags} ; ou \all{täglich, wöchentlich}. Pas de préposition, pas de cas.
    \item \textbf{Antériorité / postériorité :} \all{vor} + datif (avant / il y a) ; \all{nach} + datif (après).
\end{enumerate}
\textbf{Tableau récapitulatif (à connaître par cœur) :}
\begin{center}
\begin{tabularx}{\linewidth}{|X|X|X|X|}
\hline
\rowcolor{popblueL} \textbf{Fonction} & \textbf{Préposition} & \textbf{Cas} & \textbf{Exemple} \\ \hline
Jour / date / partie du jour & \all{am} (an dem) & Datif & \all{am Samstag, am 15. Juni, am Morgen} \\ \hline
Heure précise & \all{um} & Akkusativ & \all{um 10.30 Uhr} \\ \hline
Mois / saison / année & \all{im} (in dem) & Dativ & \all{im Juni, im Winter, im Jahr 2025} \\ \hline
Intervalle (début-fin) & \all{von ... bis (zu)} & Dativ & \all{von 10 bis 12 Uhr, vom 1. bis zum 5. Mai} \\ \hline
Durée prévue & \all{für} & Akkusativ & \all{für zwei Stunden} \\ \hline
Depuis (passé $\rightarrow$ présent) & \all{seit} & Dativ & \all{seit 2020, seit Montag} \\ \hline
À partir de (futur) & \all{ab} & Dativ & \all{ab nächstem Montag} \\ \hline
Jusqu'à (échéance) & \all{bis} / \all{bis zu} & Akk. / Dat. & \all{bis Freitag, bis zum 15. Juni} \\ \hline
Il y a (passé révolu) & \all{vor} & Dativ & \all{vor zwei Tagen} \\ \hline
Après (événement) & \all{nach} & Dativ & \all{nach der Sendung} \\ \hline
\end{tabularx}
\end{center}
\end{methode}

\begin{exoresolu}[Prépositions de temps : complétion et traduction (thème)]
\textbf{Exercice A :} Complétez avec la préposition correcte (et l'article contracté si nécessaire : \all{am, im, vom, zum}).
\begin{enumerate}
    \item Die Sendung beginnt \_\_\_\_\_\_\_\_\_\_\_\_ 10.00 Uhr.
    \item Der Informationstag ist \_\_\_\_\_\_\_\_\_\_\_\_ Samstag, \_\_\_\_\_\_\_\_\_\_\_\_ 15. Juni.
    \item \_\_\_\_\_\_\_\_\_\_\_\_ Juni findet die Veranstaltung statt.
    \item Die Anmeldung ist möglich \_\_\_\_\_\_\_\_\_\_\_\_ 1. Juni \_\_\_\_\_\_\_\_\_\_\_\_ 10. Juni.
    \item Wir senden \_\_\_\_\_\_\_\_\_\_\_\_ 20 Jahren aus Yaoundé.
    \item Bitte kommen Sie \_\_\_\_\_\_\_\_\_\_\_\_ 09.30 Uhr.
    \item Die Ausstellung dauert \_\_\_\_\_\_\_\_\_\_\_\_ zwei Wochen.
\end{enumerate}

\textbf{Exercice B (thème) :} Traduisez en allemand.
\begin{enumerate}
    \item L'annonce dure deux minutes.
    \item Nous écoutons la radio depuis ce matin.
    \item L'événement a lieu le lundi 20 mai à 14 heures.
    \item Vous pouvez vous inscrire jusqu'au 30 mai.
    \item Après l'émission, il y a de la musique.
\end{enumerate}

\tcblower
\textbf{Solution détaillée :}

\textbf{Exercice A (complétion) :}
\begin{enumerate}
    \item \all{um} 10.00 Uhr (\textbf{heure précise} $\rightarrow$ \all{um} + Akk.).
    \item \all{am} Samstag, \all{am} 15. Juni (\textbf{jour / date} $\rightarrow$ \all{an dem} = \all{am} + Dat.).
    \item \all{Im} Juni (\textbf{mois} $\rightarrow$ \all{in dem} = \all{im} + Dat.).
    \item \all{vom} 1. Juni \all{bis zum} 10. Juni (\textbf{intervalle} : \all{von dem} = \all{vom} + Dat. ; \all{bis zu dem} = \all{bis zum} + Dat.).
    \item \all{seit} 20 Jahren (\textbf{depuis} $\rightarrow$ \all{seit} + Dat. pluriel : \all{Jahren}).
    \item \all{um} 09.30 Uhr (\textbf{heure précise}).
    \item \all{zwei} Wochen (\textbf{sans préposition} : le verbe \all{dauern} se construit avec un accusatif de durée direct : \all{Die Ausstellung dauert zwei Wochen.}).
\end{enumerate}

\textbf{Exercice B (traduction — thème) :}
\begin{enumerate}
    \item \all{Die Durchsage dauert zwei Minuten.}
    \emph{Verbe \all{dauern} + accusatif de durée, sans préposition.}
    \item \all{Wir hören seit heute Morgen Radio.}
    \emph{\all{seit} + datif ; \all{heute Morgen} est une expression adverbiale figée. Présent + \all{seit} = « depuis » en français.}
    \item \all{Die Veranstaltung findet am Montag, dem 20. Mai, um 14.00 Uhr statt.}
    \emph{Verbe séparable \all{stattfinden} : \all{findet} en position 2, \all{statt} en fin. Ordre des compléments de temps : jour (\all{am Montag}) — date (\all{dem 20. Mai}) — heure (\all{um 14.00 Uhr}).}
    \item \all{Sie können sich bis zum 30. Mai anmelden.}
    \emph{\all{bis zu dem} = \all{bis zum}. Verbe réfléchi séparable \all{sich anmelden} : avec le modal \all{können}, l'infinitif \all{anmelden} va en fin de phrase.}
    \item \all{Nach der Sendung gibt es Musik.}
    \emph{\all{nach} + datif : \all{der Sendung} (fém.). Structure impersonnelle \all{es gibt} + accusatif (\all{Musik}).}
\end{enumerate}

\textbf{Erreurs fréquentes évitées :}
\begin{itemize}
    \item \all{*an Montag} $\rightarrow$ \all{am Montag}.
    \item \all{*in 10 Uhr} $\rightarrow$ \all{um 10 Uhr}.
    \item \all{*Die Ausstellung dauert für zwei Wochen} $\rightarrow$ \all{Die Ausstellung dauert zwei Wochen} (pas de \all{für} avec \all{dauern}).
    \item \all{*seit 20 Jahre} $\rightarrow$ \all{seit 20 Jahren} (datif pluriel en \textbf{-n}).
    \item \all{*bis der 30. Mai} $\rightarrow$ \all{bis zum 30. Mai} (\all{bis zu} + datif).
\end{itemize}
\end{exoresolu}

\begin{methode}[Rédiger un communiqué radio / résumer / répondre (expression guidée)]
\textbf{Objectif :} Produire un texte court (50-80 mots), structuré, authentique, respectant les conventions du genre « communiqué radio » (style clair, impersonnel, incitatif).
\begin{enumerate}
    \item \textbf{Analyser la consigne :} Quel rôle ? (Speaker, organisateur). Quel public ? (Élèves, grand public). Quel but ? (Informer, inviter, alerter).
    \item \textbf{Structurer selon le modèle standard :}
    \begin{enumerate}
        \item \textbf{Indicatif / accroche :} \all{Hier ist [Sender].} / \all{Achtung, eine wichtige Durchsage!}
        \item \textbf{Cible :} \all{Für alle [Zielgruppe]...} / \all{An alle [Zielgruppe]...}
        \item \textbf{Information noyau (W-Fragen) :} \all{Was? Wann? Wo? Wer?} (Verbes : \all{stattfinden, beginnen, dauern, einladen}).
        \item \textbf{Détails / programme :} \all{Um X Uhr: Vortrag von ... / Pause / Diskussion.}
        \item \textbf{Appel à l'action :} \all{Melden Sie sich an unter ... / Kommen Sie zahlreich! / Infos unter ...}
        \item \textbf{Formule de clôture :} \all{Wir wünschen viel Spaß / viel Erfolg / einen schönen Tag.}
    \end{enumerate}
    \item \textbf{Grammaire du style radio :}
    \begin{itemize}
        \item \textbf{Tournures impersonnelles :} \all{Es findet statt ... / Der Eintritt ist frei.} (Évite le « je/nous » trop familier.)
        \item \textbf{Prépositions contractées :} \all{am, im, zum, zur, vom, beim}.
        \item \textbf{Verbes séparables :} particule en fin de phrase principale (\all{stattfinden, anmelden, zuhören, mitmachen}).
    \end{itemize}
    \item \textbf{Relire (check-list) :} Majuscules aux noms ? Umlaute (ä, ö, ü) ? ß ou ss ? Verbe en position 2 ? Cas corrects (Akk./Dat.) ? Particules séparables en fin de phrase ?
\end{enumerate}
\textbf{Astuce Bac :} Pour un résumé (\all{Zusammenfassung}), utilise le \textbf{présent}, la 3\textsuperscript{e} personne, des phrases courtes, pas de « ich ». Pour une réponse personnelle, utilise « ich » + verbes modaux (\all{ich möchte, ich kann, ich werde}).
\end{methode}

\begin{exoresolu}[Production écrite : rédiger un communiqué radio]
\textbf{Consigne :} Vous êtes le speaker de \all{Radio Lycée Bilingue de Douala}. Rédigez un communiqué (env. 60-80 mots) pour annoncer la finale du tournoi de football inter-classes.
\textbf{Informations à inclure :}
\begin{itemize}
    \item Match : classe de 1\textsuperscript{ère} A contre Terminale A.
    \item Date : samedi 22 juin. Heure : 15 h 00.
    \item Lieu : stade scolaire (\all{das Schulstadion}).
    \item Entrée libre. Buvette au profit de l'association des parents d'élèves.
    \item Appel à venir nombreux.
\end{itemize}

\tcblower
\textbf{Proposition de corrigé (niveau A2) :}

\begin{textebox}[Communiqué radio — proposition de rédaction]
„Hier ist Radio Lycée Bilingue Douala. Wir bringen eine wichtige Durchsage für die gesamte Schülerschaft. Am Samstag, dem 22. Juni, findet das Finale des Fußballturniers statt. Die Klasse 1A spielt gegen die Terminale A. Anpfiff ist um 15.00 Uhr im Schulstadion. Der Eintritt ist frei. Für das leibliche Wohl sorgt ein Getränkestand des Elternvereins. Wir laden alle Schüler, Lehrer und Eltern herzlich ein, die Mannschaften anzufeuern. Kommt zahlreich! Wir wünschen ein spannendes und faires Spiel.“
\end{textebox}

\textbf{Traduction :} « Ici Radio Lycée Bilingue Douala. Nous vous apportons une annonce importante pour toute la communauté scolaire. Le samedi 22 juin aura lieu la finale du tournoi de football. La classe de 1A joue contre la Terminale A. Le coup d'envoi est à 15 h au stade scolaire. L'entrée est gratuite. Une buvette de l'association des parents assure le ravitaillement. Nous invitons chaleureusement tous les élèves, professeurs et parents à encourager les équipes. Venez nombreux ! Nous vous souhaitons un match passionnant et fair-play. »

\textbf{Analyse grammaticale et stylistique du corrigé :}
\begin{enumerate}
    \item \textbf{Indicatif :} \all{Hier ist Radio Lycée Bilingue Douala.} (Formule standard.)
    \item \textbf{Cible :} \all{für die gesamte Schülerschaft} (\all{für} + Akk.). \all{die Schülerschaft} (fém., singulier collectif).
    \item \textbf{Noyau :} \all{Am Samstag, dem 22. Juni, findet das Finale ... statt.}
    \begin{itemize}
        \item \all{Am} = \all{an dem} (Dat., jour/date) ; \all{dem 22. Juni} = apposition au datif.
        \item \all{findet ... statt} : verbe séparable, \all{findet} en position 2, particule \all{statt} en fin.
        \item Sujet : \all{das Finale} (neutre).
    \end{itemize}
    \item \textbf{Détails du match :} \all{Die Klasse 1A spielt gegen die Terminale A.} (\all{spielen gegen} + Akk.).
    \item \textbf{Heure / lieu :} \all{Anpfiff ist um 15.00 Uhr im Schulstadion.}
    \begin{itemize}
        \item \all{um 15.00 Uhr} (heure précise, Akk.).
        \item \all{im Schulstadion} = \all{in dem} (Dat., position \all{Wo?}).
    \end{itemize}
    \item \textbf{Infos pratiques :} \all{Der Eintritt ist frei.} (Adjectif prédicatif.) \all{Für das leibliche Wohl sorgt ein Getränkestand des Elternvereins.}
    \begin{itemize}
        \item \all{sorgen für} + Akk. ; sujet : \all{ein Getränkestand}.
        \item \all{des Elternvereins} : génitif masc. (\all{-s}) — « de l'association des parents ».
    \end{itemize}
    \item \textbf{Appel à l'action :} \all{Wir laden alle Schüler, Lehrer und Eltern herzlich ein, die Mannschaften anzufeuern.}
    \begin{itemize}
        \item \all{einladen} (séparable) : \all{laden ... ein} ; accusatif : \all{alle Schüler, Lehrer und Eltern}.
        \item \all{anfeuern} (séparable) $\rightarrow$ infinitif avec \all{zu} : \all{anzufeuern}, en fin de proposition infinitive.
    \end{itemize}
    \item \textbf{Clôture :} \all{Kommt zahlreich!} (Impératif pluriel.) \all{Wir wünschen ein spannendes und faires Spiel.} (Acc. neutre : \all{ein spannendes ... Spiel}, adjectifs avec terminaison forte \all{-es}.)
\end{enumerate}

\textbf{Variante « résumé » (Zusammenfassung) pour l'examen :}
\all{Radio Lycée Bilingue Douala kündigt das Finale des Fußballturniers an: Am Samstag, dem 22. Juni, spielt die Klasse 1A um 15.00 Uhr im Schulstadion gegen die Terminale A. Der Eintritt ist frei, ein Getränkestand des Elternvereins ist vor Ort. Alle sind herzlich eingeladen, die Mannschaften anzufeuern.}
« Radio Lycée Bilingue Douala annonce la finale du tournoi de football : le samedi 22 juin, la classe de 1A joue à 15 h au stade scolaire contre la Terminale A. L'entrée est gratuite, une buvette de l'association des parents est sur place. Tous sont chaleureusement invités à encourager les équipes. »
\end{exoresolu}

\begin{methode}[Traduire un extrait de communiqué (version / thème ciblé cas et prépositions)]
\textbf{Objectif :} Traduire fidèlement du français vers l'allemand (thème) ou de l'allemand vers le français (version) en gérant les pièges contrastifs (cas, prépositions, verbes séparables, faux amis).
\begin{enumerate}
    \item \textbf{Analyser la phrase source :} identifier le verbe noyau, les compléments (COD, COI, CCL, CCT), les connecteurs logiques.
    \item \textbf{Transposer la structure :}
    \begin{itemize}
        \item Sujet $\rightarrow$ nominatif.
        \item COD $\rightarrow$ \textbf{accusatif} (sauf verbes à datif : \all{helfen, danken, folgen...}).
        \item COI (à qui ?) $\rightarrow$ \textbf{datif}.
        \item Compléments de temps/lieu $\rightarrow$ préposition + cas adéquat.
    \end{itemize}
    \item \textbf{Gérer les pièges francophones :}
    \begin{itemize}
        \item \textbf{Attendre} $\rightarrow$ \all{warten auf} + Akk. (\all{warten} seul est intransitif : « patienter »).
        \item \textbf{Écouter} $\rightarrow$ \all{zuhören} + Dat. (écouter attentivement) vs \all{hören} + Akk. (entendre, percevoir).
        \item \textbf{S'inscrire} $\rightarrow$ \all{sich anmelden} (réfléchi, séparable) + \all{für} (événement, Akk.) / \all{bei} (organisme, Dat.).
        \item \textbf{Participer} $\rightarrow$ \all{teilnehmen an} + Dat.
        \item \textbf{Durée :} « pendant 2 heures » $\rightarrow$ \all{zwei Stunden (lang)} (durée effective) ou \all{für zwei Stunden} (durée prévue).
        \item \textbf{Depuis :} « j'écoute depuis 10 h » $\rightarrow$ \all{Ich höre seit 10 Uhr Radio.} (Présent + \all{seit} + Dat. — jamais de Perfekt ici !)
    \end{itemize}
    \item \textbf{Vérifier l'ordre des mots (Satzbau) :}
    \begin{itemize}
        \item Phrase principale : verbe conjugué en \textbf{position 2}.
        \item Deux compléments-noms : \textbf{datif avant accusatif}. Deux pronoms : \textbf{accusatif avant datif}.
        \item Subordonnée (\all{dass, weil, wenn...}) : verbe conjugué à la \textbf{fin}.
        \item Verbe séparable : particule à la \textbf{fin} (principale) ; verbe complet à la fin (subordonnée).
    \end{itemize}
    \item \textbf{Relecture orthographique :} majuscules (noms, \all{Sie/Ihr} de politesse), Umlaute, \all{ß} (après voyelle longue ou diphtongue : \all{Straße, Fußball}), \all{ss} (après voyelle brève : \all{müssen, Klasse}).
\end{enumerate}
\end{methode}

\begin{exoresolu}[Traduction contrastive : version \& thème « Radio »]
\textbf{Exercice A (version) :} Traduisez en français.
\all{„Seit drei Wochen sendet unser Schulsender jeden Mittwoch eine spezielle Meldung für die neuen Schüler. Der Sprecher erklärt dabei, wie man sich für die AGs anmelden kann. Viele hören regelmäßig zu, aber manche verpassen die Durchsage, weil sie zu spät kommen. Ab nächster Woche gibt es eine Wiederholung am Freitag.“}

\textbf{Exercice B (thème) :} Traduisez en allemand.
\begin{enumerate}
    \item J'écoute cette émission depuis lundi.
    \item Le speaker remercie les auditeurs de leur attention.
    \item Il faut s'inscrire avant le 1er juillet.
    \item Nous attendons le début du programme.
    \item La directrice parle aux élèves de la nouvelle règle.
\end{enumerate}

\tcblower
\textbf{Solution détaillée :}

\textbf{Exercice A (version) :}
\textbf{Traduction :} « Depuis trois semaines, notre radio scolaire diffuse chaque mercredi une information spéciale pour les nouveaux élèves. Le speaker explique à cette occasion comment on peut s'inscrire aux clubs (AG). Beaucoup écoutent régulièrement, mais certains ratent l'annonce parce qu'ils arrivent trop tard. À partir de la semaine prochaine, il y a une rediffusion le vendredi. »

\textbf{Justifications grammaticales :}
\begin{itemize}
    \item \all{Seit drei Wochen sendet...} : « depuis » + \textbf{présent} en allemand (action commencée dans le passé et qui continue). \all{seit} + datif pluriel : \all{drei Wochen}.
    \item \all{jeden Mittwoch} : accusatif de temps (répétition) — « chaque mercredi ».
    \item \all{eine spezielle Meldung} : accusatif (COD de \all{sendet}).
    \item \all{für die neuen Schüler} : \all{für} + accusatif.
    \item \all{Der Sprecher erklärt dabei,} : \all{dabei} = « à cette occasion, lors de cette annonce », adverbe placé après le verbe.
    \item \all{wie man sich ... anmelden kann} : subordonnée indirecte (\all{wie}) $\rightarrow$ verbe conjugué \all{kann} à la fin ; l'infinitif réfléchi \all{sich anmelden} le précède.
    \item \all{viele hören regelmäßig zu} : \all{zuhören} séparable $\rightarrow$ \all{hören ... zu}.
    \item \all{manche verpassen die Durchsage} : \all{verpassen} (rater, manquer) + accusatif.
    \item \all{weil sie zu spät kommen} : subordonnée en \all{weil} $\rightarrow$ verbe \all{kommen} à la fin.
    \item \all{ab nächster Woche} : \all{ab} + datif (\all{nächster Woche}, fém. dat. \all{-er}).
    \item \all{es gibt eine Wiederholung} : structure impersonnelle \all{es gibt} + accusatif.
    \item \all{am Freitag} : \all{an dem} = \all{am} (datif, jour).
\end{itemize}

\textbf{Exercice B (thème) :}
\begin{enumerate}
    \item \all{Ich höre diese Sendung seit Montag.}
    \emph{\all{seit} + datif (\all{Montag} sans article). Présent + \all{seit} = « depuis ».}
    \item \all{Der Sprecher dankt den Zuhörern für ihre Aufmerksamkeit.}
    \emph{\all{danken} + \textbf{datif} : \all{den Zuhörern} (pluriel dat. en \textbf{-n}). \all{für} + accusatif : \all{ihre Aufmerksamkeit}.}
    \item \all{Man muss sich vor dem 1. Juli anmelden.}
    \emph{\all{vor} + datif : \all{dem 1. Juli}. Verbe réfléchi séparable \all{sich anmelden} ; avec le modal \all{müssen}, l'infinitif \all{anmelden} va en fin de phrase.}
    \item \all{Wir warten auf den Beginn des Programms.}
    \emph{\all{warten auf} + \textbf{accusatif} : \all{den Beginn}. Génitif neutre : \all{des Programms} (\all{-s}).}
    \item \all{Die Direktorin spricht mit den Schülern über die neue Regel.}
    \emph{\all{sprechen mit} + \textbf{datif} : \all{den Schülern} (pluriel dat. en \textbf{-n}). \all{über} + \textbf{accusatif} : \all{die neue Regel}.}
\end{enumerate}

\textbf{Points clés du thème :}
\begin{itemize}
    \item \all{warten auf} (attendre) + Akk. — \textbf{jamais} \all{*warten für} ni \all{warten} transitif.
    \item \all{danken} (remercier qqn) + Dat. ; « remercier de/pour qqch » $\rightarrow$ \all{danken für} + Akk. : \all{Ich danke dir für deine Hilfe.}
    \item \all{sprechen mit} (parler à qqn) + Dat. vs \all{sprechen über} (parler de qqch) + Akk.
    \item \all{vor} (temporel « avant ») + Dat. vs \all{bis (zu)} (« jusqu'à »).
\end{itemize}
\end{exoresolu}

\begin{attention}[Les 5 erreurs qui coûtent des points au Bac (ALA17)]
\begin{enumerate}
    \item \textbf{Majuscules oubliées sur les noms :} \all{*die durchsage, *der sprecher, *am samstag, *im juni.} \textbf{Règle :} \cle{tout nom commence par une majuscule} (\all{der Sprecher, die Durchsage, der Samstag, der Juni}). Cela vaut aussi pour les infinitifs substantivés : \all{das Zuhören, das Anmelden}.
    \item \textbf{Confusion \all{seit} / \all{vor} / \all{für} (temps) :}
    \begin{itemize}
        \item \all{*Ich lerne Deutsch vor zwei Jahren.} (Faux si l'action continue.) $\rightarrow$ \all{Ich lerne seit zwei Jahren Deutsch.} (Action qui continue : \textbf{seit} + présent + datif.) \all{vor zwei Jahren} = « il y a deux ans » (événement ponctuel révolu) : \all{Vor zwei Jahren habe ich Deutsch gelernt.} « Il y a deux ans, j'ai appris l'allemand. »
        \item \all{*Ich bleibe für zwei Wochen lang.} (Double marque fautive.) $\rightarrow$ \all{Ich bleibe zwei Wochen (lang).} (Durée effective, sans préposition) ou \all{Ich bleibe für zwei Wochen.} (Durée prévue.)
    \end{itemize}
    \item \textbf{Cas après les verbes « pièges » (datif, pas accusatif) :}
    \all{*Ich helfe ihn / *Ich danke ihn / *Ich folge ihn / *Ich antworte ihn.}
    \textbf{Correct :} \all{ihm} (datif). \textbf{Liste à savoir par cœur :} \all{helfen, danken, folgen, glauben, gratulieren, antworten, gefallen, gehören, passen, schmecken, fehlen, begegnen, zuhören}. \emph{Astuce : la plupart expriment une relation interpersonnelle ou un état, pas une action directe sur un objet.}
    \item \textbf{Verbes à particule séparable : particule oubliée en fin de phrase.}
    \all{*Ich melde mich für den Kurs.} (Il manque \all{an}.) $\rightarrow$ \all{Ich melde mich \textbf{für den Kurs an}.}
    \all{*Die Veranstaltung stattfindet.} (Faux.) $\rightarrow$ \all{Die Veranstaltung \textbf{findet statt}.} En subordonnée, le verbe se recolle à la fin : \all{..., dass die Veranstaltung \textbf{stattfindet}.}
    \item \textbf{Faux amis et calques syntaxiques :}
    \begin{itemize}
        \item \all{Der Eintritt ist frei.} = « L'entrée est gratuite » (et non « libre »). « Je suis libre (disponible) » $\rightarrow$ \all{Ich habe frei.} / \all{Ich habe Zeit.}
        \item \all{*Ich mache mit bei dem Kurs.} (Calque de « je participe à ».) $\rightarrow$ \all{Ich nehme \textbf{an dem Kurs teil}.} (\all{teilnehmen an} + datif.)
        \item \all{*Ich interessiere mich...} (Oubli du réfléchi.) $\rightarrow$ \all{Ich \textbf{interessiere mich} für die Meldung.} (\all{sich interessieren für} + Akk.)
        \item \all{*Ich höre dem Radio.} $\rightarrow$ \all{Ich höre \textbf{Radio}.} (\all{Radio hören} : transitif direct, accusatif, sans article.)
    \end{itemize}
\end{enumerate}
\textbf{Conseil final :} Relisez votre copie en ne vérifiant \textbf{qu'une seule chose à la fois} : 1. Majuscules des noms. 2. Verbe en position 2 / à la fin. 3. Cas (Akk./Dat.) après verbes et prépositions. 4. Particules séparables. 5. Umlaute et ß.
\end{attention}

\newpage

\coursec{Exercices}

\courssub{Je vérifie mes connaissances}

\begin{exercice}[QCM : Cas et prépositions de temps]
Choisissez la bonne réponse parmi les trois propositions. Une seule réponse est correcte.
\begin{enumerate}
\item Der Sprecher informiert \_\_\_\_\_\_ Zuhörer über die Verspätung.
  \begin{enumerate}
  \item den \item dem \item das
  \end{enumerate}
\item Die Veranstaltung beginnt \_\_\_\_\_\_ 14.00 Uhr.
  \begin{enumerate}
  \item am \item um \item von
  \end{enumerate}
\item Der Lehrer hilft \_\_\_\_\_\_ Schülern bei der Anmeldung.
  \begin{enumerate}
  \item den \item dem \item der
  \end{enumerate}
\item \_\_\_\_\_\_ Montag gibt es keine Schule.
  \begin{enumerate}
  \item Am \item Um \item Seit
  \end{enumerate}
\item Wir hören \_\_\_\_\_\_ Radio seit einer Stunde.
  \begin{enumerate}
  \item das \item dem \item den
  \end{enumerate}
\end{enumerate}
\end{exercice}

\begin{exercice}[Vrai ou Faux justifié : Communiqué radio]
Lisez le communiqué suivant et décidez si les affirmations sont vraies (Richtig) ou fausses (Falsch). Justifiez chaque réponse par une citation exacte du texte.
\begin{textebox}[Communiqué de la station « Radio Yaoundé »]
„Achtung, liebe Hörer! Hier spricht Radio Yaoundé. Wir melden eine wichtige Änderung: Das Fußballspiel zwischen Canon Yaoundé und Tonnerre Kalara Club findet nicht heute um 15.30 Uhr statt. Der Platz ist wegen starker Regenfälle unbespielbar. Die Partie wird auf nächsten Samstag, den 20. Oktober, um 16.00 Uhr verschoben. Bitte beachten Sie diese Mitteilung. Vielen Dank für Ihre Aufmerksamkeit.“
\end{textebox}
\begin{enumerate}
\item Das Spiel findet heute wie geplant statt.
\item Der Grund für die Verschiebung ist das Wetter.
\item Das neue Datum ist der 20. Oktober, 15.30 Uhr.
\item Der Sprecher bittet um Aufmerksamkeit.
\end{enumerate}
\end{exercice}

\begin{exercice}[Vocabulaire : Familles de mots et articles]
Complétez le tableau en donnant le nom correspondant (avec article défini et pluriel), le verbe de base et la traduction française. Le premier est fait en exemple.
\begin{center}
\begin{tabularx}{\linewidth}{|X|X|X|X|}
\hline
\rowcolor{popblueL} \textbf{Verbe (Infinitiv)} & \textbf{Nom (Artikel + Plural)} & \textbf{Adjektiv / Partizip} & \textbf{Français (Verbe / Nom)} \\ \hline
sprechen & die Sprache, die Sprachen & sprachlich & parler / la langue \\ \hline
melden &  &  &  \\ \hline
veranstalten &  &  &  \\ \hline
informieren &  &  &  \\ \hline
durchsagen &  &  &  \\ \hline
\end{tabularx}
\end{center}
\end{exercice}

\begin{exercice}[Conjugaison : Verbes du thème]
Conjuguez les verbes entre parenthèses au temps indiqué (Präsens ou Perfekt). Attention aux verbes forts et aux verbes à particule séparable.
\begin{enumerate}
\item Der Sprecher \_\_\_\_\_\_ (ankündigen, Präsens) die Pause.
\item Wir \_\_\_\_\_\_ (hören, Perfekt) die Meldung gestern im Radio.
\item Die Schule \_\_\_\_\_\_ (beginnen, Präsens) um 7.30 Uhr.
\item Ich \_\_\_\_\_\_ (mitteilen, Perfekt) dir die Nachricht vor einer Stunde.
\item Der Lehrer \_\_\_\_\_\_ (helfen, Präsens) den Schülern.
\end{enumerate}
\end{exercice}

\courssub{Je m'entraîne}

\begin{exercice}[Traduction : Français $\rightarrow$ Allemand (Cas et prépositions de temps)]
Traduisez les phrases suivantes en allemand. Respectez l'ordre des mots (place du verbe, compléments).
\begin{enumerate}
\item Le speaker informe les auditeurs (Akkusativ).
\item La fête commence samedi à 10 heures (um + Uhrzeit).
\item La rue est fermée de 8 heures à 20 heures (von ... bis).
\item Depuis lundi, j'écoute cette station (seit + Dativ).
\item Le professeur aide l'élève (Dativ).
\item Nous attendons le bus depuis une heure (akkusativischer Zeitausdruck).
\end{enumerate}
\end{exercice}

\begin{exercice}[Transformation : Remplacement par pronoms personnels]
Remplacez les groupes nominaux soulignés par le pronom personnel correspondant (Nominativ, Akkusativ ou Dativ). Réécrivez la phrase complète.
\begin{enumerate}
\item Der Sprecher begrüßt \textbf{die Zuhörer}.
\item Die Lehrerin hilft \textbf{dem Schüler}.
\item Wir besuchen \textbf{die Veranstaltung} am Samstag.
\item Ich gebe \textbf{meinem Freund} die Information.
\item \textbf{Das Radio} spielt Musik.
\end{enumerate}
\end{exercice}

\begin{exercice}[Texte à trous : Communiqué scolaire]
Complétez ce communiqué radio avec les prépositions de temps ou les articles au bon cas (accusatif/datif) appropriés.
\begin{textebox}[Wichtige Durchsage]
„Liebe Schülerinnen und Schüler, \_\_\_\_\_\_ (1. préposition temps) Montag, dem 12. November, fällt der Unterricht \_\_\_\_\_\_ (2. préposition durée) 08.00 Uhr \_\_\_\_\_\_ (3. préposition temps) 14.00 Uhr aus. Grund: \_\_\_\_\_\_ (4. article Akk.) Aula wird \_\_\_\_\_\_ (5. préposition temps) einer Versammlung benötigt. Die Prüfung in Deutsch wird \_\_\_\_\_\_ (6. préposition temps) Mittwoch, \_\_\_\_\_\_ (7. préposition temps) 10.00 Uhr, nachgeholt. Bitte informieren Sie \_\_\_\_\_\_ (8. article Akk.) Eltern. Vielen Dank.“
\end{textebox}
\end{exercice}

\begin{exercice}[Appariement : Vocabulaire des médias]
Reliez chaque terme allemand à sa définition ou son équivalent français.
\begin{center}
\begin{tabular}{|l|l|}
\hline
\textbf{Allemand} & \textbf{Français / Définition} \\ \hline
1. die Durchsage & A. Personne qui parle à la radio \\ \hline
2. der Sprecher & B. Annonce officielle, brève, par haut-parleur \\ \hline
3. die Meldung & C. Émission, programme \\ \hline
4. die Sendung & D. Information, nouvelle (souvent courte) \\ \hline
5. der Hörer & E. Auditeur, celui qui écoute \\ \hline
\end{tabular}
\end{center}
Écrivez les paires correctes (ex. 1 -- B).
\end{exercice}

\courssub{Je me prépare au Bac}

\begin{exercice}[Compréhension écrite et langue : Communiqué original (Type Bac)]
Lisez attentivement le texte suivant, puis répondez aux questions en allemand (phrases complètes) et aux items Vrai/Faux.
\begin{textebox}[Communiqué de « Radio Campus Douala » – 15. Mai 2025]
„Hier ist Radio Campus Douala, Ihre Stimme für Bildung und Kultur. Wir senden live vom Campus der Universität. Liebe Studierende, wir müssen Sie leider informieren, dass die mündlichen Prüfungen im Fach Germanistik, die für morgen, den 16. Mai, ab 09.00 Uhr in Hörsaal 3 geplant waren, verschoben werden müssen. Der Prüfer, Herr Professor Meyer, ist kurzfristig erkrankt. Ein Ersatztermin steht noch nicht fest. Die neue Terminbekanntgabe erfolgt \_\_\_\_\_\_ (Präposition) spätestens Freitag \_\_\_\_\_\_ (Präposition) 14.00 Uhr über die offiziellen Aushänge und hier im Radio. Bitte verfolgen Sie \_\_\_\_\_\_ (Artikel Akk.) weiteren Mitteilungen. Wir bitten \_\_\_\_\_\_ (Präposition + Dativ) um Ihr Verständnis. Ihr Prüfungsamt.“
\end{textebox}
\textbf{A. Complétez le texte} avec les prépositions de temps (spätestens, über, um, am, bis) et les articles au cas correct (accusatif/datif) aux endroits numérotés (1 à 4 dans le texte).
\textbf{B. Questions de compréhension (en allemand, phrases complètes) :}
\begin{enumerate}
\item Wer sendet diesen Communiqué?
\item Welche Prüfungen werden verschoben?
\item Warum findet die Prüfung morgen nicht statt?
\item Wann und wie wird der neue Termin bekannt gegeben?
\item Was bittet das Prüfungsamt die Studierenden zu tun?
\end{enumerate}
\textbf{C. Vrai ou Faux ? Justifiez par une citation du texte.}
\begin{enumerate}
\item Die schriftlichen Prüfungen sind betroffen.
\item Der neue Termin ist bereits bekannt.
\item Die Information erfolgt nur per Radio.
\end{enumerate}
\end{exercice}

\begin{exercice}[Grammaire et Traduction : Version + Thème lexical]
\textbf{A. Version (Français $\rightarrow$ Allemand).} Traduisez en allemand correct (cas, prépositions, ordre des mots).
\begin{enumerate}
\item Le speaker informe les auditeurs.
\item La fête commence samedi à 10 heures.
\item La rue est fermée de 8 heures à 20 heures.
\item Depuis lundi, j'écoute cette station.
\end{enumerate}
\textbf{B. Thème lexical : Familles de mots.}
Pour chaque verbe, donnez le nom correspondant (article + pluriel), un adjectif ou participe dérivé, et la traduction française du nom.
\begin{center}
\begin{tabularx}{\linewidth}{|X|X|X|X|}
\hline
\rowcolor{popblueL} \textbf{Verbe} & \textbf{Nom (Art. + Pl.)} & \textbf{Dérivé (Adj./Part.)} & \textbf{Français (Nom)} \\ \hline
sprechen &  &  &  \\ \hline
melden &  &  &  \\ \hline
veranstalten &  &  &  \\ \hline
informieren &  &  &  \\ \hline
\end{tabularx}
\end{center}
\end{exercice}

\begin{exercice}[Expression écrite guidée : Rédaction d'un communiqué (Type Bac)]
\textbf{Consigne :} Vous êtes le speaker de « Radio Lycée Bilingue Yaoundé ». Rédigez un communiqué radio (environ \textbf{60--80 mots}) annonçant le report de l'examen d'allemand (LV2) prévu demain.
\textbf{Contraintes obligatoires :}
\begin{itemize}
\item Salutation formelle aux auditeurs (élèves, parents).
\item Annonce claire : report de l'examen d'allemand (Klausur / mündliche Prüfung).
\item Raison crédible (ex. : absence du professeur, problème technique, intempéries).
\item Nouvelle date OU indication de la publication prochaine de la nouvelle date (préposition de temps).
\item Consigne pour les élèves (réviser, consulter affichage, site web).
\item Formules de politesse de clôture.
\end{itemize}
Utilisez le vocabulaire de la leçon (\cle{die Durchsage}, \cle{verschieben}, \cle{der Prüfungstermin}, \cle{bekannt geben}, \cle{bitten um}) et soignez les cas (Akkusativ/Dativ) et les prépositions de temps (\cle{am}, \cle{um}, \cle{von ... bis}, \cle{seit}).
\end{exercice}

\newpage

\coursec{Corrigés des exercices}

\begin{corrige}[Exercice 1 — QCM : Cas et prépositions de temps]
\textbf{Réponses :}
\begin{enumerate}
\item \textbf{a) den} — \all{Der Sprecher informiert den Zuhörer über die Verspätung.} Le verbe \all{informieren} exige l'accusatif ; \all{der Zuhörer} est masculin singulier : \all{den}. « Le speaker informe l'auditeur du retard. »
\item \textbf{b) um} — \all{Die Veranstaltung beginnt um 14.00 Uhr.} Devant une heure précise, on utilise toujours \all{um}. « La manifestation commence à 14 heures. »
\item \textbf{a) den} — \all{Der Lehrer hilft den Schülern bei der Anmeldung.} Le verbe \all{helfen} exige le datif ; ici pluriel : \all{den Schülern} (datif pluriel avec \all{-n}). « Le professeur aide les élèves. »
\item \textbf{a) Am} — \all{Am Montag gibt es keine Schule.} Devant un jour de la semaine, on utilise \all{an} + datif : \all{am}. « Lundi, il n'y a pas école. »
\item \textbf{a) das} — \all{Wir hören das Radio seit einer Stunde.} \all{hören} exige l'accusatif ; \all{das Radio} est neutre, l'accusatif reste \all{das}. « Nous écoutons la radio depuis une heure. »
\end{enumerate}
\textbf{Point de vigilance :} ne pas confondre \all{helfen} (+ Dativ) avec le français « aider quelqu'un » (complément direct). En allemand : \all{Ich helfe dir}, jamais \all{ich helfe dich}.
\end{corrige}

\begin{corrige}[Exercice 2 — Vrai ou Faux justifié : Communiqué radio]
\begin{enumerate}
\item \textbf{Falsch.} Justification : „Das Fußballspiel \ldots findet \textbf{nicht heute} um 15.30 Uhr statt.“ Le match est reporté, il n'a pas lieu aujourd'hui.
\item \textbf{Richtig.} Justification : „Der Platz ist wegen \textbf{starker Regenfälle} unbespielbar.“ La cause du report est bien météorologique (fortes pluies).
\item \textbf{Falsch.} Justification : „Die Partie wird auf nächsten Samstag, den 20. Oktober, um \textbf{16.00 Uhr} verschoben.“ La date est correcte, mais la nouvelle heure est 16 h, pas 15 h 30.
\item \textbf{Richtig.} Justification : „Vielen Dank für Ihre \textbf{Aufmerksamkeit}“ et „Bitte beachten Sie diese Mitteilung.“ Le speaker demande bien l'attention des auditeurs.
\end{enumerate}
\textbf{Méthode :} pour un Vrai/Faux, citez toujours les mots exacts du texte entre guillemets allemands ; une réponse sans citation n'est pas complète.
\end{corrige}

\begin{corrige}[Exercice 3 — Vocabulaire : Familles de mots et articles]
\begin{center}
\begin{tabularx}{\linewidth}{|X|X|X|X|}
\hline
\rowcolor{popblueL} \textbf{Verbe} & \textbf{Nom (Artikel + Plural)} & \textbf{Adjektiv / Partizip} & \textbf{Français} \\ \hline
sprechen & die Sprache, die Sprachen & sprachlich & parler / la langue \\ \hline
melden & die Meldung, die Meldungen & gemeldet & annoncer, signaler / l'information, la nouvelle \\ \hline
veranstalten & die Veranstaltung, die Veranstaltungen & veranstaltet & organiser / la manifestation, l'événement \\ \hline
informieren & die Information, die Informationen & informativ, informiert & informer / l'information, le renseignement \\ \hline
durchsagen & die Durchsage, die Durchsagen & durchgesagt & annoncer (par haut-parleur) / l'annonce \\ \hline
\end{tabularx}
\end{center}
\textbf{Remarques :}
\begin{itemize}
\item Les noms en \all{-ung} et \all{-tion} sont toujours féminins (\all{die}) et forment leur pluriel en \all{-en} : \all{die Meldung, -en} ; \all{die Information, -en}.
\item \all{durchsagen} est un verbe à particule \emph{séparable} (l'accent tombe sur \all{durch-}) : au Präsens, la particule va en fin de phrase : \all{Der Sprecher sagt die Verspätung durch.} Au Perfekt, la particule se colle au participe : \all{Er hat die Verspätung durchgesagt.} « Le speaker a annoncé le retard. »
\item Même logique pour \all{mitteilen} : \all{Ich teile dir die Nachricht mit.} $\rightarrow$ \all{Ich habe dir die Nachricht mitgeteilt.}
\end{itemize}
\end{corrige}

\begin{corrige}[Exercice 4 — Conjugaison : Verbes du thème]
\begin{enumerate}
\item \all{Der Sprecher \textbf{kündigt} die Pause \textbf{an}.} — Verbe à particule séparable : au Präsens, la particule \all{an-} va en fin de phrase. « Le speaker annonce la pause. »
\item \all{Wir \textbf{haben} die Meldung gestern im Radio \textbf{gehört}.} — Perfekt avec \all{haben} + participe \all{gehört} (verbe faible). « Nous avons entendu l'information hier à la radio. »
\item \all{Die Schule \textbf{beginnt} um 7.30 Uhr.} — \all{beginnen} est un verbe fort, mais au Präsens la 3\up{e} personne est régulière : \all{beginnt}. « Les cours commencent à 7 h 30. »
\item \all{Ich \textbf{habe} dir die Nachricht vor einer Stunde \textbf{mitgeteilt}.} — \all{mitteilen} est séparable : participe \all{mitgeteilt} (\all{mit} + \all{geteilt}). « Je t'ai communiqué la nouvelle il y a une heure. »
\item \all{Der Lehrer \textbf{hilft} den Schülern.} — Verbe fort avec changement de voyelle \all{e} $\rightarrow$ \all{i} aux 2\up{e} et 3\up{e} personnes du singulier ; \all{helfen} + Dativ. « Le professeur aide les élèves. »
\end{enumerate}
\textbf{Points de vigilance :} (1) la particule séparable se détache uniquement au Präsens et au Präteritum en phrase affirmative ; (2) au Perfekt, elle se colle au participe : \all{angekündigt}, \all{mitgeteilt}, \all{durchgesagt}.
\end{corrige}

\begin{corrige}[Exercice 5 — Traduction : Français $\rightarrow$ Allemand]
\begin{enumerate}
\item \all{Der Sprecher informiert die Zuhörer.} — \all{informieren} + Akkusativ ; \all{die Zuhörer} : pluriel identique au singulier.
\item \all{Die Feier beginnt am Samstag um 10.00 Uhr.} — Jour : \all{am} ; heure : \all{um}. On peut aussi dire : \all{Am Samstag beginnt die Feier um 10.00 Uhr} (verbe en 2\up{e} position).
\item \all{Die Straße ist von 8.00 Uhr bis 20.00 Uhr geschlossen.} — Cadre \all{von \ldots bis} ; attention au ß de \all{Straße}.
\item \all{Seit Montag höre ich diesen Sender.} — \all{seit} + Dativ (\all{Montag} sans article) ; si \all{seit} ouvre la phrase, le verbe reste en 2\up{e} position : inversion du sujet. « Station » de radio = \all{der Sender}.
\item \all{Der Lehrer hilft dem Schüler.} — \all{helfen} + Dativ : \all{dem Schüler} (masculin singulier).
\item \all{Wir warten seit einer Stunde auf den Bus.} — \all{warten auf} + Akkusativ ; \all{seit einer Stunde} (datif féminin). Autre solution acceptée avec complément de durée à l'accusatif : \all{Wir warten schon eine Stunde auf den Bus.}
\end{enumerate}
\textbf{Barème indicatif (12 pts) :} 2 pts par phrase (1 pt vocabulaire et cas, 1 pt ordre des mots et orthographe).
\end{corrige}

\begin{corrige}[Exercice 6 — Transformation : Pronoms personnels]
\begin{enumerate}
\item \all{Der Sprecher begrüßt \textbf{sie}.} — \all{die Zuhörer} : pluriel, accusatif $\rightarrow$ \all{sie}. « Le speaker salue les auditeurs / il les salue. »
\item \all{Die Lehrerin hilft \textbf{ihm}.} — \all{dem Schüler} : masculin, datif $\rightarrow$ \all{ihm}. « La professeure l'aide (lui). »
\item \all{Wir besuchen \textbf{sie} am Samstag.} — \all{die Veranstaltung} : féminin, accusatif $\rightarrow$ \all{sie}. « Nous la visitons samedi. »
\item \all{Ich gebe \textbf{ihm} die Information.} — \all{meinem Freund} : masculin, datif $\rightarrow$ \all{ihm}. « Je lui donne l'information. »
\item \all{\textbf{Es} spielt Musik.} — \all{das Radio} : neutre, nominatif $\rightarrow$ \all{es}. « Elle (la radio) joue de la musique. »
\end{enumerate}
\textbf{Rappel :} le pronom suit le \emph{genre grammatical} allemand, pas le genre français : \all{das Radio} $\rightarrow$ \all{es}, \all{die Veranstaltung} $\rightarrow$ \all{sie}.
\end{corrige}

\begin{corrige}[Exercice 7 — Texte à trous : Communiqué scolaire]
Texte complété :

„Liebe Schülerinnen und Schüler, \textbf{am} Montag, dem 12. November, fällt der Unterricht \textbf{von} 08.00 Uhr \textbf{bis} 14.00 Uhr aus. Grund: \textbf{die} Aula wird \textbf{für} eine Versammlung benötigt. Die Prüfung in Deutsch wird \textbf{am} Mittwoch, \textbf{um} 10.00 Uhr, nachgeholt. Bitte informieren Sie \textbf{Ihre} Eltern. Vielen Dank.“

« Chers élèves, lundi 12 novembre, les cours de 8 h à 14 h sont annulés. Raison : la salle des fêtes est requise pour une assemblée. L'épreuve d'allemand sera rattrapée mercredi à 10 heures. Veuillez informer vos parents. Merci. »

\textbf{Justifications :}
\begin{enumerate}
\item \all{am} : jour de la semaine $\rightarrow$ \all{an} + Dativ.
\item \all{von} \ldots{} 3. \all{bis} : cadre de durée « de \ldots à \ldots ».
\item \all{die} : \all{die Aula} est féminin ; ici nominatif (sujet du passif \all{wird benötigt}). L'accusatif féminin serait identique : \all{die}.
\item \all{für} : « pour une réunion » ; \all{für} + Akkusativ (\all{eine Versammlung}).
\item \all{am} : jour $\rightarrow$ \all{an} + Dativ.
\item \all{um} : heure précise.
\item \all{Ihre} : \all{informieren} + Akkusativ ; \all{die Eltern} pluriel $\rightarrow$ \all{Ihre} (accusatif pluriel).
\end{enumerate}
\textbf{Point de vigilance :} \all{ausfallen} (être annulé) est séparable : \all{der Unterricht fällt aus} ; au Perfekt : \all{der Unterricht ist ausgefallen} (avec \all{sein}, car il s'agit d'un changement d'état).
\end{corrige}

\begin{corrige}[Exercice 8 — Appariement : Vocabulaire des médias]
\textbf{Paires correctes :} 1 -- B ; 2 -- A ; 3 -- D ; 4 -- C ; 5 -- E.
\begin{itemize}
\item \all{die Durchsage, -n} : annonce officielle brève, souvent par haut-parleur (gare, école, stade).
\item \all{der Sprecher, -} : la personne qui parle (à la radio, au micro).
\item \all{die Meldung, -en} : information courte, nouvelle.
\item \all{die Sendung, -en} : émission, programme de radio ou de télévision.
\item \all{der Hörer, -} : l'auditeur, celui qui écoute.
\end{itemize}
\textbf{Astuce :} le suffixe \all{-ung} donne des noms féminins d'action ; \all{-er} désigne la personne qui fait l'action (\all{sprechen} $\rightarrow$ \all{der Sprecher}, \all{hören} $\rightarrow$ \all{der Hörer}).
\end{corrige}

\begin{corrige}[Exercice 9 — Compréhension écrite et langue : Type Bac]
\textbf{A. Texte complété :}
\begin{enumerate}
\item \all{spätestens \textbf{am} Freitag} — jour de la semaine : \all{an} + Dativ. Formulation complète : „Die neue Terminbekanntgabe erfolgt spätestens am Freitag um 14.00 Uhr.“
\item \all{\textbf{um}} 14.00 Uhr — heure précise : toujours \all{um}.
\item „Bitte verfolgen Sie \textbf{die} weiteren Mitteilungen.“ — \all{verfolgen} + Akkusativ ; \all{die Mitteilungen} pluriel $\rightarrow$ \all{die}.
\item „Wir bitten \textbf{Sie} \textbf{um} Ihr Verständnis.“ — \all{bitten} + Akkusativ de la personne + \all{um} + Akkusativ de la chose. « Nous vous demandons votre compréhension. »
\end{enumerate}

\textbf{B. Questions de compréhension — réponses modèles :}
\begin{enumerate}
\item \all{Radio Campus Douala sendet diesen Communiqué live vom Campus der Universität.}
\item \all{Die mündlichen Prüfungen im Fach Germanistik werden verschoben.}
\item \all{Die Prüfung findet nicht statt, weil der Prüfer, Herr Professor Meyer, kurzfristig erkrankt ist.}
\item \all{Der neue Termin wird spätestens am Freitag um 14.00 Uhr über die offiziellen Aushänge und im Radio bekannt gegeben.}
\item \all{Das Prüfungsamt bittet die Studierenden, die weiteren Mitteilungen zu verfolgen, und bittet um ihr Verständnis.}
\end{enumerate}

\textbf{C. Vrai ou Faux :}
\begin{enumerate}
\item \textbf{Falsch.} „die \textbf{mündlichen} Prüfungen im Fach Germanistik \ldots werden verschoben“ — ce sont les épreuves orales, pas les écrites.
\item \textbf{Falsch.} „Ein Ersatztermin steht \textbf{noch nicht fest}.“
\item \textbf{Falsch.} „über die offiziellen \textbf{Aushänge und} hier im Radio“ — l'information passe par deux canaux, pas seulement la radio.
\end{enumerate}

\textbf{Barème indicatif (20 pts) :} A : 4 pts (1 pt par trou) ; B : 10 pts (2 pts par réponse : 1 pt contenu, 1 pt langue) ; C : 6 pts (1 pt réponse + 1 pt citation par item).

\textbf{Points de vigilance :} répondre par des phrases complètes avec le verbe en 2\up{e} position ; ne pas recopier tout le texte ; dans \all{bitten um}, la personne est à l'accusatif (\all{wir bitten Sie}), la chose est introduite par \all{um} ; dans la subordonnée avec \all{weil}, le verbe conjugué va à la fin : \all{weil der Prüfer erkrankt ist}.
\end{corrige}

\begin{corrige}[Exercice 10 — Grammaire et Traduction : Version + Thème lexical]
\textbf{A. Version :}
\begin{enumerate}
\item \all{Der Sprecher informiert die Zuhörer.}
\item \all{Die Feier beginnt am Samstag um 10.00 Uhr.}
\item \all{Die Straße ist von 8.00 Uhr bis 20.00 Uhr geschlossen.}
\item \all{Seit Montag höre ich diesen Sender.}
\end{enumerate}
Voir les justifications détaillées au corrigé de l'exercice 5 : \all{informieren} + Akkusativ ; \all{am} + jour, \all{um} + heure ; cadre \all{von \ldots bis} ; \all{seit} + Dativ avec inversion du sujet.

\textbf{B. Thème lexical :}
\begin{center}
\begin{tabularx}{\linewidth}{|X|X|X|X|}
\hline
\rowcolor{popblueL} \textbf{Verbe} & \textbf{Nom (Art. + Pl.)} & \textbf{Dérivé} & \textbf{Français (Nom)} \\ \hline
sprechen & die Sprache, -n & sprachlich & la langue \\ \hline
melden & die Meldung, -en & gemeldet & l'information, la nouvelle \\ \hline
veranstalten & die Veranstaltung, -en & veranstaltet & la manifestation, l'événement \\ \hline
informieren & die Information, -en & informativ & l'information, le renseignement \\ \hline
\end{tabularx}
\end{center}
\textbf{Barème indicatif (16 pts) :} A : 8 pts (2 pts par phrase) ; B : 8 pts (2 pts par ligne : article et pluriel corrects exigés).

\textbf{Points de vigilance :} l'article et le pluriel font partie du mot : \all{die Meldung, -en} ; sans article, la réponse est incomplète. Majuscule obligatoire à tous les noms.
\end{corrige}

\begin{corrige}[Exercice 11 — Expression écrite guidée : Rédaction d'un communiqué]
\textbf{Proposition de rédaction modèle (environ 75 mots) :}

„Liebe Schülerinnen und Schüler, liebe Eltern, hier spricht Radio Lycée Bilingue Yaoundé. Wir müssen Ihnen eine wichtige Durchsage mitteilen: Die Deutschklausur, die für morgen geplant war, wird verschoben, weil unsere Deutschlehrerin krank ist. Der neue Prüfungstermin wird am Freitag um 10.00 Uhr am schwarzen Brett und auf der Website bekannt gegeben. Bitte wiederholen Sie zu Hause den Wortschatz und die Grammatik. Wir bitten Sie um Ihr Verständnis. Vielen Dank für Ihre Aufmerksamkeit.“

\textbf{Traduction :} « Chers élèves, chers parents, ici Radio Lycée Bilingue Yaoundé. Nous devons vous communiquer une annonce importante : l'épreuve écrite d'allemand, prévue pour demain, est reportée parce que notre professeure d'allemand est malade. La nouvelle date sera annoncée vendredi à 10 heures au tableau d'affichage et sur le site web. Veuillez réviser à la maison le vocabulaire et la grammaire. Nous vous demandons votre compréhension. Merci de votre attention. »

\textbf{Barème indicatif (20 pts) :}
\begin{itemize}
\item Respect des contraintes de contenu (salutation, annonce, raison, date, consigne, clôture) : 6 pts.
\item Vocabulaire de la leçon (\all{die Durchsage}, \all{verschieben}, \all{der Prüfungstermin}, \all{bekannt geben}, \all{bitten um}) : 4 pts.
\item Grammaire : cas corrects (Akk./Dat.), prépositions de temps (\all{am}, \all{um}), place du verbe, particules séparables : 6 pts.
\item Orthographe (Umlaute, majuscules des noms), longueur respectée (60--80 mots) : 4 pts.
\end{itemize}
\textbf{Points de vigilance :}
\begin{itemize}
\item \all{verschieben} est un verbe fort séparable : \all{die Klausur wird verschoben} (passif) ou \all{wir verschieben die Klausur}.
\item \all{bekannt geben} : participe \all{bekannt gegeben} — \all{wird bekannt gegeben}.
\item \all{bitten um} : la personne à l'accusatif (\all{wir bitten Sie}), la chose après \all{um}.
\item Subordonnée avec \all{weil} : verbe conjugué à la fin : \all{weil unsere Deutschlehrerin krank ist}.
\item Éviter le faux ami : « annoncer » = \all{ankündigen} ou \all{bekannt geben}, jamais \all{annoncieren}.
\end{itemize}
\end{corrige}

\newpage

\coursec{Fiche bilan}

\begin{popfigure}
\centering
\begin{tikzpicture}[
  mindmap,
  concept color=popblue,
  level 1 concept/.append style={level distance=3.5cm, sibling angle=60},
  level 2 concept/.append style={level distance=2.5cm},
  every node/.style={font=\small, align=center},
  branch1/.style={concept color=poporange},
  branch2/.style={concept color=popgreen},
  branch3/.style={concept color=poppurple},
  branch4/.style={concept color=poppink},
  branch5/.style={concept color=popgold},
  branch6/.style={concept color=popblue!70!popdark}
]
\node[concept, text=white, font=\bfseries\large, align=center]{Le communiqué radio\\ (ALA17)}
  child[branch1] { node[concept, text=white, align=center]{Texte support\\ \emph{Das Radioprogramm}} }
  child[branch2] { node[concept, text=white, align=center]{Vocabulaire\\ Médias \& Radio} }
  child[branch3] { node[concept, text=white, align=center]{Grammaire 1\\ Accusatif / Datif} }
  child[branch4] { node[concept, text=white, align=center]{Grammaire 2\\ Prépositions de temps} }
  child[branch5] { node[concept, text=white, align=center]{Compréhension\\ \& Expression} }
  child[branch6] { node[concept, text=white, align=center]{Méthode\\ Commentaire guidé} };
\end{tikzpicture}
\legende{Carte mentale de la leçon ALA17 — Média et communication}
\end{popfigure}

\begin{bilan}

\textbf{Lexique clé — Medien und Radio}

\begin{center}
\begin{tabularx}{\linewidth}{|>{\bfseries}l|X|X|}\hline
\rowcolor{popblueL}
\all{Deutsch (Artikel, Plural)} & \textbf{Français} & \textbf{Exemple / Famille} \\ \hline
\all{das Radio, die Radios} & la radio (appareil / média) & \all{Ich höre Radio.} (J'écoute la radio.) \\ \hline
\all{die Mitteilung, die Mitteilungen} & la communication, l'avis & \all{eine wichtige Mitteilung machen} \\ \hline
\all{die Durchsage, die Durchsagen} & l'annonce (haut-parleur, radio) & \all{durchsagen} (annoncer), \all{durchgeben} \\ \hline
\all{der Sprecher, die Sprecher} & le locuteur, l'annonceur & \all{sprechen} $\to$ \all{der Vortrag} (exposé) \\ \hline
\all{die Meldung, die Meldungen} & la nouvelle, le bulletin & \all{melden} (signaler), \all{die Nachricht} \\ \hline
\all{die Veranstaltung, die Veranstaltungen} & l'événement, la manifestation & \all{veranstalten} (organiser) \\ \hline
\all{der Hörer, die Hörer} & l'auditeur & \all{hören} $\to$ \all{das Hörspiel} (pièce radiophonique) \\ \hline
\all{das Programm, die Programme} & le programme, la grille & \all{das Radioprogramm}, \all{programmieren} \\ \hline
\all{senden (sendete, gesendet)} & émettre, diffuser & \all{der Sender} (émetteur/chaîne), \all{die Sendung} (émission) \\ \hline
\all{ankündigen} & annoncer, prévoir & \all{die Ankündigung} (l'annonce) \\ \hline
\all{informieren über (+ Akk)} & informer sur & \all{Wir informieren die Hörer über das Wetter.} \\ \hline
\end{tabularx}
\end{center}

\vspace{0.5cm}
\textbf{Grammaire 1 — L'Accusatif et le Datif (Rappel synthétique)}

\begin{center}
\begin{tabularx}{\linewidth}{|>{\bfseries}l|X|X|}\hline
\rowcolor{poporangeL}
\textbf{Cas} & \textbf{Fonction principale / Question} & \textbf{Marqueurs (Articles définis)} \\ \hline
\all{Nominativ} & Sujet \emph{(Wer? Was?)} & \all{der / die / das / die} \\ \hline
\all{Akkusativ} & Objet direct \emph{(Wen? Was?)} & \all{\textbf{den} / die / das / die} \\ \hline
\all{Dativ} & Objet indirect \emph{(Wem?)} & \all{\textbf{dem} / \textbf{der} / \textbf{dem} / \textbf{den} (+n pluriel)} \\ \hline
\end{tabularx}
\end{center}

\begin{itemize}
\item \textbf{Verbes à double objet :} Sujet (Nom) + Objet \emph{personne} (Dat) + Objet \emph{chose} (Akk). \all{Der Sprecher gibt \textbf{den Hörern} (Dat) \textbf{eine Information} (Akk).}
\item \textbf{Verbes « datifs » fréquents :} \all{helfen, danken, gefallen, folgen, glauben, gratulieren, fehlen, passen, schmecken, wehtun}. Ex. \all{Das Programm gefällt \textbf{mir} (Dat).}
\item \textbf{Prépositions casuelles :}
  \begin{itemize}
  \item \cle{Akkusativ seul :} \all{durch, für, gegen, ohne, um} (mnemonic : \emph{dfguo}).
  \item \cle{Datif seul :} \all{aus, bei, mit, nach, seit, von, zu} (mnemonic : \emph{abmnsvz}).
  \item \cle{Wechselpräpositionen (Akk/Dat) :} \all{an, auf, hinter, in, neben, über, unter, vor, zwischen}. \textbf{Règle :} Mouvement $\to$ \emph{Wo hin?} = Akk ; Position $\to$ \emph{Wo?} = Dat.
  \end{itemize}
\end{itemize}

\vspace{0.5cm}
\textbf{Grammaire 2 — Les prépositions de temps (Zeitpräpositionen)}

\begin{center}
\begin{tabularx}{\linewidth}{|>{\bfseries}l|X|X|}\hline
\rowcolor{popgreenL}
\all{Präposition} & \textbf{Cas / Emploi} & \textbf{Exemples (avec traduction)} \\ \hline
\all{an} & \textbf{Datif} — Jours de la semaine, dates, fêtes & \all{am Montag} (lundi), \all{am 17. Januar}, \all{an Weihnachten} \\ \hline
\all{in} & \textbf{Datif} — Mois, saisons, années, parties de journée & \all{im Januar}, \all{im Sommer}, \all{im Jahr 2025}, \all{morgens / abends} (mais \all{am Morgen}) \\ \hline
\all{um} & \textbf{Datif} — Heure précise & \all{um 14.30 Uhr} (à 14h30), \all{um Mitternacht} \\ \hline
\all{von ... bis / an} & \textbf{Datif} — Début / Fin (durée bornée) & \all{von Montag bis Freitag}, \all{von 8 bis 12 Uhr} \\ \hline
\all{seit} / \all{ab} & \textbf{Datif} — Depuis (début passé encore vrai) / À partir de (futur) & \all{seit 2020}, \all{ab morgen} \\ \hline
\all{für} & \textbf{Akkusativ} — Durée prévue (pour) & \all{für zwei Wochen} (pour deux semaines) \\ \hline
\all{während} & \textbf{Genitif} (langue soutenue) / \textbf{Datif} (parlé) — Pendant & \all{während der Sendung} / \all{während der Pause} \\ \hline
\end{tabularx}
\end{center}

\cle{Attention :} \all{am} = \all{an dem} ; \all{im} = \all{in dem} ; \all{zum} = \all{zu dem} ; \all{zur} = \all{zu der}.

\vspace{0.5cm}
\textbf{Méthode express — Commenter un communiqué radio}

\begin{enumerate}
\item \textbf{Contexte :} Identifier la source (quelle chaîne ?), la date, le type d'émission (info, météo, culture, sport).
\item \textbf{Sujet principal :} De quoi parle l'annonce ? (\all{Thema: Veranstaltung, Warnung, Information...})
\item \textbf{Structure :} Repérer \emph{Einleitung} (accroche), \emph{Hauptteil} (détails : qui, quoi, où, quand), \emph{Schluss} (rappel, contact).
\item \textbf{Lexique fonctionnel :} Noter les verbes d'action (\all{ankündigen, einladen, warnen, bitten}) et les marqueurs de temps (\all{am, um, von...bis}).
\item \textbf{Résumé / Production :} Rédiger 3--5 phrases en allemand (Présent/Perfekt) : \all{Der Sender informiert über... Es findet statt am... Interessierte können...}
\end{enumerate}

\end{bilan}

\begin{aretenir}[Checklist avant le Bac — ALA17]
$\square$ Je connais le vocabulaire de la radio (\all{die Durchsage, der Sprecher, die Meldung, senden, ankündigen}) avec articles et pluriels.\\
$\square$ Je distingue l'Accusatif (objet direct \emph{Wen?}) du Datif (objet indirect \emph{Wem?}) dans une phrase.\\
$\square$ Je décline correctement l'article défini aux trois cas (der/den/dem, die/die/der, das/das/dem, die/die/den+n).\\
$\square$ J'identifie les verbes à double objet (Akk + Dat) et les verbes « datifs » (\all{helfen, gefallen, danken}).\\
$\square$ Je choisis le bon cas après les prépositions : \all{für} + Akk, \all{mit} + Dat, \all{in} (temps) + Dat.\\
$\square$ J'utilise les prépositions de temps : \all{am} (jour), \all{im} (mois/saison), \all{um} (heure), \all{von...bis} (durée), \all{seit} (depuis).\\
$\square$ Je sais contracter \all{an dem $\to$ am, in dem $\to$ im, zu dem $\to$ zum, zu der $\to$ zur}.\\
$\square$ Je comprends globalement un communiqué radio authentique (sujet, lieu, date, public visé).\\
$\square$ Je peux répondre aux questions \emph{W-Fragen} (Wer? Was? Wann? Wo? Warum?) sur le texte.\\
$\square$ Je sais résumer un communiqué en 3--5 phrases simples (Présent / Perfekt).\\
$\square$ Je produis un court communiqué type (annonce événement scolaire) en respectant la structure : Titre -- Infos clés -- Contact.
\end{aretenir}

\findecours

\end{document}
